%%%%%%%%%%%%%%%%%%%%%%%%%%%%%%%%%%%%%%%%%%%%%%%%%%%%%%%%%%%%%%%%%
% MEMORIA TFG - DOCUMENTO ÚNICO (listo para Overleaf)
% TFG: Knowledge Distillation y despliegue eficiente de modelos
% de lenguaje en entornos de supercomputación
% Autor: Marc Romera Peral
%
% Estado: redactado hasta el Diseño del sistema inclusive.
% Implementación, experimentación y conclusiones quedan como
% placeholders explícitos para completar cuando el trabajo técnico
% esté suficientemente cerrado.
%%%%%%%%%%%%%%%%%%%%%%%%%%%%%%%%%%%%%%%%%%%%%%%%%%%%%%%%%%%%%%%%%
\documentclass[11pt,a4paper]{article}

%%%%%%%%%%%%%%%%%%%%%%%%%%%%%%%%%%%%%%%%%%%%%%%%%%%%%%%%%%%%%%%%%
% PAQUETES
%%%%%%%%%%%%%%%%%%%%%%%%%%%%%%%%%%%%%%%%%%%%%%%%%%%%%%%%%%%%%%%%%
\usepackage[utf8]{inputenc}
\usepackage[T1]{fontenc}
\usepackage[spanish]{babel}
\addto\captionsspanish{
    \renewcommand{\refname}{Bibliografía}
    \renewcommand{\tablename}{Tabla}
    \renewcommand{\listtablename}{Índice de tablas}
    \renewcommand{\listfigurename}{Índice de figuras}
    \renewcommand{\contentsname}{Índice}
}
\usepackage[margin=2cm]{geometry}
\usepackage{microtype}
\usepackage{hyperref}
\usepackage{graphicx}
\usepackage{booktabs}
\usepackage{enumitem}
\usepackage{array}
\usepackage{xcolor}
\usepackage{parskip}
\setlength{\parskip}{7pt plus 2pt minus 1pt}
\usepackage{float}
\usepackage{tabularx}
\usepackage{ltablex}
\usepackage{colortbl}
\usepackage{pdflscape}
\usepackage{needspace}
\usepackage{tikz}
\usepackage{pgf-pie}
\usetikzlibrary{positioning,arrows.meta,fit,backgrounds,calc,shapes.geometric,decorations.pathreplacing}
\usepackage{eurosym}
\usepackage{amsmath}
\usepackage{amssymb}
\usepackage{titlesec}
\usepackage{tocloft}
\keepXaliases

% Espaciado de títulos
\titlespacing*{\section}{0pt}{10pt}{4pt}
\titlespacing*{\subsection}{0pt}{8pt}{3pt}
\titlespacing*{\subsubsection}{0pt}{6pt}{2pt}
\titleformat{\section}{\normalfont\large\bfseries}{\thesection.}{0.4em}{}
\titleformat{\subsection}{\normalfont\normalsize\bfseries}{\thesubsection.}{0.4em}{}
\titleformat{\subsubsection}{\normalfont\small\bfseries}{\thesubsubsection.}{0.4em}{}

% Anchos del índice
\setlength{\cftsecnumwidth}{2.8em}
\setlength{\cftsubsecnumwidth}{3.8em}
\setlength{\cftsubsubsecnumwidth}{4.8em}
\setlength{\cftfignumwidth}{2.8em}
\setlength{\cfttabnumwidth}{2.8em}
\setlength{\cftbeforesecskip}{1pt}

\hypersetup{
    colorlinks=true,
    linkcolor=black,
    citecolor=black,
    urlcolor=blue
}
\setcounter{secnumdepth}{2}   % numerar solo section y subsection
\setcounter{tocdepth}{2}      % mostrar en el índice solo hasta subsection

\widowpenalty=10000
\clubpenalty=10000

% Colores reutilizados en diagramas
\definecolor{bcteacher}{RGB}{220,90,90}
\definecolor{bcstudent}{RGB}{90,140,220}
\definecolor{bcrouter}{RGB}{240,170,70}
\definecolor{bcdata}{RGB}{90,180,140}
\definecolor{bcmeasure}{RGB}{160,120,200}
\definecolor{bcneutro}{RGB}{210,210,215}

%%%%%%%%%%%%%%%%%%%%%%%%%%%%%%%%%%%%%%%%%%%%%%%%%%%%%%%%%%%%%%%%%
\begin{document}
%%%%%%%%%%%%%%%%%%%%%%%%%%%%%%%%%%%%%%%%%%%%%%%%%%%%%%%%%%%%%%%%%

%%%%%%%%%%%%%%%%%%%%%%%%%%%%%%%%%%%%%%%%%%%%%%%%%%%%%%%%%%%%%%%%%
% PORTADA
%%%%%%%%%%%%%%%%%%%%%%%%%%%%%%%%%%%%%%%%%%%%%%%%%%%%%%%%%%%%%%%%%
\begin{titlepage}
\centering

\includegraphics[width=1\textwidth]{upc.png}\\[1cm]

{\Large \textbf{Trabajo de Fin de Grado}}\\[0.5cm]
{\large Grado en Ingeniería Informática}\\[0.3cm]
{\large Especialidad: Computación}\\[1cm]

{\LARGE \textbf{Knowledge Distillation y despliegue eficiente de modelos de lenguaje en entornos de supercomputación}}\\[3cm]

\begin{center}
\begin{tabular}{r l}
\textbf{Autor:} & Marc Romera Peral \\[0.3cm]
\textbf{Dirección:} & Josep Lluís Berral García \\
                    & Pol Garcia Recasens \\[0.3cm]
\textbf{Departamento:} & Arquitectura de Computadores \\[0.3cm]
\textbf{Universidad:} & Universitat Politècnica de Catalunya \\[0.3cm]
\textbf{Centro:} & Facultat d'Informàtica de Barcelona \\[0.3cm]
\textbf{Curso:} & 2025--2026 \\
\end{tabular}
\end{center}
\end{titlepage}

%%%%%%%%%%%%%%%%%%%%%%%%%%%%%%%%%%%%%%%%%%%%%%%%%%%%%%%%%%%%%%%%%
% ABSTRACTS (es / ca / en)
%%%%%%%%%%%%%%%%%%%%%%%%%%%%%%%%%%%%%%%%%%%%%%%%%%%%%%%%%%%%%%%%%
\newpage
\section*{Resumen}
\addcontentsline{toc}{section}{Resumen}

Los modelos de lenguaje de gran tamaño (Large Language Models, LLMs) han alcanzado en los últimos años un nivel de calidad que los hace útiles en muchísimas tareas, pero su despliegue en producción sigue siendo caro. El coste dominante no es entrenarlos, sino servirlos, pues cada token generado requiere una ejecución completa del modelo y los modelos más potentes no caben ni siquiera en una única GPU de gama alta sin técnicas de gestión de memoria sofisticadas.

Este Trabajo de Fin de Grado estudia cómo reducir el coste computacional de la inferencia de LLMs sin degradar significativamente la calidad de sus respuestas. La hipótesis de partida es sencilla. No todas las peticiones necesitan el modelo más grande; muchas pueden resolverse con un modelo más pequeño y más barato, y conviene reservar al modelo grande solo para las peticiones que realmente lo justifican.

A partir de esa idea se diseña un sistema de cascada entre \textit{teacher} y \textit{student}, construido con modelos obtenidos mediante destilación del conocimiento (knowledge distillation), que compara experimentalmente tres políticas de servicio (baseline con el modelo grande o \textit{teacher}, cascada forzada que intenta primero con el modelo pequeño o \textit{student} y escala si no cumple un criterio mínimo, y cascada por confianza con un predictor que decide a priori qué modelo atiende cada petición). El sistema se despliega sobre el supercomputador MareNostrum~5 del Barcelona Supercomputing Center utilizando vLLM como motor de serving, y se mide tanto en métricas de calidad (accuracy en benchmarks estándar) como en métricas de servicio (TTFT, latencia p95, throughput y memoria GPU).

El objetivo último es caracterizar de forma reproducible el compromiso entre calidad y eficiencia de las tres políticas, y discutir bajo qué condiciones la cascada es preferible al uso directo del teacher.

\textbf{Palabras clave.} Modelos de lenguaje, destilación del conocimiento, inferencia eficiente, cascada de modelos, supercomputación, vLLM, MareNostrum~5.

\newpage
\section*{Resum}
\addcontentsline{toc}{section}{Resum}

Els models de llenguatge de gran mida (Large Language Models, LLMs) han assolit en els darrers anys un nivell de qualitat que els fa útils en moltíssimes tasques, però desplegar-los en producció continua sent car. El cost dominant no és entrenar-los, sinó servir-los, ja que cada token generat requereix una execució completa del model, i els models més potents no caben ni tan sols en una GPU d'alta gamma sense tècniques de gestió de memòria sofisticades.

Aquest Treball de Fi de Grau estudia com reduir el cost computacional de la inferència de LLMs sense degradar significativament la qualitat de les seves respostes. La hipòtesi de partida és senzilla. No totes les peticions necessiten el model més gran; moltes es poden resoldre amb un model més petit i barat, i convé reservar el gran només per a les peticions que realment ho justifiquen.

A partir d'aquesta idea es dissenya un sistema de cascada entre \textit{teacher} i \textit{student}, construït amb models obtinguts mitjançant destil·lació del coneixement (knowledge distillation), que compara experimentalment tres polítiques de servei (baseline amb el model gran o \textit{teacher}, cascada forçada que prova primer amb el model petit o \textit{student} i escala si no compleix un criteri mínim, i cascada per confiança amb un predictor que decideix a priori quin model atén cada petició). El sistema es desplega sobre el supercomputador MareNostrum~5 del Barcelona Supercomputing Center utilitzant vLLM com a motor de serving, i es mesura tant en mètriques de qualitat (accuracy en benchmarks estàndard) com en mètriques de servei (TTFT, latència p95, throughput i memòria GPU).

L'objectiu final és caracteritzar de manera reproduïble el compromís entre qualitat i eficiència de les tres polítiques, i discutir en quines condicions la cascada és preferible a l'ús directe del teacher.

\textbf{Paraules clau.} Models de llenguatge, destil·lació del coneixement, inferència eficient, cascada de models, supercomputació, vLLM, MareNostrum~5.

\newpage
\section*{Abstract}
\addcontentsline{toc}{section}{Abstract}

In recent years Large Language Models (LLMs) have reached a quality level that makes them useful in a wide range of tasks, but deploying them in production is still expensive. The dominant cost is not training, but serving, because every generated token requires a full forward pass of the model, and the most capable models do not fit on a single high-end GPU without sophisticated memory management techniques.

This Bachelor's Thesis studies how to reduce the computational cost of LLM inference without significantly degrading the quality of its answers. The starting hypothesis is simple. Not every request needs the largest model; many of them can be handled by a smaller and cheaper model, and the large one should be reserved only for requests that actually require it.

Building on this idea, we design a cascade between \textit{teacher} and \textit{student}, composed of models obtained through knowledge distillation, and we experimentally compare three serving policies (baseline with the large model or \textit{teacher}, forced cascade that tries first with the small model or \textit{student} and escalates if the answer does not meet a minimum quality criterion, and confidence-based cascade with a predictor to decide a priori which model should handle each request). The system is deployed on MareNostrum~5 at the Barcelona Supercomputing Center using vLLM as the serving engine, and is measured both in quality metrics (accuracy on standard benchmarks) and in serving metrics (TTFT, p95 latency, throughput and GPU memory).

The ultimate goal is to characterise, in a reproducible way, the quality--efficiency trade-off of the three policies, and to discuss under which conditions the cascade is preferable to using the teacher directly.

\textbf{Keywords.} Large language models, knowledge distillation, efficient inference, model cascade, supercomputing, vLLM, MareNostrum~5.

%%%%%%%%%%%%%%%%%%%%%%%%%%%%%%%%%%%%%%%%%%%%%%%%%%%%%%%%%%%%%%%%%
% AGRADECIMIENTOS
%%%%%%%%%%%%%%%%%%%%%%%%%%%%%%%%%%%%%%%%%%%%%%%%%%%%%%%%%%%%%%%%%
\newpage
\section*{Agradecimientos}
\addcontentsline{toc}{section}{Agradecimientos}

A mis directores, Josep Lluís Berral y Pol Garcia, por hacer posible este proyecto, por su guía técnica y por el tiempo dedicado a las sesiones de seguimiento. Al Barcelona Supercomputing Center y, en particular, al grupo Data Centric Computing, por facilitar el acceso a MareNostrum~5 y el entorno experimental necesario para llevar a cabo los experimentos.

A la Facultat d'Informàtica de Barcelona y a los profesores con los que he coincidido durante el Grado, cuya formación ha hecho posible afrontar este trabajo. Y, finalmente, a mi familia y amigos por el apoyo durante todo el proceso.

%%%%%%%%%%%%%%%%%%%%%%%%%%%%%%%%%%%%%%%%%%%%%%%%%%%%%%%%%%%%%%%%%
% ÍNDICES
%%%%%%%%%%%%%%%%%%%%%%%%%%%%%%%%%%%%%%%%%%%%%%%%%%%%%%%%%%%%%%%%%
\newpage
\tableofcontents
\newpage
\listoftables
\newpage
\listoffigures
\newpage

\small

%%%%%%%%%%%%%%%%%%%%%%%%%%%%%%%%%%%%%%%%%%%%%%%%%%%%%%%%%%%%%%%%%
\part{Gestión de la tesis}
\label{part:gestion-tesis}
%%%%%%%%%%%%%%%%%%%%%%%%%%%%%%%%%%%%%%%%%%%%%%%%%%%%%%%%%%%%%%%%%

\section{Introducción y contextualización}

Los modelos de lenguaje de gran tamaño (Large Language Models, LLMs) han transformado el procesamiento del lenguaje natural en los últimos años. Modelos como LLaMA o Mistral muestran capacidades notables en generación de texto, comprensión lectora, traducción y razonamiento. Sin embargo, su entrenamiento y despliegue requieren recursos computacionales y costes de operación que dificultan su adopción fuera de organizaciones con infraestructuras significativas. Esta barrera de coste es uno de los principales obstáculos para la democratización de estas tecnologías.

El problema central que aborda este TFG puede formularse así. \textbf{¿Cómo reducir el coste computacional de inferencia de los modelos de lenguaje de gran tamaño (LLMs) sin degradar significativamente la calidad de sus respuestas?}

Este problema es relevante por varias razones. El coste elevado de inferir con LLMs grandes limita su accesibilidad a organizaciones con recursos considerables. Además, no todas las peticiones requieren la capacidad completa del modelo más potente. Muchas pueden resolverse satisfactoriamente con modelos más pequeños y eficientes. Esta observación sugiere que pueden diseñarse estrategias de asignación dinámica que seleccionen el modelo más adecuado en función de la dificultad estimada de cada petición.

Una de las aproximaciones posibles consiste en utilizar un sistema de cascada de modelos, respondiendo con el modelo pequeño cuando basta y escalando al grande solo cuando es necesario. En aplicaciones interactivas, donde el tiempo de respuesta importa, este tipo de enfoque puede mejorar la latencia media al evitar el uso del modelo grande en peticiones sencillas.

Para validar experimentalmente el equilibrio entre calidad y coste, se desplegará la propuesta en el Barcelona Supercomputing Center (BSC), concretamente en MareNostrum~5, utilizando vLLM como motor de serving. Es importante destacar que BSC y vLLM constituyen la infraestructura de experimentación y medición, no el objetivo principal del TFG. Proporcionan la plataforma para desplegar las distintas políticas de serving y registrar métricas de servicio (TTFT, latencia, throughput, memoria GPU), que se compararán con métricas de calidad en tareas NLP estándar.

\subsection{Marco institucional y encaje académico}

Este proyecto se enmarca en el Trabajo de Fin de Grado del Grado en Ingeniería Informática de la Facultat d'Informàtica de Barcelona (FIB), dentro de la especialidad de Computación. La modalidad del TFG es A. El tema aborda un problema central de la ingeniería de sistemas inteligentes, combinando arquitectura de sistemas, optimización de recursos y aprendizaje automático. El entorno de ejecución en MareNostrum~5 permite trabajar con infraestructura de supercomputación real, lo que aporta una dimensión práctica y profesionalizadora que complementa la formación teórica recibida en asignaturas como Arquitectura de Computadores, Aprendizaje Automático y Computación de Altas Prestaciones.

La colaboración con el BSC facilita el acceso a recursos de cómputo que no estarían disponibles en un entorno académico convencional. Esta colaboración permite validar las estrategias propuestas en condiciones realistas y genera resultados transferibles a otros entornos HPC o de producción con infraestructura similar.

\subsection{Actores implicados y beneficiarios}

Los usuarios del resultado serán investigadores y equipos técnicos que despliegan LLMs con restricciones de coste en entornos HPC. Los scripts y documentación podrán reutilizarse por estudiantes o investigadores del BSC en líneas similares. La operación recae sobre el autor con supervisión de los codirectores.

El coste computacional es asumido por el BSC mediante la cuota del grupo de investigación \textbf{Data Centric Computing} (con identificador bsc98 en el ms5). Los beneficiarios directos son las organizaciones que implementen cascadas similares, reduciendo uso de GPUs y consumo energético. La comunidad académica se beneficia de un pipeline documentado y reproducible.

%%%%%%%%%%%%%%%%%%%%%%%%%%%%%%%%%%%%%%%%%%%%%%%%%%%%%%%%%%%%%%%%%
% 2. JUSTIFICACIÓN, OBJETIVOS Y VISIÓN (agrupado para índice compacto)
%%%%%%%%%%%%%%%%%%%%%%%%%%%%%%%%%%%%%%%%%%%%%%%%%%%%%%%%%%%%%%%%%
\section{Justificación, objetivos y visión de la solución}

\subsection{Estado del arte y líneas de trabajo actuales}

El problema del coste de inferencia de los LLMs es hoy una línea de investigación muy activa, y la comunidad ha ido consolidando varias familias de respuestas que conviene conocer antes de presentar la propuesta de este TFG. Tres grandes bloques resumen razonablemente bien el panorama.

El primero agrupa las \textbf{técnicas de compresión de modelos}. La \textbf{cuantización} reduce la precisión numérica de los pesos (INT8, INT4) para acelerar cálculos y reducir memoria, aunque puede degradar la calidad de forma difícil de predecir. El \textbf{pruning} elimina conexiones redundantes para obtener modelos sparse más eficientes, pero requiere reentrenamiento y no siempre es compatible con hardware estándar. Técnicas de \textbf{PEFT} (\textit{parameter-efficient fine-tuning}), como LoRA, permiten adaptar modelos grandes con pocos parámetros entrenables, pero no reducen el coste de inferencia del modelo base. Una línea distinta, y la que se adopta en este TFG, es la de \textbf{Knowledge Distillation}~\cite{hinton2015distilling}, donde un modelo grande (teacher) transfiere conocimiento a uno más pequeño (student), que después se sirve a un coste muy inferior.

El segundo bloque es el de los \textbf{sistemas de serving eficiente}. Aquí el exponente más relevante es vLLM~\cite{vllm-docs}, un motor de inferencia con \emph{continuous batching} y \emph{PagedAttention}~\cite{vllm2023} que ha mejorado drásticamente el \textit{throughput} alcanzable con modelos medianos y grandes. Otras piezas del mismo bloque son FlashAttention~\cite{dao2022flashattention}, el \textit{speculative decoding}~\cite{leviathan2023speculative} y optimizaciones de bajo nivel sobre la atención. En este TFG, vLLM se usa como plataforma común de despliegue y medición, para que las comparaciones entre políticas se hagan en igualdad de condiciones.

El tercer bloque, más reciente y el que conecta directamente con la propuesta, es el de los \textbf{sistemas de decisión o routing entre modelos}, es decir, trabajar no con un único modelo sino con un conjunto de modelos de distinto tamaño y decidir, para cada petición, cuál es el más adecuado. FrugalGPT~\cite{chen2023frugalgpt} populariza la idea de cascada con verificación, y RouteLLM~\cite{dekoninck2024routing} explora el entrenamiento de un router que estima a priori qué modelo debería responder. El TFG se sitúa en este tercer bloque y, concretamente, propone usar KD para \emph{generar} el student que alimenta el sistema y después comparar, sobre la misma infraestructura de serving, tres políticas (\textbf{Política~A} en modo baseline, que sirve siempre con el teacher; \textbf{Política~B} en cascada forzada, que intenta primero con el modelo más pequeño y escala al grande si el resultado no pasa un filtro de calidad; y \textbf{Política~C} en cascada por confianza, que con un predictor externo decide \textit{a priori} qué modelo responde). La contribución del TFG es, por tanto, no tanto inventar una técnica nueva como \textbf{medir rigurosamente el trade-off entre calidad y coste} que esta combinación permite, en un entorno HPC real (MareNostrum~5) y con un protocolo de evaluación unificado.

\subsection{Relevancia práctica}

Reducir el coste de inferencia de LLMs es una prioridad en centros de investigación con recursos HPC compartidos y en organizaciones que despliegan asistentes conversacionales o sistemas de análisis de texto. Estas organizaciones buscan optimizar el uso de GPUs sin degradar la experiencia de usuario. Las cascadas de modelos responden a esta necesidad, usando el modelo pequeño cuando basta y escalando al grande solo cuando es necesario. Aunque el proyecto se ejecuta en MareNostrum~5 del BSC, los resultados son transferibles a otros entornos HPC o de producción con infraestructura similar.

\subsection{Relevancia académica}

El proyecto contribuye al ámbito académico mediante la documentación de un pipeline reproducible de evaluación, el análisis comparativo del balance entre calidad y eficiencia para las políticas mínimas A y B (y C si es viable), y la exploración de mecanismos de routing basados en confianza, un área con literatura emergente~\cite{chen2023frugalgpt, dekoninck2024routing}.

\subsection{Análisis de alternativas y decisión de diseño}

Se consideraron tres alternativas para reducir el coste de inferencia.

\textbf{Alternativa~1 (cuantización del teacher, INT8/INT4).} Reduce memoria y acelera inferencia sin modelos adicionales. Riesgo técnico bajo, pero ganancia limitada al factor de compresión y degradación difícil de controlar.

\textbf{Alternativa~2 (solo modelos pequeños preentrenados, contingencia).} Implementación sencilla, sin lógica de decisión. Riesgo muy bajo, pero calidad fija y sin flexibilidad para escalar en peticiones complejas.

\textbf{Alternativa~3 (elegida, sistema de cascada con múltiples políticas).} Combina modelos de distintos tamaños y compara políticas de serving (baseline, cascada forzada, cascada por confianza). Mayor complejidad, pero es la única que permite estudiar experimentalmente el equilibrio entre calidad y coste, objetivo central del TFG.

La Tabla~\ref{tab:alternativas} resume los puntos fuertes y débiles de cada alternativa.

\begin{table}[H]
\centering
\caption{Comparación de alternativas de diseño. Fuente: elaboración propia.}
\label{tab:alternativas}
\footnotesize
\begin{tabularx}{\textwidth}{>{\raggedright\arraybackslash}p{2.8cm} >{\raggedright\arraybackslash}X >{\raggedright\arraybackslash}X >{\raggedright\arraybackslash}p{3cm}}
\toprule
\textbf{Alternativa} & \textbf{Puntos fuertes} & \textbf{Puntos débiles} & \textbf{Encaje con el TFG} \\
\midrule
1. Cuantización del teacher & Bajo riesgo técnico; sin modelos adicionales; ganancia directa en memoria e inferencia & Ganancia limitada al factor de compresión; degradación de calidad difícil de controlar & Parcial: mejora eficiencia pero no permite analizar routing \\
\addlinespace
2. Solo modelos pequeños & Implementación sencilla; sin lógica de decisión & Calidad fija; sin flexibilidad ante peticiones complejas & Bajo: no aborda el equilibrio dinámico calidad--coste \\
\addlinespace
3. Cascada (elegida) & Permite comparar políticas de serving; estudia experimentalmente el equilibrio calidad--coste & Mayor complejidad de implementación y evaluación & Alto: responde directamente a la pregunta de investigación \\
\bottomrule
\end{tabularx}
\end{table}

Se adopta la tercera alternativa porque responde a la pregunta de investigación y mantiene como eje el marco entre \textit{teacher} y \textit{student} mediante knowledge distillation. Las otras resuelven parcialmente la eficiencia pero no permiten analizar el comportamiento dinámico de la cascada ni la comparación de políticas. Con acceso a MareNostrum~5, la complejidad adicional es asumible.

\subsection{Objetivo general}

El proyecto se estructura en torno a un objetivo principal y seis objetivos específicos (OE1--OE6), cada uno con criterios de aceptación cuantitativos y auditables.

El objetivo general de este TFG es \textbf{diseñar, implementar y evaluar un sistema de inferencia de LLMs basado en cascada de modelos entre \textit{teacher} y \textit{student}}, comparando experimentalmente las políticas mínimas de serving (baseline y cascada forzada) e incorporando la cascada por confianza si resulta viable, para analizar el equilibrio entre calidad y coste mediante despliegue en MareNostrum~5.

\subsection{Objetivos específicos y criterios de aceptación}
\label{subsec:oe-especificos}

A continuación, se detallan los objetivos específicos (OE1--OE6) con sus criterios de aceptación auditables.

\textbf{OE1. Desplegar infraestructura de serving con vLLM.}\\
Configurar y ejecutar modelos como servidores de inferencia en MareNostrum~5 usando vLLM.\\
\textit{Indicador:} servidor operativo sin errores críticos (caída, HTTP~5xx sostenidos, OOM).\\
\textit{Método:} revisión de logs del servidor (\texttt{vllm.log}).\\
\textit{Condición mínima:} 1~hora de ejecución continua tras warmup de 5~minutos.\\
\textit{Umbral de éxito:} 0~errores críticos en la ventana de medida.

\textbf{OE2. Diseñar generador de peticiones y recolección de datos.}\\
Implementar scripts Python que envíen peticiones, recojan respuestas con metadatos (tiempos, tokens, log~probs) y almacenen resultados.\\
\textit{Indicador:} fichero de resultados completo.\\
\textit{Método:} script \texttt{benchmark\_client.py} exportando a JSON/CSV.\\
\textit{Condición mínima:} 100~peticiones consecutivas sin pérdida de datos.\\
\textit{Umbral de éxito:} 100\% de respuestas registradas con metadatos (timestamp, TTFT, latencia, tokens).

\textbf{OE3. Medir métricas de servicio.}\\
Registrar Time To First Token (TTFT), latencia total, throughput (tokens/s) y memoria GPU para al menos dos configuraciones de modelo.\\
\textit{Indicador:} ficheros de métricas con p50 y p95.\\
\textit{Método:} script de análisis sobre logs; memoria vía \texttt{nvidia-smi}.\\
\textit{Condición mínima:} 3~repeticiones por configuración, warmup de 10~peticiones, misma config de decoding.\\
\textit{Umbral de éxito:} disponer de p50, p95 y throughput medio para $\geq2$ configuraciones.

\textbf{OE4. Evaluar calidad en tareas de razonamiento matemático.}\\
Medir calidad en GSM8K y MATH usando exact match sobre la respuesta final.\\
\textit{Indicador:} métricas de calidad (accuracy/EM) documentadas para ambos benchmarks.\\
\textit{Método:} \texttt{lm-evaluation-harness} o scripts propios comparando con ground truth.\\
\textit{Condición mínima:} evaluación en GSM8K y MATH, misma plantilla de prompt, temperatura y top\_p fijos.\\
\textit{Umbral de éxito:} informe con accuracy/EM para GSM8K y MATH.

\textbf{OE5. Obtener modelos student para la cascada.}\\
La vía principal es obtener un student mediante knowledge distillation del teacher; solo si la ejecución no es viable con la cuota disponible, seleccionar un modelo preentrenado equivalente como contingencia y documentar la decisión.\\
\textit{Indicador:} modelo student desplegado y operativo.\\
\textit{Método:} despliegue con vLLM y prueba de inferencia.\\
\textit{Condición mínima:} al menos 1~student funcional.\\
\textit{Umbral de éxito:} student responde correctamente a 100~peticiones de prueba.

\textbf{OE6. Implementar y evaluar políticas de cascada.}\\
Implementar la Política~B (cascada forzada) como mínimo exigible y la Política~C (cascada por confianza) si los recursos lo permiten. Comparar ambas contra la Política~A (baseline).\\
\textit{Indicador:} reducción de uso del teacher manteniendo calidad.\\
\textit{Método:} comparar políticas en el mismo conjunto de prueba.\\
\textit{Condición mínima:} Política~A y Política~B evaluadas con 3~repeticiones.\\
\textit{Umbral de éxito:} Política~B reduce $\geq$20\% las peticiones al teacher sin degradar calidad más de 5~puntos porcentuales respecto a Política~A.

\subsection{Descripción general de la solución}

La solución se estructura en siete fases técnicas. En la \textbf{Fase~1} se selecciona y despliega el modelo teacher con vLLM, que proporciona un servidor HTTP con PagedAttention~\cite{vllm2023}. En la \textbf{Fase~2} se implementa la generación de peticiones y recolección de datos (tiempos, tokens, log~probs) mediante scripts Python. En la \textbf{Fase~3} se miden métricas de servicio (TTFT, latencia, throughput, memoria GPU). En la \textbf{Fase~4} se evalúa el teacher en los benchmarks seleccionados para establecer la referencia de calidad (Política~A).

\begin{figure}[H]
\centering
\resizebox{\textwidth}{!}{%
\begin{tikzpicture}[
>=Latex,
node distance=8mm and 14mm,
sys/.style={draw, rounded corners=2.2pt, align=center, minimum width=3.1cm, minimum height=9mm, font=\small, thick},
flow/.style={->, thick}
]
\node[sys, fill=yellow!18] (req) {Petición de usuario};
\node[sys, right=of req, fill=orange!20] (router) {Módulo de routing};

\node[sys, above right=10mm and 14mm of router, fill=blue!16] (student) {Student};
\node[sys, below right=10mm and 14mm of router, fill=red!14] (teacher) {Teacher};

\node[sys, right=45mm of router, fill=gray!14, minimum width=3cm] (merge) {Output};
\node[sys, right=of merge, fill=teal!16, minimum width=4.0cm] (metrics) {Registro de métricas\\TTFT, latencia, throughput, EM};

\draw[flow] (req) -- (router);
\draw[flow] (router) -- node[above, sloped, font=\tiny, align=center] {petición\\simple} (student);
\draw[flow] (router) -- node[below, sloped, font=\tiny, align=center] {petición\\compleja} (teacher);
\draw[flow] (student) -- node[right, font=\tiny, align=center] {cascada} (teacher);
\draw[flow] (student.east) -- ++(7mm,0) |- (merge.north west);
\draw[flow] (teacher.east) -- ++(7mm,0) |- (merge.south west);
\draw[flow] (merge) -- (metrics);

\begin{scope}[on background layer]
\node[draw=orange!45, rounded corners=4pt, fill=orange!5, fit=(router)(student)(teacher), inner sep=6pt] {};
\end{scope}
\end{tikzpicture}%
}%resizebox
\caption{Arquitectura operativa del sistema de cascada: enrutado por dificultad, fallback al teacher y trazabilidad de métricas. Fuente: elaboración propia.}
\label{fig:arquitectura-sistema}
\end{figure}

\textbf{Benchmarks seleccionados (GSM8K y MATH).} Se han elegido dos benchmarks de razonamiento matemático~\cite{cobbe2021gsm8k, hendrycks2021math}. GSM8K contiene problemas de aritmética escolar con razonamiento en varios pasos; MATH incluye problemas de competición de mayor dificultad. Ambos permiten medir si el modelo pequeño resuelve problemas sencillos y si el routing detecta cuándo escalar al grande. La métrica es \textit{exact match} (EM) sobre la respuesta final.

En la \textbf{Fase~5} se selecciona al menos un student y se obtiene su línea base aplicando el mismo protocolo de evaluación que al teacher. La planificación prioriza la destilación entre \textit{teacher} y \textit{student} como núcleo metodológico; si esta no resulta viable por restricciones de cuota, se activa la contingencia con un modelo preentrenado equivalente y se documenta explícitamente. En la \textbf{Fase~6} se implementan y evalúan las políticas de cascada. La Política~B (cascada forzada) es el mínimo exigible y consiste en intentar primero con el modelo más pequeño y escalar si no cumple el criterio. La Política~C (cascada por confianza) utiliza un componente predictivo basado en señales como log~probs o entropía para decidir a priori qué modelo servir, evitando escalados innecesarios. La Política~C es la versión avanzada, viable si los recursos lo permiten, pero no es requisito indispensable del éxito del proyecto.

En la \textbf{Fase~7} se consolidan los resultados y se cierra la documentación del proyecto. Esto comprende la preparación de los scripts y configuraciones reproducibles (README, entornos y ejemplos de ejecución), la redacción final de la memoria del TFG y la preparación del material de defensa.

%%%%%%%%%%%%%%%%%%%%%%%%%%%%%%%%%%%%%%%%%%%%%%%%%%%%%%%%%%%%%%%%%
% 5. ALCANCE
%%%%%%%%%%%%%%%%%%%%%%%%%%%%%%%%%%%%%%%%%%%%%%%%%%%%%%%%%%%%%%%%%
\section{Alcance}

\subsection{Elementos incluidos y excluidos}

La Tabla~\ref{tab:scope} detalla qué elementos están dentro y fuera del alcance del proyecto.

\begin{table}[H]
\centering
\caption{Alcance del proyecto: elementos incluidos y excluidos. Fuente: elaboración propia.}
\label{tab:scope}
\small
\begin{tabular}{>{\raggedright\arraybackslash}p{6.5cm} >{\raggedright\arraybackslash}p{6.5cm}}
\toprule
\textbf{Incluido (IN scope)} & \textbf{Excluido (OUT of scope)} \\
\midrule
Despliegue de 1 o 2 modelos teacher con vLLM & Entrenamiento de LLMs desde cero \\
\addlinespace
Obtención de student por destilación entre \textit{teacher} y \textit{student} (con alternativa preentrenada solo como contingencia) & Optimización a nivel de kernel CUDA \\
\addlinespace
Implementación de Política~A (baseline) y Política~B (cascada forzada) & Desarrollo de nuevas arquitecturas de modelos \\
\addlinespace
Exploración de Política~C (cascada por confianza) si recursos lo permiten & Sistema de routing en producción con alta disponibilidad \\
\addlinespace
Medición de métricas de serving (TTFT, latencia, throughput, memoria GPU) & Cobertura exhaustiva de todos los benchmarks existentes \\
\addlinespace
Evaluación de calidad en GSM8K y MATH & Fine-tuning extensivo con múltiples técnicas \\
\addlinespace
Documentación y scripts reproducibles & Interfaz de usuario gráfica \\
\addlinespace
Análisis en entorno BSC con recursos asignados & Despliegue en cloud comercial (AWS, GCP, Azure) \\
\bottomrule
\end{tabular}
\end{table}

\subsection{Entregables}

Los entregables constituyen los resultados verificables del proyecto. Permiten cerrar el alcance de forma auditable y dan soporte a la trazabilidad entre objetivos, tareas y validación. Incluyen documentación técnica y artefactos de código.

\textbf{D1. Informe de configuración.} Documentación del despliegue de vLLM en BSC (entorno, dependencias, scripts).

\textbf{D2. Scripts de generación de peticiones.} Código Python para enviar peticiones y registrar metadatos.

\textbf{D3. Scripts de evaluación de métricas.} Herramientas para calcular métricas de servicio y calidad.

\textbf{D4. Informe comparativo de políticas.} Análisis cuantitativo del equilibrio calidad--eficiencia para A y B, incorporando C si su implementación es viable, y contraste entre \textit{teacher} y \textit{student}.

\textbf{D5. Prototipo de cascada.} Implementación y evaluación de la cascada entre \textit{teacher} y \textit{student} (con contingencia preentrenada, si aplica).

\textbf{D6. Memoria final del TFG.} Documento completo con metodología, experimentos, resultados y conclusiones.

\textbf{D7. Repositorio de código.} Código fuente organizado y documentado en Git.

\subsection{Matriz de trazabilidad}

La Tabla~\ref{tab:trace} relaciona objetivos, entregables, criterios y evidencias verificables.

\begin{table}[H]
\centering
\caption{Trazabilidad OE $\leftrightarrow$ entregables $\leftrightarrow$ evidencia. Fuente: elaboración propia.}
\label{tab:trace}
\footnotesize
\setlength{\tabcolsep}{3pt}
\begin{tabularx}{\textwidth}{>{\raggedright\arraybackslash}p{0.8cm} >{\raggedright\arraybackslash}p{1.8cm} X >{\raggedright\arraybackslash}p{3.8cm}}
\toprule
\textbf{OE} & \textbf{Entregables} & \textbf{Criterio resumido} & \textbf{Evidencia verificable} \\
\midrule
OE1 & D1, D7 & Servidor 1h sin errores críticos & \texttt{vllm.log} sin OOM/5xx \\
OE2 & D2, D7 & 100 peticiones registradas & \texttt{results.json} con 100 filas completas \\
OE3 & D3, D4 & Métricas p50/p95 para $\geq$2 configs & Ficheros CSV con TTFT, latencia, throughput \\
OE4 & D3, D4, D6 & Calidad en GSM8K y MATH & Tabla con accuracy/EM por benchmark \\
OE5 & D1, D4, D7 & Student obtenido por distilación (o contingencia justificada), desplegado y operativo & Log de 100 respuestas del student y registro de decisión técnica \\
OE6 & D5, D6 & Política~B reduce $\geq$20\% uso teacher, $\leq$5pp pérdida & Tabla comparativa políticas A, B (y C si viable) \\
\bottomrule
\end{tabularx}
\end{table}

\subsection{Requisitos operativos y criterios de aceptación}

Dado que este TFG no desarrolla un producto software convencional, sino un pipeline experimental de evaluación, los requisitos funcionales y no funcionales tradicionales no aplican de forma directa. En su lugar, los criterios de aceptación definidos para cada objetivo específico (Sección~\ref{subsec:oe-especificos}) actúan como requisitos operativos del pipeline, estableciendo condiciones cuantitativas que la infraestructura, los scripts y el protocolo experimental deben satisfacer para considerar válido cada entregable. Esta formulación es coherente con la naturaleza experimental del proyecto, donde la calidad se mide por la reproducibilidad y la trazabilidad de los resultados, no por funcionalidades de usuario final.

Los criterios detallados se encuentran en la Sección~\ref{subsec:oe-especificos} junto a cada objetivo específico. La Tabla~\ref{tab:acceptance} ofrece una vista resumida.

\begin{table}[H]
\centering
\caption{Resumen de requisitos operativos y criterios de aceptación. Fuente: elaboración propia.}
\label{tab:acceptance}
\footnotesize
\setlength{\tabcolsep}{4pt}
\begin{tabularx}{\textwidth}{>{\raggedright\arraybackslash}p{1cm} X >{\raggedright\arraybackslash}p{4.5cm}}
\toprule
\textbf{OE} & \textbf{Criterio} & \textbf{Evidencia} \\
\midrule
OE1 & Servidor 1h sin errores críticos (OOM, 5xx) & Log \texttt{vllm.log} \\
OE2 & 100 peticiones registradas con metadatos & \texttt{results.json} \\
OE3 & p50/p95 TTFT, latencia, throughput para $\geq$2 configs & Ficheros CSV \\
OE4 & Accuracy/EM en GSM8K y MATH & Tabla comparativa por benchmark \\
OE5 & Student por distilación (o contingencia justificada) desplegado y operativo (100 respuestas) & Log de respuestas y decisión documentada \\
OE6 & Política~B reduce $\geq$20\% uso teacher, $\leq$5pp pérdida & Tabla comparativa políticas \\
\bottomrule
\end{tabularx}
\end{table}

\subsection{Restricciones y supuestos}

\textbf{Restricciones.} Ejecución en BSC sujeta a políticas SLURM y cuota GPU asignada. Posibles restricciones de red para servicios HTTP. Dependencia de versiones compatibles vLLM/PyTorch/CUDA. Duración limitada (un cuatrimestre).

\textbf{Supuestos.} Acceso a modelos públicos (LLaMA~\cite{touvron2023llama}, Mistral~\cite{jiang2023mistral}) y datasets gratuitos. Compatibilidad vLLM con modelos seleccionados. La destilación entre \textit{teacher} y \textit{student} es la vía principal; si resulta inviable, se activará la contingencia con modelos preentrenados pequeños equivalentes.

\subsection{Riesgos iniciales y mitigación}

La identificación temprana de riesgos forma parte del alcance porque condiciona directamente los objetivos y las decisiones de diseño. La Tabla~\ref{tab:risks} presenta los riesgos identificados con disparadores, responsables y planes de contingencia.

\begin{table}[H]
\centering
\caption{Riesgos, disparadores, responsables y contingencia. Fuente: elaboración propia.}
\label{tab:risks}
\scriptsize
\setlength{\tabcolsep}{2pt}
\begin{tabularx}{\textwidth}{>{\raggedright\arraybackslash}p{2cm} c c >{\raggedright\arraybackslash}p{2.2cm} c >{\raggedright\arraybackslash}X}
\toprule
\textbf{Riesgo} & \textbf{Prob.} & \textbf{Imp.} & \textbf{Disparador} & \textbf{Resp.} & \textbf{Plan B / Contingencia} \\
\midrule
Colas GPU saturadas & Media & Alto & Espera $>$48h & Est. & Ejecutar en baja demanda; priorizar experimentos críticos. \\
\addlinespace
Destilación inviable & Media & Alto & Cuota GPU insuficiente & Est. & Activar contingencia con modelo preentrenado equivalente y documentar intento de destilación. \\
\addlinespace
Señal de confianza no discrimina & Media & Medio & Router no distingue aciertos/fallos & Est. & Probar entropía u otras señales; aceptar Política~B como resultado válido. \\
\addlinespace
Predictor no generaliza & Media & Medio & Precisión $<$60\% en validación & Est. & Simplificar modelo predictivo; priorizar Política~B como resultado mínimo. \\
\addlinespace
Latencia extra del routing & Media & Medio & Overhead $>$10\% de latencia base & Est. & Optimizar predictor; documentar overhead como parte del análisis. \\
\addlinespace
Errores de calibración & Media & Medio & Confianza no correlaciona con acierto & Est. & Aplicar calibración post-hoc; aceptar resultado exploratorio. \\
\addlinespace
Complejidad multi-modelo & Media & Medio & Errores al servir varios modelos & Est. & Simplificar a 2 modelos; serializar si hay conflictos de memoria. \\
\addlinespace
Incompatibilidad software & Baja & Medio & Error vLLM/PyTorch & Est. & Usar versiones alternativas; consultar soporte BSC. \\
\bottomrule
\end{tabularx}
\end{table}

%%%%%%%%%%%%%%%%%%%%%%%%%%%%%%%%%%%%%%%%%%%%%%%%%%%%%%%%%%%%%%%%%
% 6. METODOLOGÍA, DATOS Y RECURSOS (agrupado)
%%%%%%%%%%%%%%%%%%%%%%%%%%%%%%%%%%%%%%%%%%%%%%%%%%%%%%%%%%%%%%%%%
\section{Metodología, datos y recursos del proyecto}
\label{sec:metodologia-trabajo}

\subsection{Enfoque metodológico}

Se seguirá \textbf{Scrum adaptado} con sprints de dos semanas, complementado con un \textbf{tablero Kanban} (To~do / Doing / Done) para el seguimiento diario del flujo de tareas~\cite{schwaber2020scrum}. Scrum adaptado encaja en un TFG experimental porque el ciclo de iterar, medir y decidir es inherente al trabajo, ya que los resultados de cada experimento condicionan los siguientes pasos y las iteraciones cortas permiten incorporar hallazgos antes de comprometer más recursos. El tablero Kanban proporciona visibilidad continua sobre el estado de cada tarea e identifica bloqueos (por ejemplo, esperas en colas GPU) sin necesidad de esperar a la revisión de sprint. Además, al tratarse de un proyecto unipersonal, Scrum se adapta simplificando las ceremonias al mínimo operativo que aporta valor real.

\textbf{Ceremonias.}
\begin{itemize}[nosep,leftmargin=*]
\item \textbf{Sprint Planning (inicio de sprint):} seleccionar tareas del Product Backlog derivado del WBS de este documento, definir el objetivo del sprint y establecer qué métricas o evidencias validan su cumplimiento.
\item \textbf{Seguimiento semanal:} revisar bloqueos y avance contra el Sprint Backlog. Las reuniones quincenales con la codirección actúan como punto de control externo; en las semanas intermedias, la revisión es individual sobre el tablero Kanban.
\item \textbf{Sprint Review (final de sprint):} presentar los resultados medibles del sprint (tablas, CSV, logs, scripts) y compararlos con los criterios de aceptación definidos en los objetivos específicos (Sección~\ref{subsec:oe-especificos}).
\item \textbf{Retrospectiva:} registrar qué aspectos del proceso o de los experimentos se pueden mejorar y ajustar el backlog en consecuencia.
\end{itemize}

\textbf{Artefactos.}
\begin{itemize}[nosep,leftmargin=*]
\item \textbf{Product Backlog:} lista priorizada de tareas derivada del WBS definido en este documento. Cada entrada está ligada a un objetivo específico (OE1--OE6) y a un entregable (D1--D7).
\item \textbf{Sprint Backlog:} subconjunto de tareas seleccionadas para el sprint en curso, con estado visible en el tablero Kanban.
\item \textbf{Evidencias:} resultados reproducibles generados por cada tarea: scripts, configuraciones, logs de ejecución, tablas de métricas y gráficos.
\item \textbf{CHANGELOG / registro de decisiones:} fichero \texttt{CHANGELOG.md} que documenta qué se cambió y por qué en cada iteración, asegurando la trazabilidad de las decisiones de diseño.
\end{itemize}

\textbf{Definition of Done.} Una tarea se considera terminada cuando cumple todos los criterios siguientes.
\begin{enumerate}[nosep,leftmargin=*]
\item Código, scripts o configuración commitados en el repositorio Git y ejecución reproducible desde cero.
\item Métricas o resultados guardados en formato persistente (CSV, JSON o log) y referenciados desde la tarea.
\item Breve interpretación del resultado documentada (en el issue o en \texttt{CHANGELOG.md}).
\item Checklist de la tarea marcada como completada y vinculada al milestone correspondiente.
\end{enumerate}

\textbf{Conexión con la planificación temporal.} El diagrama de Gantt y el WBS definidos en este documento constituyen el macro-plan del proyecto, organizando las fases (F1--F7) y tareas transversales (TG) con sus dependencias y fechas objetivo. Los sprints de dos semanas ejecutan ese plan: en cada Sprint Planning se seleccionan tareas del WBS según prioridad y dependencias. Si se producen desviaciones, se replanifica el backlog de los sprints siguientes manteniendo los objetivos y restricciones del proyecto.

\subsection{Herramientas de soporte y seguimiento}

\textbf{Control de versiones.} Git se utiliza como herramienta de control de versiones. Todo el código, los scripts de experimentación y la documentación se versionan en un repositorio con commits atómicos y mensajes descriptivos. Git no es la metodología del proyecto, sino el soporte técnico que garantiza la trazabilidad.

\textbf{Seguimiento de tareas.} El avance se registra mediante issues y milestones en el repositorio, cada uno ligado a un objetivo específico o entregable del WBS. Un checklist de entregables permite verificar en cada punto de control qué tareas están completadas y cuáles presentan desviación.

\textbf{Artefactos de validación.} Para cada objetivo específico se define una evidencia verificable (logs, ficheros CSV, tablas) que debe estar en el repositorio antes de cerrar la tarea correspondiente. El fichero \texttt{CHANGELOG.md} documenta decisiones de diseño y cambios relevantes.

La planificación temporal detallada, con diagrama de Gantt, hitos, dependencias y asignación de recursos, se desarrolla en la Sección~\ref{sec:planificacion-wbs-control} de este documento.

\subsection{Protocolo de medida y comparabilidad}

Para garantizar reproducibilidad se aplicará el siguiente protocolo en todos los experimentos. Este mismo protocolo se aplica a la evaluación de modelos individuales (teacher, students) y a la evaluación de las políticas de cascada.

\begin{figure}[H]
\centering
\begin{tikzpicture}[
>=Latex,
node distance=6.5mm,
step/.style={draw, rounded corners=2pt, align=center, minimum width=8.8cm, minimum height=9mm, font=\footnotesize, thick}
]
\node[step, fill=blue!14] (s1) {1) Preparación de datos y prompts (GSM8K, MATH)};
\node[step, below=of s1, fill=cyan!14] (s2) {2) Ejecución baseline con teacher (Política A)};
\node[step, below=of s2, fill=green!16] (s3) {3) Evaluación del student (baseline o destilado)};
\node[step, below=of s3, fill=orange!18] (s4) {4) Aplicación de routing (Políticas B y C)};
\node[step, below=of s4, fill=violet!14] (s5) {5) Medida controlada (warmup + 3 repeticiones)};
\node[step, below=of s5, fill=gray!18] (s6) {6) Análisis final de calidad-coste (TTFT, p95, throughput, EM)};

\draw[->, thick] (s1) -- (s2);
\draw[->, thick] (s2) -- (s3);
\draw[->, thick] (s3) -- (s4);
\draw[->, thick] (s4) -- (s5);
\draw[->, thick] (s5) -- (s6);
\end{tikzpicture}
\caption{Pipeline experimental simplificado en flujo secuencial vertical para garantizar trazabilidad y comparabilidad. Fuente: elaboración propia.}
\label{fig:pipeline-experimental}
\end{figure}

\begin{enumerate}[nosep,leftmargin=*]
\item \textbf{Warmup:} descartar las primeras 10~peticiones antes de medir.
\item \textbf{Ventana de medida:} mínimo 100~peticiones o 30~min.
\item \textbf{Repeticiones:} 3~ejecuciones independientes; reportar media y dispersión.
\item \textbf{Métricas:} p50, p95 de TTFT y latencia; throughput medio (tokens/s).
\item \textbf{Control de variables:} misma plantilla de prompt, temperatura~0, top\_p~1, max\_tokens fijo.
\item \textbf{Comparabilidad:} teacher, students y políticas, evaluados en idénticas condiciones.
\end{enumerate}

\subsection{Datos globales del proyecto}

Este TFG se desarrolla durante 19 semanas, del 9 de febrero al 18 de junio de 2026, y suma un total de 540 horas. De ellas, 450 corresponden al trabajo técnico (fases F1--F7 y tareas transversales TG, equivalentes a 15 ECTS) y las 90 restantes a la gestión del proyecto dentro de la asignatura GEP, contabilizadas aparte y sin solapamiento. El cierre técnico está previsto para el 28 de mayo, la entrega de la memoria para el 18 de junio y la defensa para la última semana de junio.

Mi dedicación prevista es de 5--6 horas diarias, cinco días a la semana, lo que supone unas 28 horas semanales. Es un ritmo asumible porque el TFG es mi actividad principal del semestre y no curso otras asignaturas de peso. Cada lunes reviso la planificación semanal y reparto bloques entre experimentación GPU y redacción.

El periodo del 29 de mayo al 18 de junio queda como holgura para imprevistos, revisión de la memoria y preparación de la defensa.

\subsection{Recursos humanos}

El único recurso humano del proyecto soy yo, Marc Romera Peral. Las 540 horas del proyecto son íntegramente mías, repartidas entre las fases técnicas (450\,h) y la gestión del proyecto (90\,h). Las reuniones de seguimiento quincenales (TG.1, 12\,h en total) se contabilizan como mi tiempo de dedicación.

Aunque el proyecto lo realiza una sola persona, el trabajo abarca competencias muy distintas, desde la ingeniería de infraestructura HPC hasta el diseño experimental propio de la ciencia de datos, pasando por la implementación de modelos y la redacción técnica. Para reflejar esta diversidad, las horas se distribuyen entre cuatro perfiles profesionales representativos de un proyecto de ML en entornos de supercomputación (Tabla~\ref{tab:roles}). La descomposición responde a las categorías de trabajo que el WBS pone de manifiesto y permite dimensionar correctamente el esfuerzo dedicado a cada área de competencia.

{\footnotesize
\begin{tabularx}{\textwidth}{l r r >{\raggedright\arraybackslash}X}
\caption{Distribución de horas por perfil profesional. Elaboración propia.}\label{tab:roles}\\
\toprule
\textbf{Perfil (rol)} & \textbf{Horas} & \textbf{\%} & \textbf{Tareas principales} \\
\midrule
\endfirsthead
\toprule
\textbf{Perfil (rol)} & \textbf{Horas} & \textbf{\%} & \textbf{Tareas principales} \\
\midrule
\endhead
\bottomrule
\endlastfoot
Cap de projecte (CP) & 140 & 25,9 & Planificación, seguimiento y control del proyecto (GP, 90\,h), reuniones quincenales de seguimiento (TG.1, 12\,h), coordinación y contribución a la documentación continua (TG.2, 10\,h) y estructura general de la memoria y defensa (T7.1, T7.5, T7.6) \\
Data scientist (DS) & 120 & 22,2 & Diseño experimental, selección y evaluación de modelos en benchmarks, análisis de métricas de calidad y rendimiento, documentación continua de experimentos y resultados (TG.2, 8\,h) y redacción de los capítulos de resultados (T7.3) \\
ML engineer (ML) & 172 & 31,9 & Implementación de pipelines, instrumentación de métricas, entrenamiento, destilación, desarrollo de las políticas de routing, documentación continua de código y decisiones técnicas (TG.2, 8\,h) y redacción de los capítulos de routing y discusión (T7.4) \\
DevOps-HPC (DV) & 108 & 20,0 & Aprovisionamiento de MareNostrum~5, scripts SLURM, despliegue de vLLM, monitorización GPU, documentación continua de infraestructura y entorno (TG.2, 4\,h), mantenimiento del repositorio y redacción del capítulo de metodología y entorno (T7.2) \\
\midrule
\textbf{Total} & \textbf{540} & \textbf{100} & \\
\end{tabularx}
}

\noindent\textit{\footnotesize Nota: todas las horas corresponden al autor. Los perfiles reflejan el tipo de trabajo realizado en cada tarea y sirven como base de valoración económica en la sección de presupuesto de este documento. La documentación continua (TG.2) se reparte entre los cuatro perfiles: CP coordina y mantiene el registro (10\,h), DS documenta experimentos y resultados (8\,h), ML documenta código y decisiones técnicas (8\,h) y DV documenta infraestructura y entorno (4\,h). Las tareas de F7 se distribuyen igualmente entre perfiles según la naturaleza de cada capítulo.}

\subsection{Recursos materiales y tecnológicos}

La infraestructura principal del proyecto es MareNostrum~5~\cite{mn5}, el supercomputador del Barcelona Supercomputing Center (BSC). El clúster proporciona acceso a GPUs NVIDIA Hopper (H100) a través del gestor de colas SLURM~\cite{slurm}, y la cuota asignada al proyecto garantiza disponibilidad suficiente para los experimentos planificados.

El motor de inferencia utilizado es vLLM~\cite{vllm2023}, que permite servir modelos de lenguaje con alta eficiencia. La evaluación de calidad se apoya en los benchmarks matemáticos GSM8K~\cite{cobbe2021gsm8k} y MATH~\cite{hendrycks2021math}, seleccionados por ser estándar en la comunidad y ofrecer métricas objetivas de razonamiento.

Para el control de versiones y la documentación se utilizan Git y LaTeX respectivamente, lo que asegura trazabilidad completa y reproducibilidad. Todos estos recursos están disponibles desde el inicio del proyecto, ya sea porque el BSC los proporciona directamente o porque son herramientas de acceso libre, lo que confirma la viabilidad técnica de la planificación. El proyecto no requiere financiación externa ni adquisición de recursos adicionales, dado que toda la infraestructura computacional está concedida y el resto de herramientas son de acceso libre.

Además de la infraestructura computacional, el proyecto utiliza recursos de espacio de trabajo que se valorarán económicamente en la sección de presupuesto de este documento:

\begin{itemize}[nosep]
\item \textbf{Espacio de trabajo:} domicilio particular, utilizado como oficina durante las 19 semanas del proyecto. Se imputará una fracción proporcional del coste mensual.
\item \textbf{Internet:} conexión doméstica de fibra óptica, necesaria para el acceso remoto a MareNostrum~5 y la gestión del repositorio. Se imputará la fracción proporcional del coste mensual.
\item \textbf{Teléfono:} uso puntual para comunicaciones. Se imputará un coste simbólico mensual.
\end{itemize}

%%%%%%%%%%%%%%%%%%%%%%%%%%%%%%%%%%%%%%%%%%%%%%%%%%%%%%%%%%%%%%%%%
\section{Planificación temporal, WBS y control}
\label{sec:planificacion-wbs-control}

\subsection{Desglose de tareas (WBS)}

He organizado el proyecto en siete fases técnicas (F1--F7) y un bloque transversal (TG). La estructura es deliberadamente incremental, de modo que cada fase construye sobre la anterior y permite detectar desviaciones antes de que se propaguen~\cite{gpi}. Esta secuenciación refleja la lógica natural del trabajo, empezando por la infraestructura y terminando con el routing y la redacción. La Tabla~\ref{tab:wbs} recoge el desglose completo.

{\scriptsize
\setlength{\tabcolsep}{3.5pt}
\renewcommand{\arraystretch}{0.93}
\begin{tabularx}{\textwidth}{>{\ttfamily}l >{\raggedright\arraybackslash}p{4.5cm} r c >{\raggedright\arraybackslash}p{1.1cm} >{\raggedright\arraybackslash}X}
\caption{Desglose estructurado de tareas del proyecto (WBS). Elaboración propia.}\label{tab:wbs}\\
\toprule
\textbf{Cód.} & \textbf{Tarea} & \textbf{h} & \textbf{Dep.} & \textbf{Rol} & \textbf{Recursos materiales} \\
\midrule
\endfirsthead
\multicolumn{6}{l}{\textit{(continuación de la Tabla~\ref{tab:wbs})}} \\
\toprule
\textbf{Cód.} & \textbf{Tarea} & \textbf{h} & \textbf{Dep.} & \textbf{Rol} & \textbf{Recursos materiales} \\
\midrule
\endhead
\midrule
\multicolumn{6}{r}{\textit{Continúa en la página siguiente}} \\
\endfoot
\bottomrule
\endlastfoot
T1.1 & Acceso y cuenta BSC & 4 & -- & DV & MN5 \\
T1.2 & Entorno y scripts SLURM base & 10 & T1.1 & DV & MN5, SLURM \\
T1.3 & Instalar vLLM y dependencias & 12 & T1.2 & DV & MN5, vLLM, H100 \\
T1.4 & Preparar teacher y cachés & 10 & T1.3 & ML, DV & MN5, H100 \\
T1.5 & Desplegar API teacher y healthchecks & 12 & T1.4 & DV & MN5, vLLM, H100 \\
\multicolumn{2}{r}{\textit{Subtotal F1}} & \textit{48} & & & \\
\addlinespace[2pt]
T2.1 & Definir esquema req-resp y logs & 6 & T1.5 & DS, ML & MN5, Python \\
T2.2 & Impl.\ cliente benchmark y conc. & 12 & T2.1 & ML & MN5, Python \\
T2.3 & Muestreador de prompts y norm. & 6 & T2.2 & ML & MN5, Python \\
T2.4 & Validación 10\,000 peticiones & 8 & T2.3 & ML & MN5, vLLM, H100 \\
\multicolumn{2}{r}{\textit{Subtotal F2}} & \textit{32} & & & \\
\addlinespace[2pt]
T3.1 & Instrumentación latencia y tokens & 8 & T2.4 & ML, DV & MN5, Python \\
T3.2 & Experimentos de carga (sweep) & 14 & T3.1 & ML & MN5, H100 \\
T3.3 & Registro memoria y uso GPU & 6 & T3.2 & DV & MN5, H100, nvidia-smi \\
T3.4 & Agregación p50/p95 y tablas & 8 & T3.3 & DS & MN5, Python \\
\multicolumn{2}{r}{\textit{Subtotal F3}} & \textit{36} & & & \\
\addlinespace[2pt]
T4.1 & Preparar eval.\ GSM8K y MATH & 8 & T2.4 & DS & MN5, GSM8K, MATH \\
T4.2 & Ejecutar evaluación teacher & 14 & T4.1 & ML & MN5, H100, vLLM \\
T4.3 & Análisis errores y baseline calidad & 12 & T4.2 & DS & MN5, Python \\
\multicolumn{2}{r}{\textit{Subtotal F4}} & \textit{34} & & & \\
\addlinespace[2pt]
T5.1 & Selección student candidato & 6 & T4.3 & DS & MN5 \\
T5.2 & Baseline student (mismo prot.) & 12 & T5.1 & ML & MN5, H100, vLLM \\
T5.3 & Decisión viabilidad destilación & 6 & T5.2 & DS, CP & MN5 \\
T5.4 & Preparar pipeline destilación & 12 & T5.3 & ML & MN5, H100, vLLM \\
T5.5 & Ejecutar destilación o alternativa & 22 & T5.4 & ML & MN5, H100, datasets \\
T5.6 & Evaluar student destilado & 12 & T5.5 & DS, ML & MN5, H100 \\
\multicolumn{2}{r}{\textit{Subtotal F5}} & \textit{70} & & & \\
\addlinespace[2pt]
T6.1 & Definir criterio quality gate & 6 & T4.3,T5.6 & DS & MN5 \\
T6.2 & Impl.\ Política B (cascada forc.) & 14 & T6.1 & ML & MN5, H100, Python \\
T6.3 & Evaluar Política B vs A & 12 & T6.2 & DS, ML & MN5, H100 \\
T6.4 & Preparar features y dataset pred. & 10 & T6.3 & DS, ML & MN5, Python \\
T6.5 & Entrenar predictor (Política C) & 16 & T6.4 & ML & MN5, H100, Python \\
T6.6 & Evaluar Política C vs A y B & 10 & T6.5 & DS, ML & MN5, H100 \\
\multicolumn{2}{r}{\textit{Subtotal F6}} & \textit{68} & & & \\
\addlinespace[2pt]
T7.1 & Esqueleto memoria y reprod. & 8 & T2.4 & CP & LaTeX, Git \\
T7.2 & Redacción metodología y entorno & 12 & T4.3 & DV & LaTeX, Git \\
T7.3 & Redacción resultados teacher/stud. & 16 & T5.6 & DS & LaTeX, Git \\
T7.4 & Redacción routing y discusión & 16 & T6.6 & ML & LaTeX, Git \\
T7.5 & Presentación y guion defensa & 10 & T7.4 & CP & LaTeX, Beamer \\
T7.6 & Revisión final y entrega & 10 & T7.5 & CP & LaTeX, Git \\
\multicolumn{2}{r}{\textit{Subtotal F7}} & \textit{72} & & & \\
\addlinespace[2pt]
TG.1 & Reuniones de seguimiento & 12 & -- & CP & -- \\
TG.2 & Documentación continua & 30 & -- & CP, DS, ML, DV & Git, Markdown \\
TG.3 & Gestión repositorio y versionado & 48 & -- & DV & Git \\
\multicolumn{2}{r}{\textit{Subtotal TG}} & \textit{90} & & & \\
\addlinespace[2pt]
GP & Gestión del proyecto (asignatura) & 90 & -- & CP & LaTeX \\
\midrule
\multicolumn{2}{r}{\textit{Total TFG (F1--F7 + TG)}} & \textit{450} & & & \\
\multicolumn{2}{r}{\textit{Total gestión del proyecto}} & \textit{90} & & & \\
\multicolumn{2}{r}{\textbf{TOTAL GLOBAL}} & \textbf{540} & & & \\
\end{tabularx}
}

\begin{center}
\fcolorbox{gray!30}{gray!5}{\parbox{0.85\textwidth}{\scriptsize\centering
CP = Cap de projecte \enspace$\cdot$\enspace DS = Data scientist \enspace$\cdot$\enspace ML = ML engineer \enspace$\cdot$\enspace DV = DevOps-HPC\\[1pt]
MN5 = MareNostrum~5 \enspace$\cdot$\enspace H100 = GPUs NVIDIA Hopper \enspace$\cdot$\enspace vLLM = motor de inferencia \enspace$\cdot$\enspace SLURM = gestor de colas HPC\\[1pt]
GSM8K/MATH = benchmarks matemáticos}}
\end{center}

Las estimaciones siguen un enfoque bottom-up~\cite{mcconnell06}, descomponiendo cada fase en subtareas con entregable concreto y asignando horas con el apoyo de pruebas preliminares y la validación en las revisiones de seguimiento. Las tareas con mayor incertidumbre, T5.5 (destilación, 22\,h) y T6.5 (predictor, 16\,h), incorporan un margen interno superior al 20\,\% para absorber variaciones en convergencia y tiempos de cola. El total de 450\,h es coherente con los 15 ECTS en 19 semanas~\cite{gep23}, y la holgura de tres semanas (aproximadamente 60\,h útiles) cubre el peor caso acumulado de riesgos (+44\,h).

\subsubsection{Descripción de las fases}

\textbf{F1. Infraestructura y despliegue del teacher (48\,h).} Comprende desde la obtención de la cuenta en MareNostrum~5 hasta tener la API del modelo teacher respondiendo con healthchecks operativos, incluyendo la configuración del entorno, los scripts SLURM y la instalación de vLLM.

\textbf{F2. Cliente de peticiones (32\,h).} Desarrollo del cliente de benchmarking con esquema request-response, módulo de concurrencia y muestreo de prompts, validado con 10\,000 peticiones de extremo a extremo.

\textbf{F3. Métricas de servicio (36\,h).} Instrumentación de latencia, throughput y consumo de GPU bajo distintos niveles de carga. El entregable son las tablas de percentiles (p50, p95) que caracterizan el rendimiento del teacher en condiciones realistas.

\textbf{F4. Evaluación del teacher, Política~A (34\,h).} Evaluación de calidad en los benchmarks GSM8K y MATH. Los resultados constituyen la línea base de calidad contra la que se compararán las políticas de routing posteriores.

\textbf{F5. Student y destilación (70\,h).} La fase más pesada del proyecto. Incluye seleccionar el student, obtener su baseline, decidir si la destilación es viable con la cuota disponible, ejecutarla y evaluar el resultado. T5.3 actúa como punto de decisión, de forma que si la destilación no es factible se recurre a un modelo preentrenado equivalente.

\textbf{F6. Routing inteligente (68\,h).} Implementación y evaluación de las Políticas~B (cascada forzada) y~C (predictor automático). Es la fase que aporta la contribución diferencial del TFG, al decidir cuándo usar el teacher y cuándo el student para optimizar coste y calidad.

\textbf{F7. Documentación y defensa (72\,h).} Redacción progresiva de la memoria, presentación y revisión final. Las subtareas están ligadas a los entregables de cada fase para no dejar toda la escritura para el final.

\textbf{TG. Tareas transversales (90\,h).} Reuniones quincenales (12\,h), documentación continua (30\,h) y gestión del repositorio Git (48\,h), repartidas a lo largo de las 19 semanas.

\subsubsection{Distribución de las horas de gestión del proyecto (GP)}

Las 90 horas de gestión del proyecto se concentran en el primer mes y se contabilizan de forma independiente a las 450 horas técnicas. Durante ese periodo se lleva a cabo todo el trabajo de definición y planificación que sienta las bases del proyecto antes de que la ejecución técnica alcance su máxima intensidad. Esto incluye delimitar el alcance y el contexto del TFG, elaborar la planificación temporal con su WBS y su Gantt, estimar el presupuesto y los recursos necesarios, realizar el análisis de sostenibilidad y el impacto social, e identificar los riesgos con sus planes de mitigación.

Además de la planificación en sí, estas horas cubren la redacción y revisión de los informes parciales de gestión que documentan cada uno de esos aspectos, así como la integración final de toda la documentación de gestión en un documento cohesionado. Es un esfuerzo intensivo pero acotado en el tiempo, equivalente a unas 22 horas semanales durante las cuatro primeras semanas, que permite arrancar las fases técnicas con una base sólida y bien definida~\cite{gep23}.

\subsection{Planificación temporal y camino crítico}

La Figura~\ref{fig:gantt} muestra el diagrama de Gantt generado con ProjectLibre a partir de la Tabla~\ref{tab:wbs}. Este diagrama funciona como \textbf{línea base temporal} del proyecto, a partir del cual comparo semanalmente las horas ejecutadas contra lo planificado~\cite{pmbok,iso21502}. Si la desviación acumulada supera el 15\,\%, convoco reunión extraordinaria de seguimiento para decidir si replanificar o activar algún plan alternativo.

La ejecución de este plan se vehicula mediante \textbf{sprints de dos semanas} dentro del marco Scrum adaptado descrito en la Sección~\ref{sec:metodologia-trabajo}. En cada Sprint Planning se seleccionan tareas del WBS según prioridad y dependencias, y el progreso se controla mediante issues, milestones y checklists vinculados a los hitos de la Tabla~\ref{tab:hitos}.

\begin{figure}[H]
\centering
\vspace{-4pt}
\includegraphics[width=\textwidth]{gannt.png}
\vspace{-6pt}
\caption{Diagrama de Gantt del proyecto (09/02/2026 -- 18/06/2026). Elaboración propia con ProjectLibre.}
\label{fig:gantt}
\end{figure}

La planificación no es secuencial pura, ya que TG.2 y TG.3 se reparten a lo largo de todo el proyecto y TG.1 se ejecuta quincenalmente, aprovechando las esperas en colas GPU para documentación. El \textbf{camino crítico} recorre la cadena F1, F2, F3, F4, F5, F6 y la parte final de F7 (T7.4 en adelante, redacción del routing y defensa). Esta secuencia suma 288 horas de las 450 técnicas y abarca aproximadamente 16 de las 19 semanas del proyecto, ya que la dedicación semanal se reparte entre el trabajo crítico y las tareas transversales que conviven con él. Cualquier retraso en alguna de estas fases se traslada íntegramente a la fecha de entrega, sin que exista holgura intermedia entre ellas. F7 aporta el principal solapamiento de la planificación, ya que cada subtarea de redacción arranca en cuanto termina la fase técnica correspondiente (T7.1 tras F2, T7.2 tras F4, T7.3 tras F5), permitiendo avanzar en documentación sin esperar a que acabe toda la experimentación.

Los hitos verificables del proyecto se recogen en la Tabla~\ref{tab:hitos} de forma estructurada para facilitar el seguimiento.

\needspace{10\baselineskip}
{\footnotesize
\setlength{\tabcolsep}{5pt}
\renewcommand{\arraystretch}{1.05}
\begin{tabularx}{\textwidth}{c c >{\raggedright\arraybackslash}X}
\caption{Hitos del proyecto con fechas objetivo y entregables verificables. Elaboración propia.}\label{tab:hitos}\\
\toprule
\textbf{Hito} & \textbf{Fecha} & \textbf{Entregable verificable} \\
\midrule
\endfirsthead
\toprule
\textbf{Hito} & \textbf{Fecha} & \textbf{Entregable verificable} \\
\midrule
\endhead
\bottomrule
\endlastfoot
H1 & 27/02/2026 & API del teacher desplegada con healthchecks operativos en MN5 \\
H2 & 13/03/2026 & Cliente de benchmarking validado con 10\,000 peticiones \\
H3 & 27/03/2026 & Tablas de métricas p50/p95 de latencia y throughput \\
H4 & 10/04/2026 & Evaluación completa del teacher en GSM8K y MATH (Política~A) \\
H5 & 01/05/2026 & Student destilado evaluado contra baseline \\
H6 & 22/05/2026 & Comparativa de políticas de routing A y B (y C si es viable) \\
H7 & 28/05/2026 & Cierre técnico con todos los experimentos reproducibles \\
H8 & 18/06/2026 & Memoria final entregada y presentación preparada \\
\end{tabularx}
}

\subsection{Gestión del riesgo}

He identificado cuatro riesgos principales, cada uno con un criterio de activación cuantitativo y un plan de mitigación concreto. La Tabla~\ref{tab:riesgos} los resume y a continuación desarrollo el análisis de cada uno.

\needspace{6\baselineskip}
{\footnotesize
\setlength{\tabcolsep}{3pt}
\renewcommand{\arraystretch}{1.08}
\begin{tabularx}{\textwidth}{c >{\raggedright}p{1.3cm} c c c c >{\raggedright}p{2cm} r >{\raggedright\arraybackslash}X}
\caption{Registro de riesgos ordenado por nivel de criticidad, con criterios de activación y planes de mitigación. Elaboración propia.}\label{tab:riesgos}\\
\toprule
\textbf{Cód.} & \textbf{Riesgo} & \textbf{Prob.} & \textbf{Imp.} & \textbf{Nivel} & \textbf{Fases} & \textbf{Activación} & \textbf{+h (rol)} & \textbf{Mitigación} \\
\midrule
\endfirsthead
\toprule
\textbf{Cód.} & \textbf{Riesgo} & \textbf{Prob.} & \textbf{Imp.} & \textbf{Nivel} & \textbf{Fases} & \textbf{Activación} & \textbf{+h (rol)} & \textbf{Mitigación} \\
\midrule
\endhead
\bottomrule
\endlastfoot
R1 & Colas GPU saturadas & Media & Alto & Alto & F3--F6 & $>$48\,h $\times$ 3 ciclos consec. & +20 (ML) & Jobs nocturnos; reducir repeticiones \\
R2 & Destilación inviable & Media & Alto & Alto & F5 & Cuota agotada, 2 intentos & +6 (ML) & Modelo preentrenado equivalente \\
R3 & Predictor no generaliza & Media & Medio & Medio & F6 & Acc.\,$<$60\,\% tras 3 rondas & +10 (ML) & Simplificar señales; Política~B \\
R4 & Latencia routing & Media & Medio & Medio & F6 & $>$10\,\% en 2 exp.\ indep. & +8 (ML) & Optimizar router; priorizar A+B \\
\end{tabularx}
}

\textbf{R1. Colas GPU saturadas.} MareNostrum~5 es un recurso compartido, así que la congestión es un escenario probable. Considero la probabilidad media porque el proyecto tiene cuota asignada, pero el impacto es alto dado que F3--F6 dependen directamente del acceso a GPU. Si se activa el criterio, reprogramo los jobs a franjas nocturnas y, en última instancia, reduzco repeticiones experimentales (+20\,h de ML engineer). Es el riesgo que más horas puede consumir, pero queda dentro de la holgura.

\textbf{R2. Destilación inviable.} Si la cuota no permite completar el entrenamiento del student tras dos intentos en días distintos, selecciono un modelo preentrenado de la misma familia y valido su compatibilidad con el pipeline, tal como contempla T5.3. El impacto es alto sobre F5, pero el plan alternativo mantiene el alcance completo del TFG (+6\,h de ML engineer).

\textbf{R3. Predictor no generaliza.} Si la accuracy del predictor de Política~C no supera el 60\,\% tras tres rondas de ajuste, simplifico las señales de entrada o presento Política~B como resultado principal. El impacto es medio, puesto que F6 ya contempla Política~B como alternativa viable y la reducción de alcance sería parcial (+10\,h de ML engineer).

\textbf{R4. Latencia extra del routing.} Si el overhead del router supera el 10\,\% de la latencia base en dos experimentos independientes, optimizo la implementación o, en caso extremo, priorizo Políticas~A y~B en la memoria. Impacto medio (+8\,h de ML engineer).

La tabla está ordenada por nivel de criticidad combinado (probabilidad $\times$ impacto) y, dentro de un mismo nivel, por coste potencial en horas. R1 encabeza la lista porque, además de presentar nivel alto, es el riesgo con mayor impacto temporal (+20\,h) y el que afecta a más fases simultáneamente (F3 a F6). R2 ocupa el segundo lugar por su nivel alto, aunque su coste en horas es menor gracias al plan alternativo previsto en T5.3. Los riesgos R3 y R4, ambos de nivel medio, se ordenan por horas descendentes y afectan exclusivamente a F6. Esta priorización orienta la vigilancia semanal, concentrando la atención en los riesgos que más condicionan la holgura disponible~\cite{pmbok}. La Figura~\ref{fig:matriz-pi} representa visualmente esta distribución.

\begin{figure}[H]
\centering
\resizebox{0.52\textwidth}{!}{%
\begin{tikzpicture}
\definecolor{cellhigh}{RGB}{190,210,230}
\definecolor{cellmed}{RGB}{222,234,244}
\definecolor{cellempty}{RGB}{242,242,242}
\fill[cellhigh] (0,1.6) rectangle (3.2,3.2);
\fill[cellempty] (3.2,1.6) rectangle (6.4,3.2);
\fill[cellmed] (0,0) rectangle (3.2,1.6);
\fill[cellempty] (3.2,0) rectangle (6.4,1.6);
\draw[gray!70, line width=0.6pt] (0,0) rectangle (6.4,3.2);
\draw[gray!70, line width=0.6pt] (3.2,0) -- (3.2,3.2);
\draw[gray!70, line width=0.6pt] (0,1.6) -- (6.4,1.6);
\node[font=\large\bfseries] at (1.6,2.4) {R1,~R2};
\node[font=\large\bfseries] at (1.6,0.8) {R3,~R4};
\node[font=\normalsize, gray] at (4.8,2.4) {--};
\node[font=\normalsize, gray] at (4.8,0.8) {--};
\node[font=\normalsize, anchor=east] at (-0.25,2.4) {Alto};
\node[font=\normalsize, anchor=east] at (-0.25,0.8) {Medio};
\node[font=\normalsize, anchor=north] at (1.6,-0.3) {Media};
\node[font=\normalsize, anchor=north] at (4.8,-0.3) {Alta};
\node[font=\normalsize\bfseries] at (3.2,3.75) {Probabilidad};
\node[font=\normalsize\bfseries, rotate=-90] at (7.05,1.6) {Impacto};
\end{tikzpicture}%
}
\caption{Matriz probabilidad--impacto con ubicación de los riesgos identificados. Elaboración propia.}\label{fig:matriz-pi}
\end{figure}

La suma de impactos en el peor caso (+44\,h) es inferior a la holgura disponible (aproximadamente 60\,h), lo que garantiza la viabilidad temporal incluso activando todos los planes simultáneamente~\cite{pmbok}. Esto no significa que los riesgos sean despreciables, sino que la planificación los ha absorbido de forma explícita. La Figura~\ref{fig:holgura} visualiza esta relación.

\begin{figure}[H]
\centering
\begin{tikzpicture}
\definecolor{barblue}{RGB}{70,130,180}
\definecolor{barlight}{RGB}{215,228,240}
\fill[barlight] (0,0) rectangle (10,0.7);
\fill[barblue] (0,0) rectangle (7.33,0.7);
\draw[gray] (0,0) rectangle (10,0.7);
\draw[gray, densely dashed] (7.33,0) -- (7.33,0.7);
\node[font=\footnotesize, anchor=south] at (3.67,0.8) {Impacto máximo acumulado (44\,h)};
\node[font=\footnotesize, anchor=south] at (8.67,0.8) {Margen residual (16\,h)};
\node[font=\footnotesize, anchor=north] at (5,-0.15) {Holgura total disponible (60\,h)};
\end{tikzpicture}
\caption{Comparación entre la holgura disponible y el impacto acumulado de riesgos en el peor caso. Elaboración propia.}\label{fig:holgura}
\end{figure}

\subsection{Seguimiento y control}

El seguimiento combina dos mecanismos complementarios. Por un lado, la documentación continua (TG.2), que se actualiza tras cada tarea técnica con README, decisiones de diseño y resultados. Por otro, las reuniones quincenales de seguimiento (TG.1), donde contrasto el avance real frente al plan~\cite{gep23}.

El instrumento central de registro es una hoja de seguimiento semanal mantenida en el repositorio del proyecto. En ella anoto, para cada tarea del WBS, las horas planificadas acumuladas y las horas reales ejecutadas hasta la fecha, junto con el porcentaje de avance estimado y cualquier incidencia relevante. Esta comparación sistemática entre lo previsto y lo ejecutado se inspira en el concepto de valor ganado (\textit{Earned Value}) descrito en el PMBOK~\cite{pmbok}, adaptado a la escala de un TFG individual sin la formalidad completa de un cuadro de mando corporativo.

Como indicador cuantitativo, monitorizo semanalmente el porcentaje de horas ejecutadas respecto al acumulado previsto, usando el Gantt como línea base. Si la desviación supera el 15\,\%, se convoca reunión extraordinaria para decidir replanificación o activación de planes alternativos. Este umbral está pensado para detectar desviaciones antes de que comprometan hitos posteriores del camino crítico.

Una tarea se cierra formalmente cuando su evidencia está en el repositorio y cumple el criterio de aceptación del WBS, con trazabilidad completa a través de los commits en Git y registro de decisiones en \texttt{CHANGELOG.md}. Además, tras cada hito H1--H8 se realiza una verificación formal de cumplimiento cuyo resultado queda registrado en el repositorio junto con la evidencia del entregable correspondiente. La combinación de la hoja de seguimiento, las reuniones quincenales, la documentación continua y la verificación por hitos garantiza que ninguna desviación significativa pase inadvertida durante más de dos semanas~\cite{gep23}.

\subsection{Análisis crítico y coherencia de la planificación}

Toda planificación es, por definición, una hipótesis sobre el futuro, y conviene ser honesto con sus limitaciones antes de defenderla. En mi caso, he intentado construir un plan que sea a la vez ambicioso en alcance y prudente en márgenes, pero hay factores que escapan a mi control y que podrían alterar la ejecución prevista.

La principal fuente de incertidumbre es la dependencia de las colas GPU de MareNostrum~5, un recurso compartido cuya disponibilidad efectiva no puedo garantizar. A esto se suma que la cadena F4--F5--F6 es estrictamente secuencial y acumula 172 horas, de modo que cualquier retraso en una fase se propaga íntegramente a las siguientes. Si el desfase acumulado superase las dos semanas, la decisión más probable sería prescindir de Política~C y presentar únicamente A y~B, lo que reduciría el alcance pero preservaría la contribución central del TFG. Las estimaciones de las tareas experimentales (especialmente T5.5 y T6.5) llevan además un componente de incertidumbre mayor que las de infraestructura o documentación, porque dependen de la convergencia del entrenamiento y de la curva de aprendizaje con herramientas que aún no he usado de forma intensiva. Precisamente por eso he incorporado márgenes internos superiores al 20\,\% en esas tareas y he reservado tres semanas de holgura global (60\,h), dimensionadas para absorber el peor caso acumulado de los cuatro riesgos (+44\,h)~\cite{pmbok}.

Con todo, considero la planificación realista y coherente. Las 540\,h encajan con la carga de 15 ECTS más GEP, la dedicación semanal está calibrada con mi disponibilidad real, y la integración entre WBS, Gantt, registro de riesgos y mecanismos de seguimiento cubre las cuatro dimensiones fundamentales del PMBOK (alcance, tiempo, recursos y riesgos) dentro de un sistema de control que se revisa y actualiza semanalmente~\cite{gep23,gpi}.

%%%%%%%%%%%%%%%%%%%%%%%%%%%%%%%%%%%%%%%%%%%%%%%%%%%%%%%%%%%%%%%%%
\section{Presupuesto y sostenibilidad del proyecto}

\subsection{Supuestos económicos y criterios de cálculo}

El presupuesto estima el coste real que tendría el proyecto si fuera ejecutado por un equipo profesional en un contexto de empresa o consultoría tecnológica. Se trata de un presupuesto de coste de oportunidad, no de gasto personal del autor. Los supuestos son los siguientes.

\begin{itemize}[nosep]
\item \textbf{Referencia salarial base.} La Guía del Mercado Laboral 2026 de Hays~\cite{hays2026} de perfiles TIC en Barcelona proporciona los salarios de referencia. Los salarios corresponden a perfiles senior con tres a cinco años de experiencia, coherentes con la especialización técnica requerida por el proyecto.
\item \textbf{Tasa horaria cargada.} Se aplica la fórmula estándar de consultoría tecnológica en España de la siguiente forma:
\[
\text{Tasa cargada} = \frac{\text{Bruto anual}}{1.720}\;\times\;1{,}3\;(\text{SS})\;\times\;2{,}0\;(\text{overhead + margen})
\]
El factor 2,0 recoge gastos de estructura empresarial, beneficios sociales, equipamiento, administración y margen operativo, tal como aplican las empresas de servicios tecnológicos en proyectos de investigación y desarrollo.
\item \textbf{Duración del proyecto.} El proyecto dura 19 semanas, aproximadamente cinco meses efectivos, conforme al WBS de la planificación del proyecto~\cite{gep23e2}.
\item \textbf{Horas totales.} Se presupuestan 540 horas distribuidas en cuatro perfiles: Cap de projecte con 140 horas, Data scientist con 120 horas, ML engineer con 172 horas, y DevOps con 108 horas, sin modificar la distribución por tareas.
\item \textbf{Infraestructura HPC.} Se presupuesta a coste de mercado equivalente al precio on-demand de AWS, dado que el análisis de sostenibilidad económica debe reflejar el coste real del recurso si no hubiera acceso institucional. En la sección de sostenibilidad se distingue entre el coste presupuestado y el coste real para el autor.
\item \textbf{Costes indirectos.} Se estima el 20 por ciento sobre los costes de personal, como estimación estándar de overhead de estructura para proyectos informáticos.
\item \textbf{Contingencia.} Se aplica el 15 por ciento sobre el subtotal directo, dentro del rango recomendado del 10 a 20 por ciento para proyectos informáticos~\cite{gep23}.
\item \textbf{Imprevistos.} Se cuantifica como valor esperado de los riesgos R1 a R4 definidos en la planificación~\cite{gep23e2}, calculado como el producto del impacto económico por la probabilidad.
\end{itemize}

\subsection{Costes de personal (RRHH)}

La Tabla~\ref{tab:costes-rol} muestra los cuatro perfiles del equipo con sus tasas horarias cargadas.

\begin{table}[H]
\centering
\caption{Perfiles del equipo y tasa horaria cargada. Fuente: elaboración propia a partir de~\cite{hays2026}.}
\label{tab:costes-rol}
\footnotesize
\begin{tabular}{l l r r r}
\toprule
\textbf{Perfil (rol)} & \textbf{Equiv.\ Hays 2026 (senior)} & \textbf{Bruto anual (\euro{})} & \textbf{Bruto/h\,+\,SS (\euro{})} & \textbf{Tasa cargada (\euro{}/h)} \\
\midrule
Cap de projecte (CP)   & Project Manager Senior  & 60.000 & 45,35 & 87,00 \\
Data scientist (DS)    & Data Scientist Senior   & 55.000 & 41,57 & 82,00 \\
ML engineer (ML)       & ML / AI Engineer Senior & 65.000 & 49,13 & 98,00 \\
DevOps-HPC (DV)        & DevOps Engineer Senior  & 55.000 & 41,57 & 82,00 \\
\bottomrule
\end{tabular}
\end{table}

\noindent\textit{\scriptsize Tasa cargada = (Bruto anual / 1.720) $\times$ 1,3 (SS empleador) $\times$ 2,0 (factor overhead y margen). Referencia salarial: Hays España 2026~\cite{hays2026}, sección TIC, Barcelona, perfil senior (3--5 años de experiencia). El factor 2,0 es coherente con el modelo de facturación de proyectos de I+D en empresas de servicios tecnológicos.}

\noindent Los perfiles, la distribución de horas por tarea y los totales por rol coinciden con la planificación del proyecto~\cite{gep23e2}.

\subsection{Costes de personal por actividad (CPA)}

Cada tarea del WBS tiene un coste económico asociado, calculado como el producto de las horas estimadas por la tasa horaria cargada del perfil responsable:
\[
C_{\text{tarea}_i} = h_{\text{tarea}_i} \times \text{Tasa}_{\text{perfil}}
\]
Esta vinculación directa garantiza la trazabilidad entre la planificación temporal (diagrama de Gantt), la asignación de recursos humanos y la estimación económica del proyecto.

La Tabla~\ref{tab:cpa} desglosa el coste de personal por cada tarea del WBS, con subtotales por fase, usando las tasas cargadas de la Tabla~\ref{tab:costes-rol}. En las tareas con doble rol (marcadas con~$\dagger$), el coste se imputa al rol predominante. La tarea TG.2 (documentación continua) distribuye sus 30\,h entre los cuatro perfiles tal como se definió en la planificación~\cite{gep23e2}.

{\scriptsize
\setlength{\tabcolsep}{3pt}
\renewcommand{\arraystretch}{0.88}
\begin{tabularx}{\textwidth}{>{\ttfamily}l >{\raggedright\arraybackslash}X r l r r}
\caption{Costes por actividad del WBS (CPA). Fuente: elaboración propia.}\label{tab:cpa}\\
\toprule
\textbf{Cód.} & \textbf{Tarea} & \textbf{h} & \textbf{Rol} & \textbf{\euro{}/h} & \textbf{Coste (\euro{})} \\
\midrule
\endfirsthead
\multicolumn{6}{l}{\textit{(continuación de la Tabla~\ref{tab:cpa})}} \\
\toprule
\textbf{Cód.} & \textbf{Tarea} & \textbf{h} & \textbf{Rol} & \textbf{\euro{}/h} & \textbf{Coste (\euro{})} \\
\midrule
\endhead
\midrule
\multicolumn{6}{r}{\textit{Continúa en la página siguiente}} \\
\endfoot
\bottomrule
\endlastfoot
T1.1 & Acceso y cuenta BSC & 4 & DV & 82,00 & 328,00 \\
T1.2 & Entorno y scripts SLURM base & 10 & DV & 82,00 & 820,00 \\
T1.3 & Instalar vLLM y dependencias & 12 & DV & 82,00 & 984,00 \\
T1.4 & Preparar teacher y cachés$\,^\dagger$ & 10 & ML & 98,00 & 980,00 \\
T1.5 & Desplegar API teacher y healthchecks & 12 & DV & 82,00 & 984,00 \\
\multicolumn{2}{r}{\textit{Subtotal F1}} & \textit{48} & & & \textit{4.096,00} \\
\addlinespace[1.5pt]
T2.1 & Definir esquema req-resp y logs$\,^\dagger$ & 6 & DS & 82,00 & 492,00 \\
T2.2 & Impl.\ cliente benchmark y conc. & 12 & ML & 98,00 & 1.176,00 \\
T2.3 & Muestreador de prompts y norm. & 6 & ML & 98,00 & 588,00 \\
T2.4 & Validación 10\,000 peticiones & 8 & ML & 98,00 & 784,00 \\
\multicolumn{2}{r}{\textit{Subtotal F2}} & \textit{32} & & & \textit{3.040,00} \\
\addlinespace[1.5pt]
T3.1 & Instrumentación latencia y tokens$\,^\dagger$ & 8 & ML & 98,00 & 784,00 \\
T3.2 & Experimentos de carga (sweep) & 14 & ML & 98,00 & 1.372,00 \\
T3.3 & Registro memoria y uso GPU & 6 & DV & 82,00 & 492,00 \\
T3.4 & Agregación p50/p95 y tablas & 8 & DS & 82,00 & 656,00 \\
\multicolumn{2}{r}{\textit{Subtotal F3}} & \textit{36} & & & \textit{3.304,00} \\
\addlinespace[1.5pt]
T4.1 & Preparar eval.\ GSM8K y MATH & 8 & DS & 82,00 & 656,00 \\
T4.2 & Ejecutar evaluación teacher & 14 & ML & 98,00 & 1.372,00 \\
T4.3 & Análisis errores y baseline calidad & 12 & DS & 82,00 & 984,00 \\
\multicolumn{2}{r}{\textit{Subtotal F4}} & \textit{34} & & & \textit{3.012,00} \\
\addlinespace[1.5pt]
T5.1 & Selección student candidato & 6 & DS & 82,00 & 492,00 \\
T5.2 & Baseline student (mismo prot.) & 12 & ML & 98,00 & 1.176,00 \\
T5.3 & Decisión viabilidad destilación$\,^\dagger$ & 6 & DS & 82,00 & 492,00 \\
T5.4 & Preparar pipeline destilación & 12 & ML & 98,00 & 1.176,00 \\
T5.5 & Ejecutar destilación o alternativa & 22 & ML & 98,00 & 2.156,00 \\
T5.6 & Evaluar student destilado$\,^\dagger$ & 12 & DS & 82,00 & 984,00 \\
\multicolumn{2}{r}{\textit{Subtotal F5}} & \textit{70} & & & \textit{6.476,00} \\
\addlinespace[1.5pt]
T6.1 & Definir criterio quality gate & 6 & DS & 82,00 & 492,00 \\
T6.2 & Impl.\ Política B (cascada forc.) & 14 & ML & 98,00 & 1.372,00 \\
T6.3 & Evaluar Política B vs A$\,^\dagger$ & 12 & DS & 82,00 & 984,00 \\
T6.4 & Preparar features y dataset pred.$\,^\dagger$ & 10 & DS & 82,00 & 820,00 \\
T6.5 & Entrenar predictor (Política C) & 16 & ML & 98,00 & 1.568,00 \\
T6.6 & Evaluar Política C vs A y B$\,^\dagger$ & 10 & DS & 82,00 & 820,00 \\
\multicolumn{2}{r}{\textit{Subtotal F6}} & \textit{68} & & & \textit{6.056,00} \\
\addlinespace[1.5pt]
T7.1 & Esqueleto memoria y reprod. & 8 & CP & 87,00 & 696,00 \\
T7.2 & Redacción metodología y entorno & 12 & DV & 82,00 & 984,00 \\
T7.3 & Redacción resultados teacher/stud. & 16 & DS & 82,00 & 1.312,00 \\
T7.4 & Redacción routing y discusión & 16 & ML & 98,00 & 1.568,00 \\
T7.5 & Presentación y guion defensa & 10 & CP & 87,00 & 870,00 \\
T7.6 & Revisión final y entrega & 10 & CP & 87,00 & 870,00 \\
\multicolumn{2}{r}{\textit{Subtotal F7}} & \textit{72} & & & \textit{6.300,00} \\
\addlinespace[1.5pt]
TG.1 & Reuniones de seguimiento & 12 & CP & 87,00 & 1.044,00 \\
TG.2 & Doc.\ continua (coord.\ y registro) & 10 & CP & 87,00 & 870,00 \\
TG.2 & Doc.\ continua (experimentos) & 8 & DS & 82,00 & 656,00 \\
TG.2 & Doc.\ continua (código y decisiones) & 8 & ML & 98,00 & 784,00 \\
TG.2 & Doc.\ continua (infra.\ y entorno) & 4 & DV & 82,00 & 328,00 \\
TG.3 & Gestión repositorio y versionado & 48 & DV & 82,00 & 3.936,00 \\
\multicolumn{2}{r}{\textit{Subtotal TG}} & \textit{90} & & & \textit{7.618,00} \\
\addlinespace[1.5pt]
GP & Gestión del proyecto & 90 & CP & 87,00 & 7.830,00 \\
\midrule
\multicolumn{2}{r}{\textbf{TOTAL CPA}} & \textbf{540} & & & \textbf{47.732,00} \\
\end{tabularx}
}

\noindent\textit{\scriptsize $\dagger$\,Tarea con doble rol en la planificación; coste imputado al rol predominante. Verificación por rol: CP\,140\,h\,$\times$\,87,00\,=\,12.180,00; DS\,120\,h\,$\times$\,82,00\,=\,9.840,00; ML\,172\,h\,$\times$\,98,00\,=\,16.856,00; DV\,108\,h\,$\times$\,82,00\,=\,8.856,00. Suma\,=\,47.732,00\,\euro{}.}

\subsection{Costes de infraestructura}

Aunque el acceso a MareNostrum~5 está concedido por el BSC sin coste directo, el presupuesto refleja el coste equivalente de mercado para garantizar la validez del análisis económico.

Las horas de GPU se derivan del diseño experimental. Se contemplan entre tres y cinco configuraciones de modelo evaluadas sobre GSM8K y MATH, con dos a tres repeticiones por configuración. El proyecto se centra en inferencia y destilación supervisada, no en preentrenamiento desde cero, por lo que la demanda de GPU es considerablemente inferior a la de proyectos de entrenamiento completo. Se estiman 75 GPU-h planificadas más un 33 por ciento de reserva, lo que da un total de 100 GPU-h presupuestadas.

\begin{table}[H]
\centering
\caption{Costes de infraestructura. Fuente: elaboración propia a partir de~\cite{lambda2026,aws2026}.}
\label{tab:infraestructura}
\footnotesize
\begin{tabular}{l l r r r}
\toprule
\textbf{Recurso} & \textbf{Justificación} & \textbf{Cantidad} & \textbf{\euro{}/unidad} & \textbf{Coste (\euro{})} \\
\midrule
GPU H100 on-demand (AWS \texttt{p5.xlarge})  & 100\,GPU-h\,$\times$\,10,98\,\euro{}/h~\cite{aws2026} & 100\,h & 10,98 & 1.098,00 \\
Almacenamiento NFS (500\,GB, 5 meses)        & datasets, checkpoints, logs    & 5\,meses & 80,00 & 400,00 \\
Transferencia de red y egress                & descarga de modelos, backups   & estimación & -- & 150,00 \\
Monitorización y logging (cloud)             & métricas, alertas, dashboards  & 5\,meses & 40,00 & 200,00 \\
Cuota BSC imputada (simbólica)               & fracción cuota grupo inv.      & 5\,meses & 30,00 & 150,00 \\
\midrule
\textbf{Total infraestructura} & & & & \textbf{1.998,00} \\
\bottomrule
\end{tabular}
\end{table}

\noindent\textit{\scriptsize La GPU H100 se presupuesta a precio on-demand de AWS (instancia \texttt{p5.xlarge}, us-east-1)~\cite{aws2026}. Se estiman 100\,GPU-h: 75\,h planificadas más un 33\,\% de reserva para fallos, reruns y tuning. El almacenamiento cubre: dataset GSM8K/MATH ($\approx$\,100\,MB), checkpoints de modelos 7B--14B ($\approx$\,30\,GB), logs experimentales y versiones del código. La cuota BSC es una imputación simbólica que refleja el uso proporcional del acceso institucional.}

\subsection{Costes de software y herramientas}

Todo el software utilizado es de naturaleza intrínsecamente gratuita. El stack técnico se distribuye bajo licencias de código abierto como Apache 2.0 y BSD, y las herramientas complementarias son igualmente libres u open source. El coste de software se declara como cero euros, tal como recoge la Tabla~\ref{tab:software}.

\begin{table}[H]
\centering
\caption{Costes de software y herramientas. Fuente: elaboración propia.}
\label{tab:software}
\footnotesize
\begin{tabular}{l r r}
\toprule
\textbf{Recurso} & \textbf{Coste/mes (\euro{})} & \textbf{Coste final (\euro{})} \\
\midrule
vLLM, PyTorch, Python, NumPy, Transformers (HuggingFace) & 0 & 0,00 \\
Git, GitHub (plan Free / Student Developer Pack)         & 0 & 0,00 \\
VS~Code, Linux, Slurm (software del BSC)                 & 0 & 0,00 \\
Overleaf (plan Free, usuario único)                      & 0 & 0,00 \\
\midrule
\textbf{Total software} & & \textbf{0,00} \\
\bottomrule
\end{tabular}
\end{table}

\noindent\textit{\scriptsize Todo el software del proyecto tiene coste cero. El stack técnico (vLLM, PyTorch, Python, NumPy, Transformers) es open source bajo licencias Apache 2.0 y BSD. Git es open source y GitHub ofrece planes gratuitos. VS~Code es software libre y las herramientas de sistema del BSC (Linux, Slurm) son igualmente open source. Overleaf ofrece un plan gratuito suficiente para este proyecto. Ningún componente requiere pago por licencias comerciales.}

\subsection{Costes indirectos y estructurales}

Los costes indirectos incluyen los gastos del entorno de trabajo y un overhead de estructura empresarial del 20 por ciento sobre los costes de personal. Este overhead recoge administración, coordinación, infraestructura de oficina y soporte organizativo, y es un valor conservador habitual en proyectos tecnológicos de pequeña escala~\cite{pmbok}.

\begin{table}[H]
\centering
\caption{Costes estructurales directos. Fuente: elaboración propia.}
\label{tab:estructurales}
\footnotesize
\begin{tabular}{l l r}
\toprule
\textbf{Concepto} & \textbf{Cálculo} & \textbf{Subtotal (\euro{})} \\
\midrule
Espacio de trabajo (home office)    & 5 meses $\times$ 400\,\euro{}/mes               & 2.000,00 \\
Internet (fibra doméstica)          & 5 meses $\times$ 15\,\euro{}/mes                & 75,00  \\
Teléfono / comunicaciones           & 5 meses $\times$ 10\,\euro{}/mes                & 50,00  \\
Material fungible                   & estimación global                               & 50,00  \\
Impresión y encuadernación TFG      & 2 copias $\times$ 15\,\euro{}                   & 30,00  \\
Desplazamientos                     & trabajo íntegramente remoto (BSC vía SSH)       & 0,00   \\
\midrule
\textbf{Subtotal estructurales directos} & & \textbf{2.205,00} \\
\addlinespace[4pt]
Overhead de estructura empresarial  & 20\,\% $\times$ CPA (47.732,00\,\euro{})        & 9.546,40 \\
\midrule
\textbf{Total costes indirectos y estructurales} & & \textbf{11.751,40} \\
\bottomrule
\end{tabular}
\end{table}

\noindent\textit{\scriptsize Overhead del 20 por ciento: estimación conservadora para I+D con equipo reducido. En proyectos de mayor escala puede alcanzar entre el 30 y el 40 por ciento.}

\subsection{Amortización de hardware}

Los equipos de desarrollo se amortizan linealmente, imputando únicamente la fracción proporcional a los cinco meses del proyecto según la siguiente fórmula:

\[
\text{Coste amortizado} = \frac{\text{Precio de adquisición}}{\text{Vida útil en meses}} \times \text{Meses imputados}
\]

\begin{table}[H]
\centering
\caption{Amortización de hardware. Fuente: elaboración propia.}
\label{tab:hardware}
\footnotesize
\begin{tabular}{l r r r r}
\toprule
\textbf{Dispositivo} & \textbf{Precio (\euro{})} & \textbf{Vida útil (meses)} & \textbf{Meses imput.} & \textbf{Coste amort.\ (\euro{})} \\
\midrule
Portátil de desarrollo & 1.200,00 & 48 & 5 & 125,00 \\
Monitor externo (24'') & 200,00   & 60 & 5 & 16,67  \\
\midrule
\textbf{Total hardware} & & & & \textbf{141,67} \\
\bottomrule
\end{tabular}
\end{table}

\noindent\textit{\scriptsize Vida útil del portátil: cuatro años, equivalentes a 48 meses, estándar fiscal para equipos informáticos en España. Monitor: cinco años, equivalentes a 60 meses. Los precios corresponden a precio de venta al público de referencia en el momento de la adquisición.}

\subsection{Imprevistos asociados a riesgos}

Los riesgos R1 a R4 de la planificación~\cite{gep23e2} se cuantifican como el producto de su impacto económico por su probabilidad de ocurrencia. El impacto se valora a la tasa del ML engineer de 98 euros por hora, por ser el perfil más afectado, y la probabilidad Media se fija en 0,4. La Tabla~\ref{tab:imprevistos} recoge el coste esperado resultante.

\begin{table}[H]
\centering
\caption{Imprevistos asociados a riesgos. Fuente: elaboración propia a partir de la planificación del proyecto~\cite{gep23e2}.}
\label{tab:imprevistos}
\footnotesize
\begin{tabular}{l l r r r r}
\toprule
\textbf{Riesgo} & \textbf{Rol} & \textbf{+h} & \textbf{Impacto (\euro{})} & \textbf{Prob.} & \textbf{Esperado (\euro{})} \\
\midrule
R1. Colas GPU saturadas     & ML & 20 & 1.960,00 & 0,4 & 784,00  \\
R2. Destilación inviable    & ML & 6  & 588,00   & 0,4 & 235,20  \\
R3. Predictor no generaliza & ML & 10 & 980,00   & 0,4 & 392,00  \\
R4. Latencia routing        & ML & 8  & 784,00   & 0,4 & 313,60  \\
\midrule
\textbf{Total} & & \textbf{44} & \textbf{4.312,00} & & \textbf{1.724,80} \\
\bottomrule
\end{tabular}
\end{table}

\noindent\textit{\scriptsize Impacto igual a horas extra multiplicadas por 98 euros por hora, que es la tasa cargada para ML engineer senior. Coste esperado igual a impacto multiplicado por probabilidad.}

\subsection{Coste total estimado del proyecto}

El presupuesto incluye una reserva global de contingencia del 15 por ciento sobre el subtotal directo, práctica habitual en proyectos de ingeniería informática para absorber desviaciones imprevistas no cubiertas por los riesgos identificados~\cite{pmbok}. Su cálculo es:
\[
\text{Contingencia} = 0{,}15 \times \text{Subtotal directo} = 0{,}15 \times 61.623{,}07 = 9.243{,}46\;\text{\euro{}}
\]
Este porcentaje se sitúa dentro del rango del 10 al 20 por ciento recomendado para proyectos de I+D de pequeña escala.

\begin{table}[H]
\centering
\caption{Coste total estimado del proyecto. Fuente: elaboración propia.}
\label{tab:total}
\footnotesize
\begin{tabular}{l r}
\toprule
\textbf{Partida} & \textbf{Importe (\euro{})} \\
\midrule
Costes de personal por actividad (CPA)  & 47.732,00 \\
Costes de infraestructura               & 1.998,00  \\
Costes de software y herramientas       & 0,00      \\
Amortización de hardware                & 141,67    \\
Costes indirectos y estructurales       & 11.751,40 \\
\cmidrule{2-2}
Subtotal directo                        & 61.623,07 \\
Contingencia (15\,\%)                   & 9.243,46  \\
Imprevistos (coste esperado)            & 1.724,80  \\
\midrule
\textbf{TOTAL PROYECTO}                 & \textbf{72.591,33} \\
\bottomrule
\end{tabular}
\end{table}

\begin{figure}[H]
\centering
\begin{tikzpicture}
\pie[text=legend, radius=2.15, color={blue!35,orange!45,gray!35}]{65.8/RRHH,2.8/Infraestructura,31.4/Costes operativos e indirectos}
\end{tikzpicture}
\caption{Distribución porcentual del presupuesto total por partidas principales. Fuente: elaboración propia.}
\label{fig:distribucion-presupuesto}
\end{figure}

\noindent El presupuesto es coherente con las 540 horas de la planificación~\cite{gep23e2}. La partida de personal representa el 66 por ciento del total, los costes de software son cero al emplear exclusivamente software libre, y la infraestructura supone el 3 por ciento, reflejando el coste de mercado de GPUs H100 en cloud. El coste total de 72.591 euros se sitúa dentro del rango razonable para proyectos de investigación aplicada en inteligencia artificial de pequeña escala.

\subsection{Control económico y seguimiento del presupuesto}

El seguimiento se realiza de forma semanal y al cierre de cada fase, registrando horas reales por tarea y rol, gastos de infraestructura y horas de GPU por experimento mediante JobID SLURM.

Los indicadores cuantitativos de control son los siguientes. La desviación absoluta y porcentual por tarea permiten detectar sobrecostes puntuales:
\[
\Delta C_{\text{tarea}} = C_{\text{real}} - C_{\text{plan}} \qquad
\Delta C\% = \frac{C_{\text{real}} - C_{\text{plan}}}{C_{\text{plan}}} \times 100
\]
La desviación de horas por rol y de horas de GPU permiten controlar el consumo de recursos:
\[
\Delta h_{\text{rol}} = h_{\text{real,rol}} - h_{\text{plan,rol}} \qquad
\Delta \text{GPU} = \text{GPU-h}_{\text{real}} - \text{GPU-h}_{\text{plan}}
\]
Adicionalmente se calcula el índice de rendimiento de costes (CPI) mediante la técnica del valor ganado (Earned Value Management), que proporciona una medida normalizada de la eficiencia económica del proyecto:
\[
\text{CPI} = \frac{\text{EV}}{\text{AC}}
\]
donde EV (Earned Value) es el valor del trabajo completado según el presupuesto planificado y AC (Actual Cost) es el coste real incurrido. Un CPI superior a 1 indica que el proyecto se ejecuta por debajo del presupuesto, mientras que un CPI inferior a 1 señala un sobrecoste que requiere acciones correctivas como la repriorización de tareas, la reasignación de recursos o el ajuste de alcance. El coste acumulado del proyecto se compara periódicamente con la curva de gasto planificada al cierre de cada fase, permitiendo detectar desviaciones de forma temprana.

Los umbrales y acciones correctivas se establecen de la siguiente forma:
\begin{itemize}[nosep]
\item Si la desviación de coste supera el 15 por ciento en una fase del camino crítico, de la F1 a la F6, se convoca una reunión extraordinaria con posible replanificación o recorte de alcance priorizando las Políticas A y B.
\item Si la desviación de coste supera el 25 por ciento en cualquier fase, se congelan las configuraciones experimentales y se reducen las repeticiones al mínimo, de tres a dos.
\item Si la desviación en GPU supera 50 horas antes de la semana 10, se activa el plan de contingencia R1 con jobs nocturnos, reducción de repeticiones y limitación de sweeps.
\item Si la desviación de horas por rol supera el 20 por ciento al cierre de fase, se rebalancea el esfuerzo en las fases siguientes.
\end{itemize}

\noindent La hoja de seguimiento semanal, las reuniones quincenales TG.1 y la verificación por hitos H1 a H8 garantizan que ninguna desviación significativa pase inadvertida más de dos semanas.

\subsection{Marco general de la sostenibilidad}

A lo largo de mis estudios en el Grado de Ingeniería Informática en la FIB he podido profundizar en los aspectos clave de la sostenibilidad a través de diversas asignaturas, lo que me ha llevado a reconocer su importancia al desarrollar este TFG. Tras completar la encuesta del proyecto de investigación EDINSOST2-ODS sobre desarrollo sostenible, se ha realizado una autoevaluación del impacto del proyecto en las tres dimensiones de sostenibilidad.

La sostenibilidad, entendida desde una perspectiva integral que abarca las dimensiones económica, social y ambiental, es un factor relevante en cualquier proyecto tecnológico. En el contexto de proyectos que hacen uso intensivo de infraestructura de computación de altas prestaciones y modelos de lenguaje de gran tamaño, estas dimensiones están estrechamente interrelacionadas dado que la eficiencia computacional reduce simultáneamente costes, consumo energético y barreras de acceso. Un análisis adecuado permite no solo reducir el impacto inmediato sino también aportar soluciones sostenibles aplicables a proyectos futuros.

Me comprometo a registrar y publicar el consumo energético real de cada experimento comparándolo con las estimaciones de este documento, y a dedicar una sección de la memoria del TFG al balance neto de sostenibilidad del sistema de cascada.

\subsection{Dimensión económica}

Después de revisar el presupuesto del proyecto, considero que es apropiado teniendo en cuenta los beneficios que puede aportar. El coste total estimado asciende a 72.591 euros como coste de oportunidad, mientras que el gasto real para el autor es prácticamente cero dado que la infraestructura está cedida por el BSC y el software es intrínsecamente gratuito. Este proyecto tiene el potencial de ahorrar una cantidad significativa de recursos computacionales a las organizaciones que adopten el sistema de cascada, lo que se traduce en un ahorro de costes a largo plazo.

Desde el punto de vista de la viabilidad en producción, un sistema de cascada que derive el 60 por ciento de las peticiones al student puede reducir las GPU horas anuales aproximadamente un 30 por ciento. La Tabla~\ref{tab:ahorro-cascada} estima el ahorro anual proyectado según proveedor cloud~\cite{lambda2026,aws2026}.

\begin{table}[H]
\centering
\caption{Estimación de ahorro anual con sistema de cascada según proveedor cloud. Fuente: elaboración propia a partir de~\cite{lambda2026,aws2026}.}
\label{tab:ahorro-cascada}
\footnotesize
\begin{tabular}{l r r r r}
\toprule
\textbf{Proveedor} & \textbf{\euro{}/GPU-h} & \textbf{GPU-h/año (solo teacher)} & \textbf{Ahorro 30\,\%} & \textbf{Ahorro anual (\euro{})} \\
\midrule
Lambda Labs  & 2,49  & 8.760 & 2.628\,h & 6.544 \\
AWS (p5.xlarge) & 10,98 & 8.760 & 2.628\,h & 28.855 \\
\bottomrule
\end{tabular}
\end{table}

\noindent\textit{\scriptsize Se asume un despliegue continuo de una GPU H100 durante un año completo, equivalente a 8.760 GPU horas. El ahorro del 30 por ciento corresponde al escenario en que el sistema de cascada redirija el 60 por ciento de las peticiones al student con una tasa de fallback inferior al 40 por ciento.}

Esta viabilidad hace el proyecto relevante para organizaciones con recursos limitados. La vida útil de los resultados se estima entre dos y tres años, aunque los principios de diseño del routing y la destilación seguirán siendo aplicables más allá de ese horizonte.

\subsection{Dimensión ambiental}

El impacto ambiental se concentra en el consumo energético de las GPUs H100 durante las fases experimentales. La Tabla~\ref{tab:consumo-fase} desglosa la estimación por fase, asumiendo 375 watts por GPU y un factor de emisión de 0,20 kg CO\textsubscript{2}/kWh~\cite{ree2024}.

\begin{table}[H]
\centering
\caption{Estimación de consumo energético y emisiones por fase experimental. Fuente: elaboración propia a partir de~\cite{ree2024}.}
\label{tab:consumo-fase}
\footnotesize
\begin{tabular}{l r r r}
\toprule
\textbf{Fase experimental} & \textbf{GPU-h estimadas} & \textbf{Consumo (kWh)} & \textbf{CO\textsubscript{2}eq (kg)} \\
\midrule
Caracterización teacher (F2 a F4)         & 15 & 5,63  & 1,13 \\
Destilación del student (F5)              & 30 & 11,25 & 2,25 \\
Evaluación políticas routing (F6)         & 20 & 7,50  & 1,50 \\
Reruns exploratorios y ajustes (F3 a F5)  & 10 & 3,75  & 0,75 \\
\midrule
\textbf{Total planificado}                & \textbf{75} & \textbf{28,13} & \textbf{5,63} \\
\addlinespace[2pt]
Reserva del 33 por ciento                & 25 & 9,38  & 1,88 \\
\midrule
\textbf{Total presupuestado}              & \textbf{100} & \textbf{37,50} & \textbf{7,50} \\
\bottomrule
\end{tabular}
\end{table}

\noindent\textit{\scriptsize Consumo calculado como GPU-h multiplicadas por 0,375 kW. Emisiones calculadas como kWh multiplicados por 0,20 kgCO\textsubscript{2}/kWh, que es el factor de emisión medio del mix eléctrico español en 2024~\cite{ree2024}. La reserva del 33 por ciento cubre fallos de jobs, reruns y ajustes de configuración.}

El consumo directo planificado ronda los 28 kWh (5,6 kg CO\textsubscript{2}eq), y con la reserva asciende a 37,5 kWh (7,5 kg CO\textsubscript{2}eq). MareNostrum 5 ofrece ventajas ambientales frente a centros de datos comerciales gracias a su eficiencia energética~\cite{mn5}.

El argumento ambiental central del proyecto es que un modelo destilado de siete mil millones de parámetros requiere significativamente menos cómputo por token que el teacher de 14 mil millones, lo que reduce directamente el consumo energético en despliegues a gran escala. Entornos como MareNostrum 5 amplifican esta ventaja gracias a sus sistemas de refrigeración optimizados y economías de escala.

En el peor caso, si todas las consultas pasan por el student sin resolverse antes de redirigirse al teacher, el coste energético por petición sería un 50 por ciento superior al del teacher solo. La Política C mitiga este riesgo con un umbral de confianza que deriva directamente al teacher las peticiones difíciles, y si la tasa de fallback supera el 70 por ciento el sistema revierte automáticamente a servir solo con el teacher. El riesgo de efecto rebote se analizará explícitamente en la memoria del TFG.

A nivel global, el sector de las Tecnologías de la Información y Comunicación representa entre el 2,5 y el 3 por ciento de las emisiones globales de gases de efecto invernadero, con una proyección de aumento para las próximas décadas. En este contexto, proyectos como este que buscan reducir el coste computacional de la inferencia de modelos de lenguaje contribuyen de forma directa a mitigar el impacto ambiental creciente de la inteligencia artificial. Considero que la huella de carbono del proyecto es reducida y proporcionada a los beneficios esperados, y que los resultados obtenidos pueden contribuir a que futuros despliegues de modelos de lenguaje sean más eficientes energéticamente.

Durante la ejecución, el consumo energético de cada job se registrará mediante sacct con el campo ConsumedEnergy, almacenando los resultados en energy\_log.csv. Los valores estimados se contrastarán con las medidas reales y se reportarán en la memoria final.

\subsection{Dimensión social}

Realizar este proyecto en el BSC me ha permitido aprender a gestionar proyectos de investigación aplicada en un entorno de supercomputación real, trabajando con tecnologías punteras como la destilación del conocimiento, la inferencia eficiente de modelos de lenguaje y la orquestación de recursos HPC. Esta experiencia me ha proporcionado una comprensión profunda de cómo optimizar el despliegue de modelos de inteligencia artificial para hacerlo más accesible y eficiente.

El proyecto tiene un impacto social positivo al contribuir a democratizar el acceso a modelos lingüísticos avanzados. La inferencia con modelos de gran tamaño exige GPUs costosas que solo organizaciones con presupuestos elevados pueden permitirse en producción continua. Al reducir los recursos computacionales necesarios mediante destilación y routing inteligente, se amplía el abanico de organizaciones beneficiarias a centros de investigación públicos, universidades y startups con recursos limitados. De manera análoga a cómo la automatización de procesos reduce tiempos de respuesta y mejora la experiencia de los usuarios finales, un sistema de cascada eficiente permite ofrecer servicios de inteligencia artificial de alta calidad a un coste operativo significativamente menor, beneficiando directamente a las organizaciones que lo adopten. La publicación abierta del código, los scripts de evaluación y la documentación del pipeline refuerzan esta transferencia de conocimiento hacia la comunidad investigadora.

A nivel personal, el TFG me ha permitido adquirir competencias en computación de altas prestaciones, reproducibilidad experimental, gestión de proyectos con infraestructura compartida y optimización de modelos de lenguaje. Estas son habilidades con alta demanda en el mercado laboral actual del sector de machine learning e inteligencia artificial, y considero que la formación práctica obtenida durante este proyecto complementa de forma significativa los conocimientos teóricos adquiridos durante el Grado.

Entre los riesgos sociales se identifican el efecto rebote por reducción del coste por petición, la posible amplificación de sesgos del teacher en el student, y la dependencia de infraestructuras de computación de altas prestaciones concentradas en pocas instituciones. Para mitigarlos, el sistema se evaluará con benchmarks objetivos y reproducibles definidos en el alcance, y el pipeline se publicará con sus limitaciones explicitadas. El proyecto no trabaja con datos personales ni toma decisiones automatizadas sobre individuos.

\subsection{Autoevaluación de sostenibilidad}

Tras realizar el análisis de las tres dimensiones, valoro positivamente el ejercicio de reflexión que ha supuesto elaborar este informe de sostenibilidad. Me ha permitido tomar conciencia de que las decisiones técnicas adoptadas durante la planificación del proyecto, como priorizar la destilación del conocimiento frente al uso directo de modelos de gran tamaño o comprometerse a publicar el código de forma abierta, tendrán consecuencias directas sobre la eficiencia económica, el consumo energético y la accesibilidad social de los resultados. Identifico como fortaleza principal del proyecto su potencial para reducir el coste computacional y las emisiones asociadas al despliegue de modelos de lenguaje. Como área de mejora, reconozco la necesidad de monitorizar activamente el riesgo de efecto rebote y de amplificación de sesgos durante la futura fase de explotación. Esta autoevaluación ha reforzado mi compromiso con integrar criterios de sostenibilidad a lo largo de la ejecución de este TFG y en mis futuros proyectos de ingeniería.

%%%%%%%%%%%%%%%%%%%%%%%%%%%%%%%%%%%%%%%%%%%%%%%%%%%%%%%%%%%%%%%%%
\part{Marco teórico}
\label{part:marco-teorico}
%%%%%%%%%%%%%%%%%%%%%%%%%%%%%%%%%%%%%%%%%%%%%%%%%%%%%%%%%%%%%%%%%

Este bloque constituye la base conceptual del trabajo. Su objetivo es presentar, de forma progresiva y autocontenida, todos los fundamentos técnicos que un lector con formación en ingeniería informática necesita para comprender las decisiones de diseño, los experimentos y los resultados que se discuten en las partes posteriores de la memoria. La exposición parte de conceptos elementales (qué es una red neuronal, qué es un token, cómo aprende un modelo) y avanza gradualmente hacia las técnicas específicas sobre las que se articula el TFG, entre ellas la arquitectura transformer, el \textit{fine-tuning} y el PEFT, la destilación de conocimiento, el serving de modelos de lenguaje y la asignación dinámica de peticiones entre modelos de distinto tamaño. Cada concepto se motiva en relación con el problema práctico del proyecto antes de desarrollarse en profundidad, de manera que la teoría no quede aislada del trabajo experimental.

%%%%%%%%%%%%%%%%%%%%%%%%%%%%%%%%%%%%%%%%%%%%%%%%%%%%%%%%%%%%%%%%%
\section{Fundamentos de deep learning}
\label{sec:fundamentos}

Antes de introducir la arquitectura de los modelos de lenguaje conviene fijar una serie de conceptos previos que se utilizarán de forma continua en el resto de la memoria. El objetivo de esta sección no es desarrollar una introducción exhaustiva a todo el campo, sino presentar de forma clara las ideas mínimas necesarias para entender cómo se construyen, entrenan y evalúan los modelos que se estudian en este trabajo. En particular, se introducirán los conceptos de red neuronal, función de pérdida, optimización por gradientes, tokenización y modelo de lenguaje.

\subsection{Redes neuronales artificiales}

Las redes neuronales artificiales constituyen la base de la mayor parte de modelos modernos de \textit{deep learning}. En esencia, una red neuronal es una función parametrizada que transforma una entrada en una salida mediante una secuencia de operaciones matemáticas sencillas. Su interés no reside en la complejidad de cada operación aislada, sino en la capacidad que emerge al combinar muchas de ellas de forma estructurada y al ajustar automáticamente sus parámetros a partir de datos.

La unidad conceptual más simple dentro de una red es la neurona. Una neurona recibe varios valores de entrada, asigna a cada uno de ellos un peso, suma el resultado y añade un sesgo. A continuación, aplica una función de activación que introduce no linealidad en el modelo. De forma simplificada, el cálculo realizado por una neurona puede escribirse como
\[
z = \sum_{i=1}^{n} w_i x_i + b,
\qquad
a = \sigma(z),
\]
donde $x_1, \dots, x_n$ son las entradas, $w_1, \dots, w_n$ los pesos asociados, $b$ es el sesgo, $\sigma$ es la función de activación y $a$ es la salida producida por la neurona. Sin la activación, la red completa se reduciría a una combinación lineal de las entradas y tendría una capacidad expresiva muy limitada.

Cuando varias neuronas se agrupan y procesan conjuntamente la misma entrada, se obtiene una capa. En una capa densa o \textit{fully connected}, cada neurona de una capa recibe información de todas las salidas de la capa anterior. Si una capa recibe un vector de entrada $x \in \mathbb{R}^{n}$ y produce una salida $h \in \mathbb{R}^{m}$, su operación puede expresarse de forma compacta como
\[
h = \sigma(Wx + b),
\]
donde $W \in \mathbb{R}^{m \times n}$ contiene los pesos de la capa y $b \in \mathbb{R}^{m}$ contiene los sesgos. Esta formulación matricial es especialmente importante porque es la forma en la que estas operaciones se implementan de manera eficiente sobre GPU.

\begin{figure}[H]
\centering
\begin{tikzpicture}[
    neuron/.style={circle, draw=black, thick, minimum size=8mm, inner sep=0pt},
    input/.style={neuron, fill=bcdata!30},
    hidden/.style={neuron, fill=bcstudent!25},
    output/.style={neuron, fill=bcteacher!25},
    conn/.style={-Latex, draw=black!60, line width=0.7pt},
    labelbox/.style={rounded corners, draw=black!35, fill=white, align=center, inner sep=5pt}
]

% Entrada
\foreach \i/\y/\txt in {1/1.8/$x_1$,2/0.6/$x_2$,3/-0.6/$x_3$,4/-1.8/$x_n$}{
    \node[input] (I\i) at (0,\y) {};
    \node[left=2mm of I\i] {\small \txt};
}

% Capa oculta 1
\foreach \i/\y in {1/1.2,2/0,3/-1.2}{
    \node[hidden] (H1\i) at (3,\y) {};
}

% Capa oculta 2
\foreach \i/\y in {1/0.8,2/-0.8}{
    \node[hidden] (H2\i) at (6,\y) {};
}

% Salida
\node[output] (O1) at (9,0) {};
\node[right=2mm of O1] {\small $\hat{y}$};

% Conexiones entrada -> oculta 1
\foreach \i in {1,2,3,4}{
    \foreach \j in {1,2,3}{
        \draw[conn] (I\i) -- (H1\j);
    }
}

% Conexiones oculta 1 -> oculta 2
\foreach \i in {1,2,3}{
    \foreach \j in {1,2}{
        \draw[conn] (H1\i) -- (H2\j);
    }
}

% Conexiones oculta 2 -> salida
\foreach \i in {1,2}{
    \draw[conn] (H2\i) -- (O1);
}

% Títulos de capas
\node[above=4mm of I1] {\small \textbf{Capa de entrada}};
\node[above=4mm of H11] {\small \textbf{Capas ocultas}};
\node[above=4mm of O1] {\small \textbf{Salida}};

% Caja explicativa
\node[labelbox] at (4.5,-3.0) {\small En cada neurona\\[2pt] $z=\sum_i w_i x_i + b$\\[2pt] $a=\sigma(z)$};

% Caja lateral
\node[labelbox] at (4.5,3.0) {\small Los pesos $w$ y los sesgos $b$\\son parámetros aprendidos};

\end{tikzpicture}
\caption{Esquema simplificado de una red neuronal densa. Cada neurona combina las señales de la capa anterior mediante pesos, añade un sesgo y aplica una activación no lineal.}
\label{fig:red-neuronal-basica}
\end{figure}

Como se observa en la Figura \ref{fig:red-neuronal-basica}, una red neuronal puede entenderse como una composición de capas. La salida de una capa pasa a ser la entrada de la siguiente, lo que permite construir transformaciones progresivamente más abstractas. Las primeras capas suelen capturar patrones relativamente simples, mientras que las capas posteriores combinan esos patrones para representar relaciones más complejas. Esta idea resulta fundamental para entender por qué las arquitecturas profundas han tenido tanto éxito en tareas de visión, voz y lenguaje.

El término \textit{profunda} hace referencia precisamente a esa acumulación de capas. Una red con pocas capas puede modelar relaciones simples, pero suele necesitar muchas neuronas para aproximar funciones complejas. En cambio, una red profunda distribuye la representación en varios niveles, lo que permite reutilizar patrones intermedios y ganar capacidad expresiva de una forma más eficiente. En la práctica, esta profundidad es uno de los factores que han hecho posible el desarrollo de modelos de gran capacidad como los LLM.

Conviene remarcar que no todas las redes neuronales tienen exactamente la estructura mostrada en la figura. La representación anterior es deliberadamente sencilla y sirve solo como punto de partida conceptual. Los modelos modernos utilizados en procesamiento del lenguaje natural, y en particular los basados en la arquitectura Transformer, no se construyen como una simple pila de capas densas. Sin embargo, siguen respondiendo a la misma idea general. Existen entradas numéricas, parámetros aprendibles, transformaciones lineales, activaciones no lineales y una composición de bloques que produce una salida final.

Finalmente, los pesos y sesgos de todas las capas forman el conjunto de parámetros del modelo. Cuando se afirma que un modelo tiene 7B parámetros, lo que se está diciendo es que contiene aproximadamente siete mil millones de valores numéricos ajustables. Estos parámetros se almacenan en memoria y se modifican durante el entrenamiento para que la red aprenda a resolver la tarea deseada. Este detalle, que puede parecer puramente técnico, es en realidad central en el contexto de este trabajo, ya que el número de parámetros condiciona de forma directa el consumo de memoria, el coste computacional, la latencia de inferencia y, en general, la viabilidad de desplegar el modelo de forma eficiente.
\subsection{Entrenamiento, función de pérdida y descenso por gradiente}

Entrenar una red neuronal consiste en encontrar un conjunto de parámetros $\theta$ (la colección de todas las matrices $W$ y vectores $b$) que minimice una \textbf{función de pérdida} $\mathcal{L}(\theta)$ sobre un conjunto de datos de entrenamiento $\mathcal{D}$. La pérdida mide cuánto se desvían las predicciones del modelo de las respuestas deseadas; si la pérdida es cero, el modelo predice perfectamente sobre $\mathcal{D}$.

En el caso de los modelos de lenguaje, la pérdida habitual es la \textbf{entropía cruzada} (\textit{cross-entropy}) entre la distribución predicha por el modelo y la distribución objetivo (que en el caso supervisado más simple es una distribución degenerada concentrada en el token correcto). Para un único ejemplo $(x, y)$ donde $y$ es el token correcto y el modelo asigna probabilidad $p_\theta(y \mid x)$ a ese token, la pérdida es
\[
\mathcal{L}(x,y;\theta) = -\log p_\theta(y \mid x).
\]
El signo negativo convierte la log-probabilidad (que es negativa o cero) en una cantidad positiva que se quiere minimizar: cuanto mayor es la probabilidad que el modelo asigna al token correcto, menor es la pérdida.

La minimización se realiza mediante \textbf{descenso por gradiente estocástico} (\textit{stochastic gradient descent}, SGD) o alguna de sus variantes con momento adaptativo, siendo \textbf{AdamW} la más utilizada en el entrenamiento de LLMs modernos. La idea es conceptualmente simple. En cada paso se calcula el gradiente de la pérdida respecto a todos los parámetros, $\nabla_\theta \mathcal{L}$, y se actualizan los parámetros dando un pequeño paso en la dirección que reduce la pérdida según
\[
\theta \;\leftarrow\; \theta - \eta\,\nabla_\theta \mathcal{L},
\]
donde $\eta$ es la \textbf{tasa de aprendizaje} (\textit{learning rate}), un hiperparámetro que controla el tamaño del paso. Si $\eta$ es demasiado grande el proceso diverge; si es demasiado pequeño, el entrenamiento es excesivamente lento. En la práctica, $\eta$ suele variar durante el entrenamiento siguiendo un esquema de \textit{warmup} inicial seguido de un decaimiento gradual.

El cálculo del gradiente en redes profundas se realiza mediante el algoritmo de \textbf{retropropagación} (\textit{backpropagation}), que no es más que la aplicación sistemática de la regla de la cadena del cálculo diferencial a través de todas las capas de la red. El coste computacional de un paso de retropropagación es del mismo orden que el del paso ``hacia adelante'' (\textit{forward pass}), pero requiere almacenar en memoria las activaciones intermedias de cada capa para poder calcular las derivadas parciales; este requisito de memoria es una de las razones por las que entrenar modelos grandes exige GPUs con mucha VRAM.

Un aspecto que conviene retener para secciones posteriores es que todo el entrenamiento de redes neuronales, incluido el \textit{fine-tuning} y la destilación, se reduce en esencia a repetir este ciclo de forward, cálculo de pérdida, backward y actualización de parámetros. Lo que varía entre técnicas es qué función de pérdida se usa, qué parámetros se actualizan y qué datos se emplean. Esta uniformidad conceptual es lo que permite implementar la destilación de conocimiento como un caso particular del \textit{fine-tuning}, tal y como se verá en las secciones correspondientes.

\subsection{Tokens y modelos de lenguaje}

Un \textbf{modelo de lenguaje} es, en su formulación más general, un modelo probabilístico que asigna probabilidades a secuencias de texto. Dado un prefijo $(t_1, t_2, \dots, t_n)$, el modelo estima la distribución condicional $p(t_{n+1} \mid t_1, \dots, t_n)$ sobre la siguiente unidad de texto, denominada \textbf{token}.

La elección de qué constituye un token no es trivial y tiene consecuencias directas sobre el rendimiento del modelo. Los primeros modelos de lenguaje operaban a nivel de palabra completa, lo que producía vocabularios enormes e incapacidad para manejar palabras fuera de vocabulario (\textit{out-of-vocabulary}, OOV). Operar a nivel de carácter resuelve el problema de OOV pero produce secuencias muy largas que son costosas de procesar. La solución adoptada por los modelos actuales es la tokenización \textbf{subword}, siendo \textit{Byte Pair Encoding} (BPE) el algoritmo más extendido. BPE construye un vocabulario de tamaño fijo (típicamente entre 30\,000 y 150\,000 unidades) fusionando iterativamente los pares de bytes o caracteres más frecuentes del corpus de entrenamiento. El resultado es que las palabras comunes se representan como un único token, mientras que las raras se descomponen en fragmentos más pequeños que el modelo ya conoce. Esta representación ofrece un compromiso favorable entre el tamaño del vocabulario, la longitud de las secuencias y la capacidad de manejar texto arbitrario.

Una vez definida la unidad token, un modelo de lenguaje autoregresivo modela la probabilidad conjunta de una secuencia factorizándola como producto de condicionales mediante
\[
p(t_1, t_2, \dots, t_N) = \prod_{i=1}^{N} p(t_i \mid t_1, \dots, t_{i-1}).
\]
Generar texto con este modelo consiste en muestrear repetidamente de $p(t_{i+1} \mid t_1, \dots, t_i)$, añadir el token generado al contexto y repetir. Esta generación secuencial es lo que se conoce como \textbf{generación autoregresiva} y es el modo de operación fundamental de todos los LLMs actuales, incluidos los que se emplean en este TFG.

El entrenamiento de un modelo de lenguaje autoregresivo consiste, por tanto, en ajustar los parámetros $\theta$ de la red neuronal para que la distribución $p_\theta$ se aproxime lo máximo posible a la distribución real del lenguaje, minimizando la entropía cruzada token a token sobre un gran corpus de texto. Este proceso se denomina \textbf{preentrenamiento} y es la fase más costosa del ciclo de vida de un LLM: los modelos de la familia Qwen 2.5 que se usan en el TFG fueron preentrenados sobre billones de tokens durante semanas o meses en clústeres de miles de GPUs~\cite{qwen2}.

\subsection{Arquitectura transformer a nivel funcional}

Con los fundamentos anteriores en mente, se enlaza aquí con la arquitectura concreta que subyace a los modelos de lenguaje modernos, el \textit{transformer}, y, más adelante en esta misma sección, con las razones por las que servir estos modelos en producción resulta tan costoso. Se fija también el vocabulario técnico (prefill, decode, KV cache, TTFT, throughput) que se emplea de forma recurrente en el resto de la memoria.

Un modelo de lenguaje grande (LLM) moderno es, en esencia, una red neuronal basada en la arquitectura \textit{transformer}~\cite{vaswani2017attention}, entrenada para predecir el siguiente token a partir de los tokens anteriores. Esta parte no pretende sustituir un manual sobre transformers, pero sí cubrir con suficiente detalle los componentes que aparecen en el resto de la memoria. La razón es sencilla: casi todas las decisiones de ingeniería del TFG (el KV cache, el balance entre prefill y decode, el diseño del router, las optimizaciones de serving) se entienden mejor cuando se tiene en mente qué está haciendo realmente el modelo capa a capa.

\paragraph{Idea general de predecir el siguiente token.} Un LLM autoregresivo aprende, durante el preentrenamiento, a modelar la distribución condicional $p_\theta(t_{n+1} \mid t_1, \dots, t_n)$. Dado un prefijo de tokens, asigna una probabilidad a cada posible continuación. Una vez entrenado, generar texto consiste en muestrear (o elegir por \textit{argmax}) el siguiente token según esa distribución, añadirlo al contexto y repetir. Por construcción, el modelo no tiene un estado interno explícito como los sistemas recurrentes clásicos: todo lo que ``recuerda'' sobre el contexto está codificado en los vectores de claves y valores que acompañan a cada token del prefijo. De esa forma, el cálculo se paraleliza muy bien durante el entrenamiento (todos los tokens se procesan a la vez) y es lo que permite escalar a modelos de miles de millones de parámetros sobre clústeres de GPUs.

\paragraph{Por qué reemplazó a las RNN.} Antes del transformer, el estado del arte en procesamiento de lenguaje eran modelos recurrentes (LSTM, GRU) que procesaban un token tras otro y mantenían un estado oculto que se iba actualizando. Esto tenía dos limitaciones serias. La primera es que el entrenamiento era inherentemente secuencial y muy difícil de paralelizar a gran escala. La segunda es que, en dependencias largas, la señal del gradiente se desvanecía antes de llegar a los tokens del principio, lo que dificultaba aprender relaciones a larga distancia. El transformer ataca las dos a la vez: el mecanismo de atención da a cada token acceso directo a cualquier otro token del contexto (no hay que ``propagar'' el estado paso a paso), y el grafo de cómputo se puede desenrollar por completo durante el entrenamiento, con la consiguiente mejora brutal en throughput de entrenamiento sobre GPUs y TPUs.

\paragraph{Tokenización y embeddings.} El primer paso al alimentar un LLM es convertir el texto crudo en una secuencia de tokens enteros. La mayoría de modelos actuales, incluidos los de la familia Qwen que se usan en este TFG, emplean una tokenización \textit{subword} de tipo BPE (\textit{Byte Pair Encoding}) o una variante suya. La intuición es que ni usar el carácter ni la palabra completa son buenas unidades: el carácter produce secuencias muy largas y la palabra sufre con términos raros o compuestos. Las unidades subword ofrecen un compromiso: las palabras frecuentes acaban siendo un único token y las raras se descomponen en trozos que el modelo ya conoce.

Una vez tokenizado, cada token $t_i$ entero se transforma en un vector denso $e_i \in \mathbb{R}^{d}$ mediante una capa de \textbf{embeddings}. Formalmente, existe una matriz $E \in \mathbb{R}^{|V|\times d}$ (con $|V|$ el tamaño del vocabulario) y $e_i = E_{t_i,:}$. Este espacio de embeddings es aprendido: durante el preentrenamiento, palabras con significados o roles gramaticales parecidos acaban cerca en él. La dimensión $d$ (a menudo llamada \textit{hidden size}) es un parámetro arquitectónico importante; en los modelos del TFG va desde 1\,024 hasta 5\,120 aproximadamente, según el tamaño. Esa misma matriz $E$ se suele reutilizar (\textit{tied embeddings}) en la cabecera de salida del modelo para proyectar de nuevo a logits sobre el vocabulario, lo que ahorra parámetros y mejora la generalización.

\paragraph{Codificación posicional.} El mecanismo de atención, tal y como se definirá después, es \textit{invariante a permutaciones}: si se barajan los tokens del contexto, la atención produce las mismas salidas salvo por un reordenamiento. Eso es claramente inadecuado para texto, donde el orden importa enormemente. La solución es añadir información posicional a la entrada. En el transformer original, Vaswani et al.~\cite{vaswani2017attention} propusieron codificaciones sinusoidales fijas que suman un patrón determinista al embedding de cada posición, de manera que las funciones trigonométricas de distinta frecuencia codifiquen el índice $i$ en un vector de la misma dimensión $d$ que el embedding. Modelos más recientes, incluidos Qwen, utilizan en cambio \textbf{codificaciones posicionales rotatorias} (\textit{Rotary Position Embeddings}, RoPE)~\cite{su2021roformer}, que en lugar de sumar una señal rotan los vectores de consulta y clave en un plano dependiente de la posición. La ventaja de RoPE es doble: codifica \emph{posiciones relativas} de forma natural y permite extrapolar a longitudes mayores que las vistas durante el entrenamiento sin necesidad de reentrenar. Para este TFG no es necesario profundizar en la matemática de RoPE, pero sí es útil saber que la capacidad del modelo para manejar contextos largos depende en gran medida de esta codificación.

\paragraph{Self-attention como pieza central.} Dentro de cada capa, el mecanismo de atención transforma una secuencia de representaciones de entrada $X \in \mathbb{R}^{n\times d}$ en otra secuencia de la misma forma, de modo que cada posición resulta enriquecida con información relevante extraída de las demás. Para ello, se proyectan linealmente las entradas en tres conjuntos distintos
\[
Q = X W_Q,\qquad K = X W_K,\qquad V = X W_V,
\]
con $W_Q, W_K, W_V \in \mathbb{R}^{d\times d_k}$ matrices aprendidas. La intuición habitual es que cada token formula una \emph{pregunta} en el espacio de queries ($Q$), cada token también expone unas \emph{pistas} en el espacio de keys ($K$) sobre qué información guarda, y los \emph{valores} ($V$) son el contenido que se recupera si la pregunta encaja con la pista. La atención calcula entonces la afinidad entre cada query y cada key mediante un producto escalar, la normaliza con un softmax y la usa para ponderar los valores mediante
\[
\mathrm{Attn}(Q, K, V) \;=\; \mathrm{softmax}\!\left( \frac{Q K^\top}{\sqrt{d_k}} \right) V.
\]
El factor $\sqrt{d_k}$ no es cosmético: sin él, al aumentar $d_k$, las magnitudes de los productos escalares crecen y el softmax se satura, de modo que casi toda la masa se concentra en un único índice y los gradientes se vuelven minúsculos. Dividir por $\sqrt{d_k}$ mantiene la varianza acotada y el softmax bien comportado.

En los LLMs autoregresivos, la atención es además \emph{causal}: a cada posición $i$ se le prohíbe mirar a las posiciones $j > i$ (futuras), porque durante el entrenamiento el objetivo es predecir el token $i+1$ a partir de los tokens $1$ a $i$ y permitir que vea el futuro sería trampa. Esta restricción se implementa con una \textit{máscara} que pone $-\infty$ en las posiciones futuras antes del softmax, de manera que queden probabilidad cero.

Hay una consecuencia práctica muy importante para el coste: la matriz $Q K^\top$ tiene tamaño $n \times n$, donde $n$ es la longitud de la secuencia. Por lo tanto la atención es \textbf{cuadrática con la longitud} tanto en memoria como en cómputo, y cuanto más largo es el contexto, más cara se vuelve. A esto se debe que, aunque los transformers sean conceptualmente muy simples, escalar a contextos de decenas o centenares de miles de tokens sea un problema de ingeniería nada trivial.

\paragraph{Multi-head attention.} La atención tal y como se ha descrito calcula un único patrón de afinidad entre tokens. Los autores del transformer observaron que permite más capacidad expresiva hacer varias ``miradas'' en paralelo, cada una en un subespacio distinto. Esto es la \textbf{multi-head attention}: se divide la dimensión oculta $d$ en $h$ cabezas de tamaño $d_k = d/h$, y cada cabeza aprende sus propias matrices $W_Q^{(i)}, W_K^{(i)}, W_V^{(i)}$. Se ejecuta la atención de forma independiente en cada cabeza y las salidas se concatenan y se proyectan con una matriz adicional $W_O \in \mathbb{R}^{d\times d}$. Intuitivamente, distintas cabezas pueden especializarse en distintas señales (por ejemplo, relaciones sintácticas a corto alcance frente a coherencia temática a largo alcance), lo cual es una de las razones por las que los transformers son tan flexibles. En la práctica, no todas las cabezas llegan a estar igualmente activas, y es frecuente encontrar cabezas redundantes que se pueden podar sin pérdida apreciable de calidad; esto es una de las bases de trabajos recientes de compresión estructural.

\begin{figure}[H]
\centering
\resizebox{0.95\textwidth}{!}{%
\begin{tikzpicture}[
>=Latex,
xbox/.style={draw, rounded corners=2pt, fill=gray!15, minimum width=9mm, minimum height=7mm, font=\scriptsize},
qbox/.style={draw, rounded corners=2pt, fill=bcstudent!35, minimum width=9mm, minimum height=7mm, font=\scriptsize},
kbox/.style={draw, rounded corners=2pt, fill=bcteacher!35, minimum width=9mm, minimum height=7mm, font=\scriptsize},
vbox/.style={draw, rounded corners=2pt, fill=bcrouter!35, minimum width=9mm, minimum height=7mm, font=\scriptsize},
cell/.style={draw, minimum width=6mm, minimum height=6mm, font=\tiny, inner sep=0},
cap/.style={font=\footnotesize, align=center}
]
% --- Columna 1: tokens de entrada X ---
\node[cap] at (0, 3.1) {Entrada};
\foreach \i/\y in {1/2.2, 2/1.4, 3/0.6, 4/-0.2}{
  \node[xbox] (x\i) at (0, \y) {$x_{\i}$};
}
\node[cap] at (0, -1.0) {$X\in\mathbb{R}^{n\times d}$};

% --- Columna 2: proyecciones QKV ---
\node[cap] at (3.0, 3.1) {Proyecciones};
\foreach \i/\y in {1/2.2, 2/1.4, 3/0.6, 4/-0.2}{
  \node[qbox] (q\i) at (2.3, \y) {$q_{\i}$};
  \node[kbox] (k\i) at (3.2, \y) {$k_{\i}$};
  \node[vbox] (v\i) at (4.1, \y) {$v_{\i}$};
}
\foreach \i in {1,2,3,4}{
  \draw[->, gray] (x\i.east) -- (q\i.west);
}
\node[cap] at (3.2, -1.0) {$Q=XW_Q,\ K=XW_K,\ V=XW_V$};

% --- Columna 3: matriz de afinidad softmax(QK^T/sqrt(d_k)) con máscara causal ---
\node[cap] at (7.3, 3.1) {Afinidad (causal)};
% Fila 1: solo la diagonal es activa (el resto queda enmascarado a -inf)
\node[cell, fill=bcmeasure!60] at (6.4, 2.3) {};
\node[cell, fill=gray!10]      at (7.0, 2.3) {};
\node[cell, fill=gray!10]      at (7.6, 2.3) {};
\node[cell, fill=gray!10]      at (8.2, 2.3) {};
% Fila 2
\node[cell, fill=bcmeasure!40] at (6.4, 1.6) {};
\node[cell, fill=bcmeasure!60] at (7.0, 1.6) {};
\node[cell, fill=gray!10]      at (7.6, 1.6) {};
\node[cell, fill=gray!10]      at (8.2, 1.6) {};
% Fila 3
\node[cell, fill=bcmeasure!30] at (6.4, 0.9) {};
\node[cell, fill=bcmeasure!50] at (7.0, 0.9) {};
\node[cell, fill=bcmeasure!70] at (7.6, 0.9) {};
\node[cell, fill=gray!10]      at (8.2, 0.9) {};
% Fila 4
\node[cell, fill=bcmeasure!25] at (6.4, 0.2) {};
\node[cell, fill=bcmeasure!40] at (7.0, 0.2) {};
\node[cell, fill=bcmeasure!55] at (7.6, 0.2) {};
\node[cell, fill=bcmeasure!75] at (8.2, 0.2) {};
\node[cap] at (7.3, -1.0) {$\mathrm{softmax}\!\big(QK^{\top}/\sqrt{d_k}\big)$};

% --- Columna 4: salida ---
\node[cap] at (10.6, 3.1) {Salida};
\foreach \i/\y in {1/2.2, 2/1.4, 3/0.6, 4/-0.2}{
  \node[xbox, fill=bcdata!40] (o\i) at (10.6, \y) {$o_{\i}$};
}
\node[cap] at (10.6, -1.0) {$O = \text{softmax}(\cdot)\, V$};

% Flechas entre bloques
\draw[->, thick] (4.3, 1.0) -- (6.0, 1.0) node[midway, above, font=\scriptsize] {$QK^{\top}$};
\draw[->, thick] (8.5, 1.0) -- (10.0, 1.0) node[midway, above, font=\scriptsize] {$\cdot V$};
\end{tikzpicture}%
}%
\caption{Flujo conceptual de self-attention para una secuencia de cuatro tokens. Las entradas $x_i$ se proyectan linealmente en \emph{queries} $q_i$, \emph{keys} $k_i$ y \emph{values} $v_i$. La matriz $QK^{\top}/\sqrt{d_k}$ se normaliza con softmax para obtener pesos de atención (aquí triangular inferior por la máscara causal) y se aplica a los \emph{values} para producir las salidas $o_i$, cada una contextualizada por los tokens anteriores. En multi-head attention este esquema se replica $h$ veces en subespacios distintos. Fuente: elaboración propia.}
\label{fig:self-attention}
\end{figure}

\paragraph{Bloque completo (atención, feed-forward, residuales y normalización).} Cada una de las $L$ capas de un transformer moderno es, en realidad, una composición de dos sub-bloques. Primero, un bloque de atención multi-cabeza envuelto en una conexión residual y una normalización de capa (\textit{LayerNorm} o \textit{RMSNorm}). Después, un bloque \textbf{feed-forward} (una MLP de dos capas con una activación no lineal de tipo GELU o SwiGLU), de nuevo con residual y normalización. La fórmula esquemática de una capa es
\[
X' = X + \mathrm{Attn}(\mathrm{Norm}(X)),\qquad
X'' = X' + \mathrm{FFN}(\mathrm{Norm}(X')).
\]
Las conexiones residuales cumplen dos funciones importantes: estabilizan el entrenamiento en redes profundas (el gradiente tiene un camino directo que no desaparece al crecer $L$) y permiten que el modelo \emph{opte} por no modificar la representación si así lo prefiere (los pesos pueden aprender una transformación cercana a cero y la residual preserva la entrada). La normalización, por su parte, mantiene controladas las escalas de activaciones capa a capa. En los modelos actuales es habitual la variante \textit{pre-norm} (normalizar \emph{antes} de cada sub-bloque), porque resulta más estable a profundidades grandes que la \textit{post-norm} original.

El bloque feed-forward es donde se concentra una parte muy significativa de los parámetros totales. Una arquitectura típica expande la dimensión oculta a $4d$ o a $8d/3$ (con SwiGLU), aplica la no linealidad y vuelve a proyectar a $d$. Conceptualmente, si la atención se encarga de ``mezclar'' información entre posiciones, el feed-forward se encarga de ``procesarla'' dentro de cada posición por separado. Ambas operaciones son complementarias y ninguna basta por sí sola.

\paragraph{Cabecera y decodificación.} Tras las $L$ capas se obtiene, para cada posición, un vector final $h_i \in \mathbb{R}^{d}$. La cabecera (o \textit{LM head}) es una proyección lineal $W_{\text{out}} \in \mathbb{R}^{d \times |V|}$ que transforma ese vector en un vector de \textbf{logits} sobre el vocabulario. Aplicando un softmax a los logits se obtiene una distribución de probabilidad $p_\theta(t \mid t_1, \dots, t_i)$ sobre el siguiente token.

A partir de esa distribución se elige el próximo token según una de varias estrategias. La más simple es la decodificación \textbf{greedy} (\textit{argmax}), que toma siempre el token con mayor probabilidad; es determinista y es el modo usado en la evaluación de calidad del TFG, precisamente porque elimina una fuente de varianza entre ejecuciones y garantiza que un mismo prompt produce siempre la misma respuesta. La alternativa es el \textbf{sampling}, donde se muestrea de la distribución, a veces truncada mediante \textit{top-$k$} o \textit{top-$p$} (nucleus sampling) y reescalada con una temperatura, para favorecer la diversidad. Este modo es habitual en aplicaciones interactivas donde se desea que el modelo produzca respuestas variadas, pero a cambio introduce no-determinismo que dificulta la comparación experimental. La elección de estrategia de decodificación importa mucho en la práctica: dos modelos idénticos con estrategias de sampling distintas pueden dar perfiles de calidad y de longitud de respuesta muy diferentes, lo cual impone pautas de comparabilidad estrictas. Por eso este TFG fija siempre decodificación greedy al evaluar calidad.

\begin{figure}[H]
\centering
\resizebox{\textwidth}{!}{%
\begin{tikzpicture}[
>=Latex,
tk/.style={draw, rounded corners=2pt, minimum width=10mm, minimum height=7mm, font=\footnotesize},
blk/.style={draw, rounded corners=2pt, fill=blue!10, minimum width=62mm, minimum height=8mm, font=\footnotesize},
kv/.style={draw, rounded corners=2pt, fill=orange!18, minimum width=14mm, minimum height=7mm, font=\footnotesize}
]
% Tokens de entrada
\node[tk] (t1) at (0,0) {$t_1$};
\node[tk, right=3mm of t1] (t2) {$t_2$};
\node[tk, right=3mm of t2] (t3) {$t_3$};
\node[tk, right=3mm of t3] (t4) {$t_4$};
\node[right=3mm of t4, font=\footnotesize] (dots) {$\dots$};
\node[tk, right=3mm of dots] (tn) {$t_n$};

% Bloque transformer centrado respecto a la fila de tokens
\node[blk, below=10mm of t3, xshift=8mm] (attn) {Atención multi-cabeza};
\node[blk, below=3mm of attn] (ffn)  {Feed-forward};
\node[font=\footnotesize, right=3mm of attn] (Lcapas) {$\times L$ capas};

% Flechas tokens -> atención
\foreach \x in {t1,t2,t3,t4,tn}{\draw[->] (\x.south) -- (attn.north);}
\draw[->] (attn) -- (ffn);

% KV cache a la derecha
\node[kv, right=22mm of ffn] (kv) {KV cache};
\draw[->, dashed] (attn.east) -| (kv.north) node[pos=0.25, above, font=\scriptsize] {$K,V$};

% Cabecera + token salida
\node[blk, below=5mm of ffn, fill=green!15] (head) {Cabecera (logits sobre vocabulario)};
\draw[->] (ffn) -- (head);

\node[tk, below=5mm of head, fill=yellow!25] (tn1) {$t_{n+1}$};
\draw[->] (head) -- (tn1);

% Realimentación por la derecha, sin cruzar los bloques
\draw[->, dashed, gray]
  (tn1.east) -- ++(2.2,0)
  node[midway, below, font=\scriptsize] {siguiente paso}
  |- (tn.east);
\end{tikzpicture}%
}%
\caption{Esquema funcional de un transformer autoregresivo: cada token pasa por $L$ capas (atención + feed-forward), las matrices $K$ y $V$ se almacenan en el KV cache para reutilizarse en los siguientes pasos, y la cabecera produce la distribución sobre el vocabulario con la que se muestrea el siguiente token. Fuente: elaboración propia.}
\label{fig:transformer-esquema}
\end{figure}

Para este TFG no hace falta profundizar más en la arquitectura. Basta con tener presente que cada token generado requiere recorrer todas las capas del modelo, y que la parte dominante del coste está en la atención y en las capas feed-forward. En los modelos de 7B a 14B parámetros que aparecen en el proyecto, cada paso forward equivale a varios teraflops de cómputo y a leer cientos de megabytes de memoria HBM por token generado.

\subsection{Inferencia autoregresiva y KV cache}

La inferencia de un LLM tiene dos fases muy distintas entre sí: \textbf{prefill} y \textbf{decode}. Entender sus diferencias es fundamental porque explica por qué un sistema que funciona bien para generar una única respuesta larga puede comportarse mal al servir muchas peticiones en paralelo, y al revés.

\textbf{Prefill (procesar el prompt).} Cuando llega una petición, el modelo procesa el prompt completo de golpe. Todos los tokens del prompt se pasan a la vez por la red, y en cada capa se calculan sus $K$ y $V$. El coste de prefill es aproximadamente lineal con la longitud del prompt y domina en términos de FLOPs: es una operación limitada por cómputo (\textit{compute-bound}) y aprovecha muy bien las GPUs modernas porque se puede paralelizar completamente.

\textbf{Decode (generar tokens).} Tras el prefill, el modelo entra en bucle autoregresivo: genera un token, lo añade al contexto, genera el siguiente y así hasta alcanzar un criterio de parada. En cada paso de decode solo se procesa \textbf{un único token}, pero la atención tiene que mirar a todos los tokens anteriores. Si el KV cache no existiese, habría que recalcular $K$ y $V$ de todo el contexto en cada paso, lo que daría un coste cuadrático con la longitud de la secuencia~\cite{vaswani2017attention}. El KV cache resuelve esto: se almacenan las $K$ y $V$ de los tokens ya procesados y, en cada paso, solo se calculan las del token nuevo.

El resultado es que el decode es una operación limitada por \textbf{memoria} (\textit{memory-bound}): hay que leer el KV cache entero de HBM cada paso, la cantidad de cómputo es pequeña en comparación, y la utilización de la GPU baja mucho frente al prefill. Esta asimetría entre una fase \textit{compute-bound} y otra \textit{memory-bound} es la raíz del problema de eficiencia en LLMs y es lo que han atacado los motores de serving modernos.

\begin{figure}[H]
\centering
\resizebox{\textwidth}{!}{%
\begin{tikzpicture}[
>=Latex,
box/.style={draw, rounded corners=2pt, minimum height=10mm, font=\footnotesize, align=center},
axlab/.style={font=\footnotesize}
]
% Bloques prefill y decode
\node[box, fill=blue!20, minimum width=50mm] (pre) at (0,0) {\textbf{Prefill}\\procesa todo el prompt en paralelo};
\node[box, fill=orange!18, minimum width=85mm, right=3mm of pre] (dec) {\textbf{Decode}\\un token por paso, $n$ pasos secuenciales};

% Eje temporal que cubre todo el ancho
\draw[->, thick] ($(pre.west)+(-0.4,-1.0)$) -- ($(dec.east)+(0.6,-1.0)$) node[right, axlab] {tiempo};

% Anotaciones bajo cada fase
\node[axlab, below=3mm of pre, align=center] (apre) {compute-bound\\alta utilización GPU};
\node[axlab, below=3mm of dec, align=center] (adec) {memory-bound\\lee KV cache en cada paso};
\end{tikzpicture}%
}%
\caption{Fases de inferencia en un LLM: prefill procesa todo el prompt de una vez (compute-bound), decode genera los tokens de uno en uno leyendo el KV cache en cada paso (memory-bound). Fuente: elaboración propia.}
\label{fig:prefill-decode}
\end{figure}

Un detalle importante es el \textbf{tamaño del KV cache}. Por cada token, por cada capa y por cada cabeza se almacena un vector $K$ y un vector $V$. En un modelo de unos 7B parámetros, 32 capas y 32 cabezas de 128 dimensiones en FP16, el KV cache ocupa aproximadamente 0,5 MB por token. Un prompt de 4\,000 tokens ya consume 2 GB solo en caché, y una petición con un contexto de 32\,000 tokens puede llegar a 16 GB. En una GPU H100 de 80 GB esto es mucho: después de reservar memoria para los pesos del modelo (unos 14 GB en FP16 para un 7B), para los activations del prefill y para márgenes de seguridad, quedan pocos gigabytes para KV cache. Servir varias peticiones largas a la vez rápidamente satura la memoria.

\subsection{Métricas estándar de servicio}

El TFG mide el coste de inferencia con cuatro métricas estándar en la literatura de serving. Es importante definirlas con precisión porque luego se usan de forma recurrente en los capítulos de diseño y resultados.

La primera es el \textbf{TTFT} (\textit{Time To First Token}), que mide el tiempo transcurrido entre que el cliente envía la petición y recibe el primer token de respuesta. Incluye la latencia de red, la espera en cola interna del motor y la fase de prefill, y es la métrica que más pesa en la percepción de \textit{responsividad} de una aplicación interactiva.

La segunda es el \textbf{TPOT} (\textit{Time Per Output Token}), que cuantifica el tiempo medio que tarda cada token de decode una vez iniciada la generación. En régimen estable depende sobre todo del ancho de banda de memoria y del nivel de batch que esté sirviendo el motor. Se suele reportar también su inverso, expresado en tokens por segundo por petición, para facilitar la comparación.

La tercera es la \textbf{latencia total}, que abarca la duración completa de una petición desde la llegada hasta la emisión del último token. Se puede descomponer como
\[
\ell_{\text{total}} \;\approx\; \text{TTFT} + (n_{\text{out}} - 1)\cdot \text{TPOT},
\]
donde $n_{\text{out}}$ es el número de tokens generados. Se reportan percentiles (p50, p95, p99) en lugar de la media, porque la media no captura bien el comportamiento de la cola de la distribución.

Finalmente, el \textbf{throughput} mide los tokens generados por segundo \emph{agregados} a nivel de sistema, no por petición individual. Un throughput alto suele implicar un TPOT por petición peor, porque el motor está sirviendo más peticiones simultáneas y compite por los mismos recursos de memoria y cómputo. Este es el clásico compromiso entre latencia y throughput propio de cualquier sistema de colas.

\begin{table}[H]
\centering
\caption{Métricas estándar de servicio de un LLM usadas en este TFG. Fuente: elaboración propia a partir de~\cite{vllm2023, vllm-docs}.}
\label{tab:metricas-servicio}
\footnotesize
\begin{tabularx}{\textwidth}{>{\raggedright\arraybackslash}p{2.4cm} >{\raggedright\arraybackslash}X >{\raggedright\arraybackslash}p{4cm}}
\toprule
\textbf{Métrica} & \textbf{Qué mide} & \textbf{Interpretación} \\
\midrule
TTFT & Tiempo hasta el primer token & Responsividad percibida por el usuario. Muy afectada por el prefill. \\
\addlinespace
TPOT & Tiempo por token de decode & Velocidad en régimen estable; depende del batch y del ancho de banda de memoria. \\
\addlinespace
Latencia total & Duración completa de la petición & Relevante para aplicaciones no interactivas; se reporta en percentiles. \\
\addlinespace
Throughput & Tokens/segundo agregados del sistema & Coste operativo: tokens servibles por euro de GPU. \\
\addlinespace
Memoria GPU & VRAM ocupada por pesos + KV cache + activations & Determina cuántos modelos caben en paralelo y qué \textit{context length} admite el sistema. \\
\bottomrule
\end{tabularx}
\end{table}

Estas cinco métricas forman la base de la evaluación de la Parte~\ref{part:diseno-sistema}. En el protocolo de medida que se describe más adelante (warmup, ventana de 100 peticiones, tres repeticiones y control de variables) se reportan siempre p50 y p95 de TTFT y latencia, y la media de throughput, porque esa combinación captura tanto el caso típico como la cola.

%%%%%%%%%%%%%%%%%%%%%%%%%%%%%%%%%%%%%%%%%%%%%%%%%%%%%%%%%%%%%%%%%
\section{Optimización del modelo. Panorama de técnicas, fine-tuning y LoRA}
\label{sec:optim-modelo}

\subsection{Panorama de compresión, adaptación y serving}
\label{sec:panorama-opt}

Antes de entrar a fondo en Knowledge Distillation, conviene situarla dentro del mapa más amplio de técnicas que la comunidad usa hoy para hacer más baratos los LLMs. Esto permite entender por qué se ha elegido KD como pieza habilitadora del TFG y, sobre todo, por qué no es incompatible con el resto, pues en un despliegue real las distintas familias se combinan.

\subsubsection{Compresión de modelos}

La \textbf{cuantización} reduce la precisión de los pesos (y a veces de las activaciones) de FP16 a INT8, INT4 o formatos mixtos. El resultado es un modelo que pesa varias veces menos en memoria y cuyo forward es más rápido en hardware con soporte para aritmética de baja precisión, a cambio de una pérdida de calidad que depende del método y del modelo. Es la técnica más barata de aplicar porque no requiere reentrenamiento (en su variante \textit{post-training}) y suele aparecer como primer paso en cualquier pipeline de serving.

El \textbf{pruning} elimina conexiones (pruning no estructurado) o unidades completas de la red (pruning estructurado) con la idea de que buena parte de los parámetros son redundantes. El pruning estructurado es el más útil desde el punto de vista de coste porque produce modelos más pequeños que aprovechan GPUs estándar sin necesidad de kernels específicos para matrices sparse. Suele combinarse con una fase de recuperación vía fine-tuning.

La \textbf{destilación de modelos} (\textit{knowledge distillation}, KD) es la técnica sobre la que se articula este TFG y se desarrolla en la sección siguiente. A diferencia de cuantización y pruning, que actúan sobre el mismo modelo comprimiéndolo, KD entrena un modelo distinto (más pequeño) para imitar el comportamiento del grande. Eso la hace conceptualmente más cara en tiempo de entrenamiento pero abre un grado de libertad adicional: permite elegir libremente la arquitectura del student, cambiar la familia de modelos, o incluso destilar varios students de distinto tamaño a partir del mismo teacher.

\subsubsection{PEFT y arquitecturas dispersas}

Las técnicas de \textbf{PEFT} (\textit{parameter-efficient fine-tuning}), como LoRA y sus derivadas, no reducen el coste de inferencia del modelo base. Sirven para adaptar un LLM grande a una tarea concreta entrenando apenas una fracción de los pesos. Son relevantes para el TFG porque, si la destilación se hace a partir de un teacher \textit{ya adaptado} a un dominio (con LoRA) o incluso si el student se afina con LoRA después de destilar, la cascada resultante mantiene los beneficios de la adaptación sin disparar el coste de entrenamiento.

Una vía distinta son las arquitecturas con \textbf{mixture-of-experts} (MoE): modelos con muchos parámetros totales pero donde cada token solo activa un subconjunto de \textit{expertos}. En inferencia, el coste efectivo es mucho menor que el tamaño nominal del modelo. MoE es una forma de compresión estructural que convive perfectamente con KD (se puede destilar un MoE grande a un denso pequeño), pero no se explora en este TFG para mantener acotado el alcance.

\subsubsection{Optimización a nivel de serving}

Las técnicas anteriores actúan sobre el modelo. A nivel de \textbf{serving} aparece otra capa de optimizaciones que no cambian los pesos pero sí cómo se ejecutan, entre ellas \textit{continuous batching}, PagedAttention, FlashAttention y \textit{speculative decoding}. Todas ellas se detallan más adelante en la Sección~\ref{sec:serving-teorico}. Lo importante en este punto es que son \textbf{ortogonales} a KD: un student destilado se sigue beneficiando de estas optimizaciones, y de hecho el TFG las aprovecha al servir los modelos vía vLLM.

\subsubsection{Dónde encaja el TFG}

Con este mapa en mente, la elección de KD como pieza central se justifica por tres motivos. Primero, KD produce un modelo nuevo que puede servirse en paralelo al teacher, condición necesaria para construir cascadas. Segundo, KD es compatible con todas las demás técnicas: el student destilado puede después cuantizarse, servirse con PagedAttention y beneficiarse de speculative decoding; nada de lo que se hace en este TFG cierra esas puertas. Y tercero, KD permite estudiar el balance entre calidad y coste con un control fino, porque al variar el tamaño del student y la receta de destilación se obtienen puntos distintos de la frontera de Pareto.

\subsection{Fine-tuning y adaptación eficiente}

Aunque el núcleo metodológico del TFG es la destilación de conocimiento, en la práctica casi todos los experimentos de destilación se implementan como una forma particular de \textit{fine-tuning} del modelo pequeño. Entender qué es realmente el \textit{fine-tuning}, qué se modifica dentro del modelo y qué alternativas eficientes existen es por tanto un prerrequisito para comprender el pipeline experimental que se describe en la parte de diseño. A continuación se presenta primero la idea general del \textit{fine-tuning}, después la variante de \textit{full fine-tuning} y finalmente la familia de técnicas PEFT, con especial atención a LoRA por ser la que se emplea en una de las ramas experimentales.

\subsubsection{Qué está pasando realmente durante el fine-tuning}

Un LLM llega a nosotros tras una fase de preentrenamiento masivo sobre grandes corpus de texto, donde ha aprendido a modelar la distribución general del lenguaje. El resultado es un conjunto de parámetros $\theta_{\text{pre}}$ que ya contiene una gran cantidad de conocimiento lingüístico, factual y de razonamiento. El fine-tuning es, simplemente, una segunda fase de entrenamiento a partir de $\theta_{\text{pre}}$ con un corpus y un objetivo específicos, con la intención de especializar el modelo sin destruir lo que ya sabe. Los modelos \textit{Instruct} que se usan como punto de partida en este TFG (por ejemplo \texttt{Qwen2.5-1.5B-Instruct}, \texttt{-7B-Instruct} y \texttt{-14B-Instruct}) son el resultado de un fine-tuning supervisado sobre datos de instrucción, seguido en muchos casos de una etapa de alineamiento por preferencias (DPO, RLHF).

Desde el punto de vista matemático, el fine-tuning es un problema de optimización estándar: dado un dataset $\mathcal{D} = \{(x_i, y_i)\}$ y una función de pérdida $\mathcal{L}$ (habitualmente entropía cruzada token a token para la predicción del siguiente token), se minimiza
\[
\theta^{\star} \;=\; \arg\min_{\theta} \; \mathbb{E}_{(x,y)\sim \mathcal{D}}\big[\, \mathcal{L}(f_\theta(x), y)\,\big],
\]
partiendo de $\theta = \theta_{\text{pre}}$ e iterando con un optimizador basado en gradientes, normalmente AdamW. El éxito del fine-tuning descansa sobre dos propiedades del preentrenamiento que conviene hacer explícitas. La primera es que \textbf{la inicialización ya es buena}: el modelo parte de una zona razonable del espacio de parámetros en la que los gradientes tempranos tienen ya una dirección útil, y el optimizador no necesita construir el conocimiento desde cero. La segunda es que \textbf{se puede aprender con poco dato}: como la base ya codifica conocimiento general del lenguaje, para especializarla en un dominio concreto (por ejemplo, matemáticas con GSM8K) basta un corpus comparativamente pequeño, del orden de miles a decenas de miles de ejemplos, para obtener resultados apreciables. Entrenar desde cero un modelo equivalente exigiría órdenes de magnitud más datos y cómputo.

\paragraph{Qué parámetros se actualizan.} En el fine-tuning \emph{completo}, todos los parámetros del modelo son susceptibles de ser modificados. Eso incluye los embeddings, todas las matrices de atención ($W_Q, W_K, W_V, W_O$) de todas las capas, todos los pesos del feed-forward, los parámetros de las normalizaciones y la cabecera de salida. Para un modelo de 7B eso significa literalmente 7\,000 millones de parámetros en modo ``entrenable'', cada uno con su gradiente y, dependiendo del optimizador, uno o dos momentos adicionales por parámetro en memoria. El coste computacional y de memoria crece muy rápido con el tamaño del modelo (véase el apartado siguiente).

\paragraph{Variantes según el tipo de objetivo.} Conviene distinguir tres regímenes de fine-tuning que aparecen con nombre propio en la literatura porque se comportan de manera cualitativamente distinta. El primero es el \textbf{continual pretraining}, donde se sigue entrenando con la misma tarea del preentrenamiento (predicción del siguiente token) pero sobre un corpus nuevo, típicamente especializado en un dominio particular; su propósito es aportar conocimiento factual o estilístico al modelo sin alterar su formato de salida. El segundo es el \textbf{supervised fine-tuning} (SFT), donde se usan pares explícitos de entrada y respuesta deseada y se entrena al modelo a generar la respuesta dada la entrada. Es el régimen natural para tareas de instrucción y diálogo y, en el contexto de este TFG, es el mecanismo mediante el cual el student aprende a reproducir las respuestas del teacher; la destilación de salidas es, a efectos prácticos, un SFT sobre un corpus sintético generado por el modelo grande. El tercero es el \textit{preference tuning}, que agrupa técnicas como RLHF y DPO, donde el modelo se entrena para preferir ciertas respuestas sobre otras usando señales de preferencia humana o sintéticas. Aunque este régimen no se emplea directamente en el TFG (no se entrena con preferencias), es útil saber que las versiones \textit{Instruct} de los modelos Qwen que se usan como punto de partida ya incorporan internamente una adaptación de este tipo.

\paragraph{Ventajas y riesgos del fine-tuning.} El fine-tuning es el mecanismo más directo y probado para especializar un LLM: permite mejorar de forma medible la calidad en un dominio concreto, corregir patrones indeseados como la verbosidad excesiva o el formato de salida erróneo, y transferir conocimiento de un modelo grande a uno pequeño cuando el corpus de entrenamiento ha sido generado por el grande.

A cambio, presenta varios riesgos que son particularmente relevantes en el contexto de este proyecto. El más conocido es el \textbf{olvido catastrófico}: actualizar todos los parámetros sobre un corpus nuevo puede degradar capacidades generales del modelo que no estaban representadas en los datos de ajuste. Un student afinado agresivamente sobre matemáticas puede perder calidad en diálogo general, lo que conviene monitorizar si el modelo va a usarse en un sistema que también recibe consultas de otros dominios. Otro riesgo es el \textbf{sobreajuste}: con corpus pequeños y muchos parámetros entrenables, el modelo puede memorizar los ejemplos de entrenamiento y perder capacidad de generalización. En los benchmarks que se utilizan en este TFG (GSM8K, MATH), donde el conjunto de evaluación es estrictamente independiente del de entrenamiento, el sobreajuste se detecta como una divergencia entre la pérdida de entrenamiento y la accuracy de evaluación. Por último, el fine-tuning es \textbf{sensible a los hiperparámetros}: tasas de aprendizaje demasiado altas pueden desestabilizar el modelo de manera irreversible. La literatura recomienda tasas de aprendizaje órdenes de magnitud menores que las del preentrenamiento, ventanas de \textit{warmup} cortas y un número limitado de épocas para evitar la degradación.

\subsubsection{Fine-tuning completo y su coste}

En el \textit{full fine-tuning} (FT), todos los parámetros del modelo son entrenables. Es la opción más expresiva y la que tiene, en principio, mayor capacidad para asimilar información nueva. Pero también es la más cara en memoria y en cómputo, y para entender por qué conviene detenerse brevemente en lo que ocupa la GPU durante un \textit{training step}.

\paragraph{Desglose de la memoria en entrenamiento.} Para un modelo de $N$ parámetros entrenado en FP16 con AdamW y \textit{mixed precision}, la memoria de GPU necesaria se puede descomponer de forma aproximada en cuatro grandes partidas. Los \textbf{pesos del modelo} ocupan $2N$ bytes en FP16. Los \textbf{gradientes} requieren otros $2N$ bytes (en FP16 o FP32, según la implementación). Los \textbf{estados del optimizador} son la partida más pesada: AdamW mantiene dos momentos (primer y segundo momento exponencial) por cada parámetro, lo que en FP32 supone $8N$ bytes adicionales; a menudo se conserva además una copia maestra de los pesos en FP32 por estabilidad numérica, que añade otros $4N$ bytes. Por último, las \textbf{activaciones intermedias} se almacenan durante el forward para poder calcular los gradientes en el backward. Su tamaño crece con el \textit{batch size}, la longitud de secuencia y la profundidad del modelo; con la técnica de \textit{gradient checkpointing} se recalculan parte de estas activaciones en lugar de almacenarlas, a cambio de aproximadamente un 30\,\% más de cómputo.

Sumando las partidas fijas, el orden de magnitud es de $16N$ a $20N$ bytes \emph{sin contar activaciones}. Para un modelo de 7B parámetros esto son unos 120--140~GB, cifra que no cabe en una única H100 de 64--80~GB de VRAM. Por eso cualquier fine-tuning completo de modelos modernos exige estrategias de paralelismo (\textit{multi-GPU} con DDP, FSDP o DeepSpeed ZeRO) para repartir esta memoria entre varios dispositivos. En la práctica, esto se traduce en solicitudes de más GPU-hora en el clúster y en configuraciones de entrenamiento más complejas, factores que pesan directamente en el presupuesto de cómputo del TFG.

\paragraph{Cuándo es necesario.} A pesar del coste, hay escenarios donde el FT tiene ventajas claras: cuando se dispone de abundante cómputo, cuando el corpus objetivo es suficientemente grande, y cuando se busca maximizar calidad sin restricciones. En el TFG, la rama experimental \textit{v6} se resolvió con FT completo precisamente porque era la configuración más ambiciosa y se quería saber, con la mejor herramienta posible, cuánto podía aprender un student de 1{,}5B de las salidas del teacher de 7B. Este FT fue el experimento más costoso del TFG en términos de GPU-hora y es el que fija la cota superior empírica de calidad alcanzable por destilación con los recursos disponibles.

\subsubsection{PEFT y la aparición de LoRA}

La idea detrás de PEFT es que muchas veces no hace falta actualizar todos los parámetros del modelo para adaptarlo a una tarea nueva. Si el modelo preentrenado ya es lo bastante bueno, lo que se necesita es una ``pequeña corrección'' concentrada en pocas direcciones del espacio de parámetros. PEFT agrupa un conjunto de técnicas que comparten una estrategia común, consistente en congelar la mayor parte del modelo y entrenar solo un \textbf{subconjunto reducido} de parámetros o, alternativamente, añadir un pequeño número de parámetros nuevos que se entrenan junto al resto del modelo congelado. Entre las más conocidas están \textit{adapter layers}, \textit{prompt tuning}, \textit{prefix tuning} y, muy especialmente, LoRA.

\paragraph{Motivación cuantitativa.} La literatura y experiencia práctica coinciden en un hecho empírico importante: la \emph{dirección efectiva} del gradiente durante el fine-tuning vive en un subespacio de dimensión mucho menor que la total. En términos intuitivos, la diferencia entre un modelo preentrenado y su versión ajustada a una tarea no es ``ruido denso en todos los pesos'', sino algo más parecido a una actualización de \emph{bajo rango}. Esta observación, formalizada por Hu et al. en el trabajo original de LoRA~\cite{hu2021lora}, es la base sobre la que se construyen los métodos PEFT modernos.

\subsubsection{LoRA (Low-Rank Adaptation)}
\label{sec:finetuning-lora}

\paragraph{Idea central.} Considérese una capa lineal cualquiera del transformer, representada por una matriz $W_0 \in \mathbb{R}^{d_{\text{in}} \times d_{\text{out}}}$. En el FT completo, esta matriz se actualiza como $W_0 + \Delta W$, con $\Delta W$ aprendida durante el entrenamiento. LoRA hace la hipótesis de que esta actualización $\Delta W$ puede aproximarse bien mediante un producto de dos matrices mucho más pequeñas:
\[
\Delta W \;\approx\; B\,A,\qquad A\in\mathbb{R}^{d_{\text{in}}\times r},\; B\in\mathbb{R}^{r\times d_{\text{out}}},
\]
con $r \ll \min(d_{\text{in}}, d_{\text{out}})$. El parámetro $r$ se llama \textbf{rango} de la adaptación. Durante el fine-tuning, $W_0$ se \textbf{congela} y solo se aprenden $A$ y $B$. El forward de la capa pasa de $y = W_0 x$ a
\[
y \;=\; W_0\, x \;+\; B\,A\, x.
\]
Al inicio del entrenamiento se inicializa $A$ con valores pequeños y $B = 0$, de modo que $\Delta W = 0$ y la salida inicial coincide con la del modelo base; el ajuste emerge poco a poco a medida que $A$ y $B$ se desvían de esa configuración.

\paragraph{Cuántos parámetros se entrenan.} Si la matriz original tiene $d_{\text{in}} d_{\text{out}}$ parámetros, la actualización LoRA añade $r(d_{\text{in}} + d_{\text{out}})$. Para $d = 4096$ y $r = 8$, pasamos de 16{,}8M de parámetros por capa a solo 65K, lo que supone una reducción de dos órdenes de magnitud. Cuando LoRA se aplica únicamente a las matrices de atención ($W_Q, W_K, W_V, W_O$) de todas las capas (configuración habitual), el número total de parámetros entrenables en un modelo de 7B suele rondar unos pocos millones, frente a los 7\,000 millones del FT completo.

\paragraph{Por qué reduce drásticamente el coste.} La caída del coste proviene sobre todo del tamaño del conjunto entrenable. Los gradientes solo se calculan (y se almacenan) para $A$ y $B$. Los momentos del optimizador, que eran la mayor parte del desglose de memoria, se reducen proporcionalmente. Y como el modelo base no cambia, puede mantenerse en formatos más agresivos de cuantización (por ejemplo, 4-bit en QLoRA) sin afectar a la calidad final, lo que reduce todavía más la memoria total. En términos prácticos, un ajuste LoRA de un modelo de 7B puede ejecutarse en una sola H100, cosa que no es factible con FT completo.

\paragraph{Qué se entrena y qué no.} Es importante retener con claridad el esquema: los pesos originales $W_0$ quedan exactamente igual que en el modelo base; lo único que aprende el entrenamiento es el contenido de las matrices $A$ y $B$ de cada capa adaptada. Eso implica que el conocimiento preexistente del modelo está, literalmente, intacto: no puede ser destruido porque los pesos que lo codifican no se tocan. La adaptación se superpone encima en forma de una corrección de bajo rango, y puede ser activada o desactivada simplemente quitando las matrices $A,B$. Esta propiedad se aprovecha a veces en producción para servir varios clientes con la misma base y adaptadores distintos, o para hacer \textit{A/B testing} entre versiones.

\paragraph{Merge en inferencia.} Una propiedad técnica relevante para este TFG es que, una vez entrenadas las matrices LoRA, la corrección $BA$ puede \textbf{fusionarse} (\textit{merge}) con la matriz original en tiempo de carga:
\[
W_{\text{merged}} = W_0 + B\,A.
\]
El modelo resultante es, a efectos de inferencia, \emph{indistinguible} de un modelo con $W_{\text{merged}}$ como peso entrenado, sin ningún coste adicional en el \textit{forward}. Esto significa que LoRA no impone penalización alguna en el serving: las ramas experimentales del TFG que usan LoRA se sirven en producción exactamente como si hubieran sido ajustadas con FT completo.

\paragraph{Trade-offs de LoRA frente a FT completo.} Entender el equilibrio entre ambos enfoques es importante porque la elección afecta directamente a los resultados que se reportan en la parte experimental.

En cuanto a \textbf{calidad}, LoRA captura empíricamente la mayor parte de la mejora que obtendría un FT completo en tareas donde la adaptación requerida es moderada. Cuando el dominio objetivo se aparta mucho de lo que el modelo base ya sabe, el FT completo tiene ventaja porque puede reconfigurar en profundidad más representaciones internas de las que caben en una corrección de bajo rango. En el caso de este TFG, GSM8K y MATH son dominios donde los modelos preentrenados ya codifican cierta competencia aritmética y algebraica, y la mejora que busca la destilación es relativamente moderada, lo que explica que LoRA dé resultados competitivos.

LoRA introduce también dos \textbf{hiperparámetros nuevos}: el rango $r$, que controla la expresividad de la corrección, y un factor de escala $\alpha$ (el \textit{scaling factor}), de manera que la contribución efectiva de cada adaptador a la capa es $(\alpha/r)\, B A\, x$. Ajustar estos valores requiere experimentación, pero el espacio de búsqueda es pequeño en comparación con el del FT completo.

El \textbf{coste de entrenamiento} es muy inferior al del FT completo. Una iteración de LoRA puede ser varias veces más rápida y consumir hasta un orden de magnitud menos de memoria, lo que permite realizar múltiples experimentos exploratorios dentro de la misma cuota de GPU. Finalmente, el \textbf{riesgo de olvido catastrófico} es muy reducido, porque los pesos originales del modelo permanecen intactos y la corrección es pequeña en magnitud.

\paragraph{Relación con la destilación en este TFG.} En el pipeline experimental del TFG aparecen dos tipos de ajuste diferenciados: ramas tempranas basadas en \emph{SFT sobre respuestas del teacher con LoRA} (por su bajo coste, apropiado para iterar rápidamente sobre configuraciones), y una rama final basada en \emph{FT completo} del student sobre el mismo tipo de datos, concebida como la configuración más ambiciosa una vez fijado el resto del diseño. La comparación entre ambas es interesante no solo por el resultado numérico sino porque ilustra el compromiso típico entre eficiencia y capacidad expresiva al que el TFG se enfrenta de manera continua.

\begin{figure}[H]
\centering
\resizebox{0.95\textwidth}{!}{%
\begin{tikzpicture}[
>=Latex,
frz/.style={draw, rounded corners=2pt, fill=gray!25, minimum width=28mm, minimum height=9mm, font=\footnotesize, align=center},
trn/.style={draw, rounded corners=2pt, fill=bcstudent!45, minimum width=28mm, minimum height=9mm, font=\footnotesize, align=center},
lor/.style={draw, rounded corners=2pt, fill=bcrouter!40, minimum width=16mm, minimum height=8mm, font=\scriptsize, align=center},
lbl/.style={font=\scriptsize, align=center}
]
% Cabecera con etiquetas
\node[font=\bfseries\footnotesize] at (0, 3.1) {Fine-tuning completo};
\node[font=\bfseries\footnotesize] at (8.2, 3.1) {LoRA};

% --- Lado izquierdo: FT completo ---
\node[trn] (w1) at (0, 2.1) {$W_0^{(1)}$ entrenable};
\node[trn] (w2) at (0, 0.9) {$W_0^{(2)}$ entrenable};
\node[trn] (w3) at (0, -0.3) {$W_0^{(3)}$ entrenable};
\node[lbl] at (0, -1.7) {Todos los parámetros\\ se actualizan};

% --- Lado derecho: LoRA ---
\node[frz] (fw1) at (7.2, 2.1) {$W_0^{(1)}$ congelada};
\node[frz] (fw2) at (7.2, 0.9) {$W_0^{(2)}$ congelada};
\node[frz] (fw3) at (7.2, -0.3) {$W_0^{(3)}$ congelada};

\node[lor] (ab1) at (10.6, 2.1) {$B^{(1)}A^{(1)}$};
\node[lor] (ab2) at (10.6, 0.9) {$B^{(2)}A^{(2)}$};
\node[lor] (ab3) at (10.6, -0.3) {$B^{(3)}A^{(3)}$};

\draw[->, thick] (fw1.east) -- node[above, font=\scriptsize] {$+$} (ab1.west);
\draw[->, thick] (fw2.east) -- node[above, font=\scriptsize] {$+$} (ab2.west);
\draw[->, thick] (fw3.east) -- node[above, font=\scriptsize] {$+$} (ab3.west);

\node[lbl] at (8.9, -1.7) {Solo las matrices $A,B$\\ de bajo rango se entrenan};
\end{tikzpicture}%
}%
\caption{Esquema conceptual de fine-tuning completo (izquierda) frente a LoRA (derecha). En el FT, todas las matrices $W_0^{(\ell)}$ de cada capa son parámetros entrenables. En LoRA, las $W_0^{(\ell)}$ permanecen congeladas y solo se entrenan pares de matrices $A,B$ de rango bajo, cuya suma $BA$ constituye la corrección efectiva sobre cada capa. Fuente: elaboración propia a partir de~\cite{hu2021lora}.}
\label{fig:lora-vs-ft}
\end{figure}

%%%%%%%%%%%%%%%%%%%%%%%%%%%%%%%%%%%%%%%%%%%%%%%%%%%%%%%%%%%%%%%%%
\section{Knowledge Distillation}
\label{sec:knowledge-distillation}

La \textit{knowledge distillation} (KD) es la técnica que hace posible la cascada de este TFG, porque sin un student que imite razonablemente bien al teacher, la idea de escalar al grande solo cuando haga falta pierde todo su sentido. En esta sección se explica qué es la destilación, cómo se formula matemáticamente y qué variantes se usan hoy en día, con especial atención a la que se emplea en los experimentos.

\subsection{Definición y motivación}

La idea central fue propuesta por Hinton, Vinyals y Dean en 2015~\cite{hinton2015distilling} para modelos de clasificación, y ha sido extendida desde entonces a prácticamente cualquier tipo de red neuronal, incluidos los LLMs generativos. El punto de partida es casi filosófico. Cuando un modelo grande ha sido entrenado sobre un corpus masivo, la información que ha adquirido no reside únicamente en cuál es su predicción correcta, sino también (y esto es lo importante) en \emph{cómo distribuye la probabilidad entre las alternativas incorrectas}. Esa distribución es mucho más rica que una simple etiqueta \emph{one-hot}.

Para entender por qué, conviene detenerse un momento en un ejemplo. Supongamos que un modelo grande se enfrenta a un problema de aritmética que debería resolverse con 42. Una etiqueta dura solo diría ``la respuesta es 42''. El teacher, en cambio, dirá algo más parecido a: ``mi respuesta es 42 con probabilidad 0,85; 24 con 0,08; 12 con 0,04; el resto del vocabulario es esencialmente cero''. Esa estructura secundaria no es ruido: contiene información sobre qué errores son \emph{plausibles} (24 y 12 son números del orden correcto y con quizá algún paso intermedio mal hecho) y cuáles son \emph{imposibles} (99 o una letra). Hinton llamó a esa información los \textit{soft labels} o \textit{dark knowledge}, por contraste con las hard labels del entrenamiento clásico. Un modelo pequeño que aprenda a imitar esas distribuciones suaves está recibiendo, en cada ejemplo, una señal mucho más densa que la que le daría una etiqueta one-hot. En lugar de un único bit de información por token, recibe una descripción probabilística completa sobre el vocabulario entero.

La consecuencia empírica es significativa. Un student entrenado con distilación suele recuperar una fracción muy apreciable de la calidad del teacher, por encima de la que se conseguiría entrenando el mismo modelo pequeño desde cero con los mismos datos. Esto se explica por dos factores complementarios. Primero, el student accede a una \emph{función objetivo mejor condicionada}: la distribución del teacher ya es una aproximación razonable del comportamiento deseado, y encajar en ella es más fácil que deducirla solo a partir de etiquetas binarias. Segundo, las soft labels contienen implícitamente conocimiento sobre la \emph{estructura} del problema: qué clases se parecen entre sí, qué errores son ``cercanos'' y cuáles ``lejanos''. El student no tiene que reconstruir esa geometría por sí solo, la recibe ya extraída.

Esta propiedad es especialmente valiosa en LLMs, donde entrenar un modelo de decenas de miles de millones de parámetros desde cero con tokens suficientes cuesta decenas o centenares de miles de horas de GPU. Destilar un student a partir de un teacher ya entrenado, sobre un corpus relativamente modesto, es órdenes de magnitud más barato y, si se hace con cuidado, conserva buena parte del comportamiento del modelo grande, especialmente en dominios bien cubiertos por los datos de destilación. Por eso la destilación aparece una y otra vez en la literatura reciente como pieza central del pipeline de entrega de modelos a producción, y por eso es la técnica habilitadora de la cascada en este TFG.

\subsection{Formulación clásica de la destilación por respuesta}

La versión original de Hinton se formula sobre las distribuciones de salida del teacher y del student. Sea un input $x$ y su etiqueta verdadera $y$. Sea $z^t(x)$ el vector de logits del teacher y $z^s(x)$ el del student. Sobre esos logits se aplica un softmax con \emph{temperatura} $T$:
\[
p^t_i(x; T) = \frac{\exp\!\big(z^t_i(x)/T\big)}{\sum_j \exp\!\big(z^t_j(x)/T\big)}, \qquad
p^s_i(x; T) = \frac{\exp\!\big(z^s_i(x)/T\big)}{\sum_j \exp\!\big(z^s_j(x)/T\big)}.
\]
La pérdida de destilación es una combinación convexa de dos términos: una pérdida de entropía cruzada contra la etiqueta real y una divergencia de Kullback--Leibler entre las distribuciones del teacher y del student,
\[
\mathcal{L}_{\mathrm{KD}}(x, y) \;=\; \alpha \, \underbrace{\mathrm{CE}\!\big(y,\, p^s(x;1)\big)}_{\text{pérdida supervisada}} \;+\; (1-\alpha)\, T^2 \, \underbrace{\mathrm{KL}\!\big(p^t(x;T)\,\|\,p^s(x;T)\big)}_{\text{pérdida de destilación}}.
\]

La expresión parece inocente pero cada pieza tiene un motivo, y conviene desgranarla con cuidado porque los mismos ingredientes (KL, temperatura, balance) aparecen modificados en todas las variantes modernas.

\paragraph{Qué mide la divergencia KL y por qué aquí.} La divergencia de Kullback--Leibler entre dos distribuciones $p$ y $q$ sobre el mismo soporte se define como
\[
\mathrm{KL}(p\,\|\,q) \;=\; \sum_{i} p_i \log\!\frac{p_i}{q_i}.
\]
No es simétrica ($\mathrm{KL}(p\|q)\neq \mathrm{KL}(q\|p)$) y admite una interpretación muy natural: mide cuántos bits (o \emph{nats}) de información se pierden de media si se usa $q$ para codificar muestras que en realidad provienen de $p$. Al ponerla en la pérdida como $\mathrm{KL}(p^t\|p^s)$, estamos imponiendo al student que su distribución $p^s$ se acerque a la del teacher $p^t$ en el sentido de \emph{cubrir la masa de probabilidad del teacher}: si el teacher asigna probabilidad alta a una clase, el student será penalizado fuertemente por no seguirlo; si el teacher asigna probabilidad cero, el término correspondiente desaparece y no impone restricción alguna. Esta asimetría es deseable en distilación, porque lo que nos importa es que el student \emph{no olvide} lo que el teacher cree probable, no al revés.

\paragraph{El papel de la temperatura $T$.} Al dividir los logits por $T>1$ antes del softmax, la distribución se ``suaviza'': las probabilidades más altas bajan y las más bajas suben. Es un mecanismo crucial porque en modelos bien entrenados la distribución del teacher suele estar muy concentrada en la clase correcta (por ejemplo, 0,99 en una clase y 0,01 repartido en el resto). En esa situación, el softmax original \emph{aplasta} toda la información secundaria: las probabilidades relativas entre las clases incorrectas son indistinguibles numéricamente y no aportan señal apenas. Aumentar $T$ tiene el efecto de magnificar esas diferencias y ponerlas en un rango donde el student puede realmente aprenderlas. A cambio, al suavizar hacemos también que los gradientes del student se vuelvan más pequeños (las diferencias entre probabilidades son más sutiles), y por eso el segundo término se multiplica por $T^2$, para compensar esta reducción y mantener el orden de magnitud del gradiente invariante con $T$~\cite{hinton2015distilling}. En la práctica, valores razonables son $T\in[2,10]$ según la tarea y el modelo; $T$ muy grande tiende a ``aplastarlo todo'' por exceso y deja de aportar información.

\paragraph{El papel de $\alpha$.} Es el hiperparámetro que balancea cuánto pesa aprender de la verdad objetiva (las hard labels) frente a cuánto pesa imitar al teacher (las soft labels). Valores bajos de $\alpha$ significan que el student casi solo imita al teacher, lo que le permite heredar comportamiento incluso en ejemplos donde la etiqueta dura no es informativa. Valores altos lo anclan a la verdad, útil cuando el teacher está sesgado o se equivoca sistemáticamente en ciertos casos. Un valor típico en la literatura es $\alpha \in [0{,}1, 0{,}5]$, pero depende fuertemente de la tarea y de la calidad del teacher.

\paragraph{Qué está aprendiendo realmente el student.} Puede resultar útil pensar la pérdida KD como un objetivo con dos ``supervisores'' complementarios. Las hard labels le dicen al student \emph{qué} responder; las soft labels del teacher le enseñan, además, \emph{cómo se parece cada posible respuesta a la correcta}. Esa segunda señal es estructural: fija cómo debe organizar el student su espacio de salida, qué errores debe considerar ``cercanos'' a la respuesta correcta y cuáles ``lejanos''. Empíricamente, esta estructura es transferible a ejemplos nuevos (generaliza mejor) y es una de las razones de que KD sea más que un simple \textit{label smoothing}. De hecho, cuando el teacher es de muy alta calidad, distilar con $\alpha$ cercano a cero y $T$ elevado puede superar en generalización al entrenamiento supervisado clásico incluso cuando el student tiene la misma capacidad que el teacher.

\begin{figure}[H]
\centering
\resizebox{\textwidth}{!}{%
\begin{tikzpicture}[
>=Latex,
mdl/.style={draw, rounded corners=2pt, align=center, minimum width=2.6cm, minimum height=9mm, font=\footnotesize, thick},
logit/.style={draw, rounded corners=2pt, align=center, minimum width=2.1cm, minimum height=7mm, font=\footnotesize, fill=gray!10},
loss/.style={draw, rounded corners=2pt, align=center, minimum width=3.2cm, minimum height=8mm, font=\footnotesize, thick},
arr/.style={->, thick}
]
% Entrada
\node[font=\footnotesize] (x1) at (0,0) {$x$};

% Modelos
\node[mdl, fill=bcteacher!40] (teacher) at (2.3, 1.3) {Teacher\\(gran modelo)};
\node[mdl, fill=bcstudent!40] (student) at (2.3, -1.3) {Student\\(modelo pequeño)};

% Logits / probabilidades
\node[logit] (pt) at (5.2, 1.3) {$p^t(x;T)$};
\node[logit] (ps) at (5.2, -1.3) {$p^s(x;T)$};

% Etiqueta dura
\node[logit, fill=bcdata!40] (y) at (5.2, -3.0) {etiqueta $y$};

% Losses parciales
\node[loss, fill=bcrouter!30] (kl) at (9.0, 1.3) {$T^{2}\,\mathrm{KL}(p^{t}\|p^{s})$};
\node[loss, fill=bcdata!25] (ce) at (9.0, -2.3) {$\mathrm{CE}(y,\, p^{s}(x;1))$};

% Loss total
\node[loss, fill=bcmeasure!35, minimum width=4.4cm] (sum) at (13.3, -0.5) {$\mathcal{L}=\alpha\,\mathrm{CE}+(1{-}\alpha)\,T^{2}\,\mathrm{KL}$};

% Flechas de entrada a modelos
\draw[arr] (x1) -- (teacher.west);
\draw[arr] (x1) -- (student.west);

% Modelos a logits
\draw[arr] (teacher.east) -- (pt.west);
\draw[arr] (student.east) -- (ps.west);

% pt -> KL (recta)
\draw[arr] (pt.east) -- (kl.west);
% ps -> KL (sube, pasando a la derecha de pt)
\draw[arr] (ps.east) -- ++(0.6,0) |- (kl.south);

% y -> CE
\draw[arr] (y.east) -- (ce.west);
% ps -> CE (pasa por debajo, dashed indica T=1)
\draw[arr, dashed] (ps.south) -- ++(0,-0.55) -| (ce.north);

% KL y CE -> L
\draw[arr] (kl.east) -| (sum.north);
\draw[arr] (ce.east) -| (sum.south);
\end{tikzpicture}%
}%
\caption{Esquema de la destilación clásica por respuesta: el teacher y el student procesan la misma entrada, y la pérdida del student combina una CE contra la etiqueta dura con una KL contra la distribución suavizada del teacher. La línea discontinua indica que el término CE usa el student evaluado a temperatura 1. Fuente: elaboración propia.}
\label{fig:kd-esquema}
\end{figure}

\subsection{De los modelos clasificadores a los modelos generativos}

La formulación anterior se pensó originalmente para clasificadores: una entrada, una salida categórica. En LLMs generativos la situación es distinta porque la ``salida'' no es una etiqueta única, sino una secuencia de tokens $y = (y_1, y_2, \dots, y_T)$ generada de forma autoregresiva. Esto obliga a repensar qué significa imitar al teacher.

Una opción es aplicar la KD clásica \emph{token a token}: en cada posición $i$ de una secuencia fijada, el teacher produce una distribución $p^t(\cdot \mid y_{<i}, x)$ sobre el siguiente token y el student se entrena para aproximar esa distribución con una KL. Esta variante se conoce como \textbf{word-level KD} o \textbf{on-logits KD} y es la extensión más directa de Hinton al entorno autoregresivo. Tiene la ventaja de transferir todo el detalle probabilístico del teacher, pero presenta inconvenientes prácticos en nuestro contexto: exige mantener el teacher \emph{cargado durante todo el entrenamiento} del student, lo que implica que dos modelos muy grandes deben convivir en memoria de GPU; requiere, además, que el vocabulario del student sea idéntico al del teacher (o al menos compatible mediante una proyección), porque las distribuciones se comparan token a token sobre el mismo vocabulario.

La alternativa, descrita por Kim y Rush~\cite{kim2016sequence}, es la \textbf{destilación a nivel de secuencia} o \textit{sequence-level KD}. En ella, el teacher primero genera respuestas completas $\hat{y}^t(x)$ para cada entrada $x$ del corpus y, después, el student se entrena para \emph{producir} esas mismas respuestas mediante entropía cruzada estándar, sin acceso en línea a las distribuciones del teacher. La pérdida pasa a ser simplemente
\[
\mathcal{L}_{\mathrm{seq-KD}}(x) \;=\; -\sum_{i=1}^{T} \log p^s\!\big(\hat{y}^t_i(x)\,\big|\,\hat{y}^t_{<i}(x),\, x\big),
\]
es decir, un fine-tuning supervisado clásico sobre un corpus sintético generado por el teacher. Conceptualmente, se puede ver como una aproximación de la KD token a token en la que la distribución del teacher se sustituye por una \emph{muestra} suya (la respuesta generada). Se pierde la riqueza de las soft labels (ahora cada posición vuelve a ser una etiqueta \emph{hard} marcada por el teacher), pero se ganan propiedades de ingeniería muy valiosas para trabajar con LLMs grandes.

En primer lugar, \textbf{el teacher se usa una única vez} para generar el corpus sintético y se descarga. El entrenamiento posterior del student es un SFT estándar que cabe en una sola GPU y reutiliza todo el ecosistema habitual de fine-tuning (transformers, TRL, PEFT). En segundo lugar, \textbf{no se requiere compatibilidad de vocabulario ni de arquitectura} entre teacher y student: basta con que ambos produzcan texto. Esto abre la puerta a destilar entre familias de modelos distintas o incluso a destilar salidas de modelos cerrados vía API. En tercer lugar, \textbf{el corpus sintético es reutilizable}: distintos students, distintas configuraciones de hiperparámetros o distintas mezclas de dominios se pueden entrenar sobre el mismo corpus, lo que abarata considerablemente las iteraciones experimentales.

El coste se paga en otro lado: el student aprende a reproducir las respuestas del teacher, no su \emph{distribución}. Si una pregunta tiene varias soluciones válidas, el student solo ``ve'' la que tocó generar el teacher y pierde información sobre alternativas razonables. Empíricamente, la literatura reciente sobre LLMs muestra que esta pérdida es a menudo aceptable cuando el corpus sintético es suficientemente grande y diverso, y que los beneficios prácticos compensan ampliamente el coste. Por ese motivo, la destilación a nivel de secuencia es hoy el esquema dominante en sistemas reales, y es también el que se adopta en este TFG.

\begin{figure}[H]
\centering
\resizebox{\textwidth}{!}{%
\begin{tikzpicture}[
>=Latex,
mdl/.style={draw, rounded corners=3pt, align=center, minimum width=30mm, minimum height=11mm, font=\footnotesize, thick},
data/.style={draw, rounded corners=2pt, align=center, minimum width=28mm, minimum height=9mm, font=\footnotesize, fill=bcdata!30},
arr/.style={->, thick}
]
% Fase 1: Generación del corpus
\node[font=\bfseries\footnotesize] at (0, 2.6) {Fase 1: Generación};
\node[data] (prompts) at (0, 1.5) {Prompts\\(GSM8K, MATH)};
\node[mdl, fill=bcteacher!40] (teacher) at (0, 0) {Teacher\\(7B / 14B)};
\node[data, fill=bcmeasure!25] (corpus) at (0, -1.5) {Corpus sintético\\$\{(x_i, \hat{y}^t_i)\}$};
\draw[arr] (prompts) -- (teacher);
\draw[arr] (teacher) -- (corpus);

% Fase 2: SFT del student
\node[font=\bfseries\footnotesize] at (6.5, 2.6) {Fase 2: SFT del student};
\node[mdl, fill=bcstudent!40] (student) at (6.5, 0) {Student\\(1.5B)};
\node[data, fill=bcrouter!25] (trained) at (6.5, -1.5) {Student\\destilado};

\draw[arr] (corpus.east) -- ++(1.2,0) |- (student.west) node[pos=0.25, above, font=\scriptsize] {entrenamiento SFT};
\draw[arr] (student) -- (trained);

% Anotaciones
\node[font=\scriptsize, align=center, text=gray] at (0, -2.5) {El teacher se descarga\\tras esta fase};
\node[font=\scriptsize, align=center, text=gray] at (6.5, -2.5) {Reutiliza tooling estándar\\(transformers, TRL, PEFT)};
\end{tikzpicture}%
}%
\caption{Pipeline de destilación \textit{sequence-level} adoptado en este TFG. En la primera fase, el teacher genera respuestas completas sobre los prompts de los benchmarks. En la segunda, el student se ajusta vía SFT sobre ese corpus sintético. El teacher no interviene en la segunda fase, lo que permite iterar sobre el student con una sola GPU. Fuente: elaboración propia.}
\label{fig:seq-kd-pipeline}
\end{figure}

\subsection{Otras variantes relevantes}

La taxonomía completa es mucho más rica. La Tabla~\ref{tab:kd-variantes} recoge las variantes que es útil conocer para situar el trabajo, agrupadas por la naturaleza de la señal que se transfiere.

\begin{table}[H]
\centering
\caption{Principales variantes de destilación de conocimiento en modelos neuronales. Fuente: elaboración propia a partir de~\cite{hinton2015distilling, sanh2019distilbert, gou2021survey, agarwal2024gkd}.}
\label{tab:kd-variantes}
\footnotesize
\begin{tabularx}{\textwidth}{>{\raggedright\arraybackslash}p{3cm} >{\raggedright\arraybackslash}X >{\raggedright\arraybackslash}p{3.5cm}}
\toprule
\textbf{Variante} & \textbf{Qué se transfiere} & \textbf{Ejemplo representativo} \\
\midrule
Response-based (\textit{word-level KD}) & Distribución del teacher sobre el vocabulario, token a token & Hinton et al.~\cite{hinton2015distilling}; DistilBERT~\cite{sanh2019distilbert} \\
\addlinespace
Sequence-level KD & Secuencia completa generada por el teacher, usada como etiqueta para un SFT del student & Kim \& Rush~\cite{kim2016sequence}; Alpaca/Vicuna~\cite{taori2023alpaca} \\
\addlinespace
Feature-based & Activaciones intermedias (representaciones ocultas) en capas seleccionadas & FitNets, hidden-state distillation \\
\addlinespace
Relation-based & Relaciones entre ejemplos (distancias o similitudes entre representaciones) & Relational KD, RKD \\
\addlinespace
On-policy / generative KD & El student genera, el teacher puntúa, y la pérdida se aplica sobre las secuencias del propio student & GKD~\cite{agarwal2024gkd} \\
\bottomrule
\end{tabularx}
\end{table}

\textbf{Response-based (word-level).} Es la destilación clásica adaptada al entorno autoregresivo: el student aprende a igualar la distribución del teacher en cada posición. Muy completa en términos de información transferida, pero cara de orquestar cuando teacher y student son ambos grandes, y restringida a pares con vocabulario compatible.

\textbf{Sequence-level.} Es la variante adoptada en este TFG y su presencia en la tabla por derecho propio refleja su importancia práctica en la era de los LLMs. El teacher genera respuestas completas y el student se ajusta vía SFT sobre ellas. Es el esquema que está detrás de la mayoría de modelos derivados que han aparecido en 2023--2024, donde modelos abiertos de unos pocos miles de millones de parámetros se entrenan sobre salidas generadas por modelos más grandes (propietarios o abiertos). En realidad, muchos de los modelos que hoy conocemos bajo la etiqueta de ``\textit{student}'' son sequence-level KD disfrazados.

\textbf{Feature-based.} El student no solo reproduce la salida final sino también activaciones intermedias en capas seleccionadas. Gana expresividad porque transfiere representaciones internas, pero necesita una correspondencia razonable entre las arquitecturas del teacher y del student. Si el student tiene menos capas o dimensiones distintas, hay que añadir proyectores entrenables que alineen los espacios. En contextos donde teacher y student son de la misma familia y difieren solo en tamaño (por ejemplo, en Qwen2.5 de 14B y 7B), feature-based es técnicamente factible pero rara vez aporta ganancias claras por encima de la sequence-level KD.

\textbf{Relation-based.} En lugar de imitar \emph{qué} dice el teacher sobre cada ejemplo, imita \emph{cómo relaciona} el teacher distintos ejemplos entre sí (distancias, similitudes, ángulos en el espacio de representación). Es útil en \textit{representation learning} (por ejemplo, para destilar encoders) pero menos directo en LLMs generativos, donde la salida es texto.

\textbf{On-policy vs off-policy.} Esta distinción es ortogonal a las anteriores y ha ganado peso en los últimos años. En la destilación \textit{off-policy} (también llamada ``por imitación'' o de corpus fijo) el teacher genera previamente y el student se entrena sobre esas secuencias. Es sencilla, reproducible y permite reutilizar el corpus entre experimentos. Pero tiene un problema conocido como \textit{train--test mismatch}: el student, una vez entrenado, genera sus propias secuencias que pueden separarse distribucionalmente de las del teacher; en esas zonas nuevas del espacio, el student nunca fue entrenado y su comportamiento puede degradarse.

En la destilación \textit{on-policy} o \textit{generative} (GKD)~\cite{agarwal2024gkd}, el student genera sus propias secuencias durante el entrenamiento, el teacher las evalúa en tiempo real y la pérdida se aplica sobre esas secuencias generadas. Esto cierra el mismatch: el student aprende precisamente en las zonas donde acaba moviéndose en inferencia. A cambio, es mucho más caro porque requiere tanto \textit{forward} del student (para generar) como \textit{forward} del teacher (para puntuar) en bucle, y exige mantener ambos modelos en memoria simultáneamente. En presupuestos de cómputo ajustados, la estrategia habitual es priorizar el off-policy y, si sobra presupuesto, añadir una fase final de refinamiento on-policy.

\subsection{Cómo encaja la destilación en este TFG}

En este TFG, el papel de la destilación es \textbf{habilitar} la cascada. Su rol no es introducir una técnica nueva de KD, sino obtener uno o más students capaces de atender un porcentaje razonable de peticiones con calidad aceptable, de modo que las políticas B y C tengan un modelo pequeño al que merezca la pena enviar tráfico antes de escalar al grande. Dentro del mapa anterior, el esquema concreto adoptado es una destilación \textbf{sequence-level, off-policy}. Esta elección es consecuencia directa de varias restricciones prácticas del proyecto que vale la pena hacer explícitas.

La primera es el \textbf{presupuesto de GPU acotado}. El corpus de destilación se genera una única vez con el teacher sobre GSM8K y MATH, y el teacher se descarga inmediatamente después. Todos los entrenamientos posteriores del student se realizan con un único modelo en GPU, lo que permite iterar con rapidez sobre configuraciones distintas sin consumir cuota de forma desproporcionada.

La segunda es la \textbf{compatibilidad de \textit{tooling}} y la simplicidad operativa. Una vez que se dispone de un corpus de pares \emph{(prompt, respuesta del teacher)}, el entrenamiento del student es un fine-tuning supervisado estándar. Esto permite aprovechar el ecosistema maduro de herramientas de SFT (transformers, TRL, PEFT) sin necesidad de infraestructura \textit{ad hoc}, y combinar la destilación con LoRA (Sec.~\ref{sec:finetuning-lora}) en las ramas de iteración rápida y con fine-tuning completo en la rama final más ambiciosa.

La tercera es la \textbf{cobertura del dominio}. Los benchmarks objetivo son matemáticos, donde la correctitud final es binaria y el razonamiento intermedio suele ser reproducible por un student si el corpus de destilación cubre los tipos de problema con suficiente diversidad. En este régimen, la literatura reciente sugiere que la pérdida de información al pasar de \emph{word-level} a \emph{sequence-level} KD es pequeña y claramente aceptable frente al ahorro de cómputo.

La cuarta es la \textbf{reusabilidad del corpus}. El mismo corpus sintético se ha usado como base para múltiples experimentos de destilación del TFG (distintas configuraciones de hiperparámetros, LoRA frente a FT completo, filtrado por corrección de respuestas, entre otros), lo que es una propiedad exclusiva del esquema off-policy.

\paragraph{Plan de contingencia.} Si el protocolo de decisión T5.3 hubiera concluido que la destilación no era viable con la cuota disponible, se habría adoptado el plan de contingencia R2 con un modelo preentrenado de la misma familia (\texttt{Qwen2.5-1.5B-Instruct} directamente, sin ajuste adicional), aceptando una cota de calidad algo más baja como precio por no tener que entrenar. Esta decisión queda explícitamente documentada y no invalida el resto del experimento, porque las políticas de cascada se pueden evaluar contra cualquier student, destilado o no. En la práctica el proyecto ha llegado a ejecutar los experimentos de destilación y, por tanto, la configuración de referencia que se reporta en la parte experimental es la del student destilado.

%%%%%%%%%%%%%%%%%%%%%%%%%%%%%%%%%%%%%%%%%%%%%%%%%%%%%%%%%%%%%%%%%
\section{Serving, cascadas de modelos e infraestructura HPC}
\label{sec:serving-teorico}

Tener un teacher y un student entrenados no basta para que la cascada funcione en la práctica. Hace falta un motor de inferencia capaz de servir ambos modelos con buena utilización de la GPU, baja latencia de primer token y throughput alto, y con una API estable sobre la que construir el router. Este bloque enlaza tres piezas que en el índice conviene tratar unidas: el \textbf{serving} moderno con vLLM~\cite{vllm2023, vllm-docs}, las \textbf{cascadas y políticas de routing} entre modelos, y el \textbf{entorno HPC} (MareNostrum~5, SLURM) donde se ejecuta el pipeline experimental.

\subsection{Serving con vLLM (memoria, paginación y batching)}

\subsubsection{El problema de memoria y la fragmentación del KV cache}

Como se ha visto en la sección anterior, el KV cache crece linealmente con el número de tokens y ocupa buena parte de la VRAM disponible. El problema se agrava cuando se sirven varias peticiones a la vez, porque cada una tiene su propio KV cache de tamaño variable y dinámico: algunas acabarán generando 50 tokens y otras 2\,000, y no se sabe a priori cuáles.

Los sistemas de serving anteriores a vLLM reservaban el KV cache como un \textbf{bloque contiguo} por petición, dimensionado al peor caso (la longitud máxima esperada). Esto produce dos formas de desperdicio. La \textbf{fragmentación interna} aparece porque cada petición reserva mucho más espacio del que acaba usando: una petición que termina con 300 tokens pero reservó espacio para 2\,048 desperdicia el 85\,\% de esa memoria. La \textbf{fragmentación externa} ocurre porque entre peticiones quedan huecos de memoria contigua que no son suficientemente grandes para una petición nueva pero sí para varias pequeñas, y el gestor clásico no es capaz de aprovecharlos.

El resultado es que, con gestión ingenua, una GPU con 80~GB efectivamente solo puede sostener una fracción reducida de ese espacio como KV cache útil, y el motor de serving acaba limitado no por cómputo sino por memoria disponible para sostener el batch. La consecuencia directa es un throughput bajo.

\subsubsection{PagedAttention y vLLM}

vLLM introduce \textbf{PagedAttention}~\cite{vllm2023}, una idea inspirada directamente en cómo los sistemas operativos gestionan la memoria virtual. La memoria del KV cache se divide en \textbf{bloques} de tamaño fijo (típicamente 16 tokens por bloque). Cada petición mantiene una tabla lógica que mapea sus tokens a bloques físicos no necesariamente contiguos. Cuando una petición crece, el motor le asigna nuevos bloques del pool libre, vengan de donde vengan. Cuando una petición termina, sus bloques vuelven al pool y pueden ser reutilizados por cualquier otra.

Los beneficios de este diseño son inmediatos. En primer lugar, se elimina la fragmentación interna porque solo se asignan bloques cuando son necesarios y no hay ``reserva al peor caso''. En segundo lugar, desaparece la fragmentación externa porque no se requiere memoria contigua: cualquier bloque libre del pool sirve para cualquier petición. En tercer lugar, PagedAttention habilita la compartición de prefijos: cuando varias peticiones comparten prompt (por ejemplo, el mismo \textit{system prompt}), los bloques correspondientes al prefijo común pueden reutilizarse sin duplicarse, empleando un mecanismo de \textit{copy-on-write} si alguna petición diverge.

El trabajo original~\cite{vllm2023} reportó mejoras de throughput de 2--4$\times$ respecto a motores de serving anteriores en los mismos escenarios, manteniendo la latencia. A día de hoy vLLM es una de las piezas más usadas en serving de LLMs abiertos en producción y es el motor de referencia en el ecosistema del BSC para LLMs.

\begin{figure}[H]
\centering
\resizebox{0.95\textwidth}{!}{%
\begin{tikzpicture}[
>=Latex,
legacyblk/.style={draw, fill=red!25, minimum height=6mm, minimum width=6mm, inner sep=0},
freeblk/.style={draw, fill=gray!20, minimum height=6mm, minimum width=6mm, inner sep=0},
paged/.style={draw, fill=blue!30, minimum height=6mm, minimum width=6mm, inner sep=0},
lbl/.style={font=\scriptsize},
rowlbl/.style={font=\footnotesize\bfseries}
]
% Fila 1: Legacy (reserva contigua con huecos)
\node[rowlbl, anchor=east] at (-0.3, 0) {Legacy};
\foreach \x/\s in {0/legacyblk, 1/legacyblk, 2/legacyblk, 3/freeblk, 4/freeblk, 5/legacyblk, 6/legacyblk, 7/freeblk, 8/freeblk, 9/freeblk, 10/legacyblk, 11/legacyblk}{
  \node[\s] at (\x*0.72, 0) {};
}
\node[lbl, anchor=west] at (11*0.72+0.55, 0) {$\sim$40\% útil};

% Fila 2: PagedAttention (todos los bloques aprovechables)
\node[rowlbl, anchor=east] at (-0.3, -1.1) {PagedAttn};
\foreach \x in {0,1,2,3,4,5,6,7,8,9,10,11}{
  \node[paged] at (\x*0.72, -1.1) {};
}
\node[lbl, anchor=west] at (11*0.72+0.55, -1.1) {$\sim$90\% útil};

% Leyenda de colores
\node[legacyblk] at (0.5, -2.3) {};
\node[lbl, anchor=west] at (0.9, -2.3) {ocupado};
\node[freeblk] at (3.0, -2.3) {};
\node[lbl, anchor=west] at (3.4, -2.3) {libre (fragmento desperdiciado)};
\node[paged] at (7.6, -2.3) {};
\node[lbl, anchor=west] at (8.0, -2.3) {bloque paginado reutilizable};
\end{tikzpicture}%
}%
\caption{Gestión del KV cache: esquema clásico (reserva contigua con huecos por fragmentación) frente a PagedAttention (bloques fijos gestionados como memoria virtual). Cada cuadrado representa un fragmento de VRAM del mismo tamaño. Fuente: elaboración propia a partir de~\cite{vllm2023}.}
\label{fig:paged-attention}
\end{figure}

\subsubsection{Batching continuo}

La otra pieza clave del serving moderno es el \textbf{continuous batching} (también llamado \textit{in-flight batching}). Los motores tradicionales procesan peticiones en batches estáticos: toman un grupo de peticiones, les hacen el prefill, generan tokens hasta que todas terminan, y entonces pasan al siguiente batch. Este esquema tiene un problema obvio: una petición corta queda bloqueada esperando a que una petición larga del mismo batch termine, desperdiciando recursos.

En continuous batching, la unidad mínima de trabajo no es ``una petición completa'' sino ``un paso de decode''. En cada iteración, el motor elige qué peticiones vivas tienen un token listo para procesar, las agrupa en un mini-batch y ejecuta un único paso forward para todas. Cuando una petición termina, se elimina sin afectar al resto; cuando llega una nueva, se le hace prefill y se une al pool sin esperar a un batch nuevo. El efecto combinado es que la GPU se mantiene ocupada casi siempre, y las peticiones cortas no pagan la latencia de las largas.

vLLM implementa continuous batching de manera nativa sobre PagedAttention: cada iteración de decode selecciona peticiones vivas, prepara los bloques correspondientes del KV cache y dispatches un kernel CUDA que las procesa en paralelo. Este bucle es lo que verás referido como el \textit{scheduler} en la documentación de vLLM.

\subsubsection{Otras técnicas relacionadas}

Aparte de PagedAttention y continuous batching, hay dos familias de técnicas que conviene mencionar porque son primas cercanas y aparecen frecuentemente en la literatura de serving eficiente.

La primera es \textbf{FlashAttention}~\cite{dao2022flashattention, dao2023flash2}, que optimiza el kernel de atención para reducir los accesos a HBM mediante \textit{tiling} y cómputo en SRAM del chip. FlashAttention es ortogonal a PagedAttention: mientras que PagedAttention ataca la organización del KV cache en memoria, FlashAttention acelera el cálculo del kernel de atención en sí. Las versiones modernas de vLLM lo integran de serie, y ambos beneficios se suman sin interferencias.

La segunda es el \textbf{speculative decoding}~\cite{leviathan2023speculative}, una técnica en la que un modelo ``borrador'' pequeño propone varios tokens por adelantado y el modelo grande los verifica en paralelo. Si los acepta, se gana velocidad porque se generan varios tokens en un solo paso de verificación; si rechaza alguno, se retrocede al punto de divergencia. Es otra forma de reducir la latencia de decode aprovechando hardware infrautilizado durante los pasos de atención del modelo grande. Guarda un parentesco conceptual con la cascada del TFG (dos modelos de distinto tamaño colaboran para servir una petición), pero opera a un nivel distinto, dentro de una misma petición, token a token, y con el modelo grande siempre implicado.

Estas técnicas no son el foco del TFG, pero conviene tenerlas presentes para situar la cascada dentro del panorama más amplio de serving eficiente. En todas ellas, la filosofía es la misma: reducir el coste de servicio sin cambiar significativamente qué se responde.

\subsection{Cascadas, routing y políticas de decisión}

Las partes anteriores de este capítulo han visto cómo entrenar un student que imita al teacher y cómo servir ambos modelos de forma eficiente. Queda la pregunta central del TFG: \textbf{¿cómo decidir, para cada petición, qué modelo debe responder?} Lo que sigue repasa el marco conceptual de cascadas y routing de modelos, y sitúa las tres políticas del TFG (A, B y C) dentro de ese mapa.

\subsubsection{Motivación. No todas las peticiones necesitan el modelo grande}

La observación empírica que da sentido a la cascada es sencilla y muy robusta: en la mayoría de distribuciones de peticiones reales, \textit{la dificultad se concentra en una cola}. Una gran parte de las peticiones son fáciles (preguntas sencillas, aritmética directa, reformulaciones de texto corto) y un modelo pequeño las resuelve correctamente con costes muy inferiores a los del modelo grande. Solo una minoría requiere razonamiento largo, conocimiento de cola o manejo de casos fuera de distribución.

Esta idea no es nueva. En visión por computador o búsqueda de documentos, los \textit{cascade classifiers} llevan décadas aplicándola~\cite{viola2001cascade}. La novedad en LLMs es que el coste por petición es mucho mayor y la diferencia de coste entre modelos de distintos tamaños es enorme: pasar de 14B a 7B puede reducir el coste por token a la mitad o más. Incluso si la política de routing solo acierta en una parte de las peticiones, el ahorro agregado es sustancial.

\subsubsection{Tres familias de políticas}

Aunque en la literatura aparecen muchos esquemas concretos, todos se pueden agrupar en tres familias según \textit{cuándo} deciden qué modelo sirve cada petición.

\textbf{Política A (baseline).} Siempre se usa el teacher. No hay decisión. Sirve como línea base de calidad máxima y coste máximo. Todo el resto se mide contra esta referencia.

\textbf{Política B (cascada forzada).} Se responde primero con el student. Sobre su respuesta se aplica un \textit{quality gate}, es decir, una función barata que decide si la respuesta parece correcta. Si lo es, se acepta; si no, se escala al teacher y se responde con la suya. La ventaja es que no requiere entrenar ningún componente adicional; la desventaja es que las peticiones que acaban escalando pagan tanto el coste del student como el del teacher. Se suele llamar también \textit{inference cascade} o \textit{sequential cascade}~\cite{chen2023frugalgpt}.

\textbf{Política C (cascada por confianza, routing predictivo).} Antes de servir, un componente predictivo mira la petición y estima a qué modelo conviene enviarla. Si predice que el student la resuelve bien, solo se usa el student; si predice que no, se envía directamente al teacher. Se evita así el coste extra del fallback de la Política B, pero a cambio hay que entrenar el predictor y pagar su latencia (aunque sea pequeña). Este es el esquema que usan sistemas como FrugalGPT~\cite{chen2023frugalgpt} y RouteLLM~\cite{dekoninck2024routing}.

\begin{figure}[H]
\centering
\resizebox{\textwidth}{!}{%
\begin{tikzpicture}[
>=Latex,
node distance=6mm and 8mm,
mdl/.style={draw, rounded corners=2pt, align=center, minimum width=22mm, minimum height=9mm, font=\footnotesize, thick},
dec/.style={draw, rounded corners=2pt, align=center, minimum width=18mm, minimum height=9mm, font=\footnotesize, fill=bcrouter!40},
req/.style={draw, rounded corners=2pt, align=center, minimum width=16mm, minimum height=8mm, font=\footnotesize, fill=yellow!30},
out/.style={draw, rounded corners=2pt, align=center, minimum width=16mm, minimum height=8mm, font=\footnotesize, fill=bcdata!35}
]
% Políticas en columnas
% --- Col A ---
\node[font=\bfseries\footnotesize] at (0,2.3) {Política A};
\node[req] (ra) at (0,1.2) {petición};
\node[mdl, fill=bcteacher!40] (ta) at (0,0) {Teacher};
\node[out] (oa) at (0,-1.2) {respuesta};
\draw[->] (ra) -- (ta);
\draw[->] (ta) -- (oa);

% --- Col B ---
\node[font=\bfseries\footnotesize] at (4.8,2.3) {Política B};
\node[req] (rb) at (4.8,1.2) {petición};
\node[mdl, fill=bcstudent!40] (sb) at (4.8,0) {Student};
\node[dec, minimum width=22mm] (gateb) at (4.8,-1.1) {quality gate};
\node[mdl, fill=bcteacher!40] (tb) at (4.8,-2.4) {Teacher};
\node[out] (ob) at (4.8,-3.6) {respuesta};
\draw[->] (rb) -- (sb);
\draw[->] (sb) -- (gateb);
\draw[->] (gateb) -- node[right, font=\scriptsize] {fallback} (tb);
\draw[->] (tb) -- (ob);
\draw[->] (gateb.west) -- ++(-7mm,0) |- node[near start, left, font=\scriptsize] {acepta} (ob.west);

% --- Col C ---
\node[font=\bfseries\footnotesize] at (10,2.3) {Política C};
\node[req] (rc) at (10,1.2) {petición};
\node[dec, minimum width=26mm] (pred) at (10,0) {predictor / router};
\node[mdl, fill=bcstudent!40] (sc) at (8.6,-1.4) {Student};
\node[mdl, fill=bcteacher!40] (tc) at (11.4,-1.4) {Teacher};
\node[out] (oc) at (10,-2.8) {respuesta};
\draw[->] (rc) -- (pred);
\draw[->] (pred.west) -- ++(-3mm,0) |- (sc.north) node[midway, left, font=\scriptsize] {fácil};
\draw[->] (pred.east) -- ++(3mm,0) |- (tc.north) node[midway, right, font=\scriptsize] {difícil};
\draw[->] (sc.south) |- (oc.west);
\draw[->] (tc.south) |- (oc.east);
\end{tikzpicture}%
}
\caption{Las tres políticas de cascada comparadas en este TFG. La A siempre usa teacher, la B usa student y escala a teacher solo si el quality gate rechaza, la C decide a priori con un router. Fuente: elaboración propia.}
\label{fig:politicas}
\end{figure}

\subsubsection{Señales usadas por un router}

La Política C requiere un componente que tome la decisión sin ejecutar el modelo grande. Las señales que usa el router pueden venir de tres fuentes distintas. La primera son \textbf{features del prompt}: propiedades que se pueden calcular sin ejecutar ningún modelo, como la longitud, la presencia de números, la densidad de palabras técnicas, el idioma o la similitud con prompts anteriores cacheados. Son baratísimas de calcular pero informan poco sobre la dificultad real de la petición.

La segunda fuente son \textbf{features del student}, que aparecen al hacer inferencia rápida con el modelo pequeño: log-probabilidades de los primeros tokens generados, entropía de la distribución de salida, perplejidad del propio prompt o margen entre el top-1 y el top-2 del último token. Estas señales son mucho más informativas que las anteriores, a costa de tener que pasar la petición por el student para obtenerlas.

La tercera opción son \textbf{features mixtas}, donde un predictor ligero (regresión logística, MLP o \textit{gradient boosting}) combina las señales anteriores para aprender un clasificador de dificultad. Este es el esquema más común cuando se busca un compromiso razonable entre coste del router y capacidad predictiva.

FrugalGPT~\cite{chen2023frugalgpt} propone una estrategia en dos fases: primero una cascada forzada (Política B en nuestros términos), después un router entrenado con las señales acumuladas, y finalmente un \textit{LLM scorer} que decide si aceptar cada respuesta. RouteLLM~\cite{dekoninck2024routing} va más lejos y entrena el router con pares de preferencias (qué modelo era mejor en cada ejemplo) usando aprendizaje preferencial. En ambos casos, la idea común es que el router no necesita ser muy complicado para ser útil: incluso clasificadores lineales sobre señales simples consiguen reducir significativamente el uso del modelo grande.

\subsubsection{Métricas para comparar políticas}

Comparar políticas de cascada no es tan directo como comparar la accuracy de dos modelos. Hay que reportar, como mínimo, tres dimensiones a la vez. La primera es la \textbf{calidad}, medida como accuracy (o \textit{exact match}) en el benchmark de referencia. La segunda es el \textbf{coste}: en el caso de LLMs, la métrica más portable es la fracción de peticiones que acabaron pasando por el teacher, aunque también se suele reportar el coste promedio por petición en términos de tokens generados o GPU-segundo. La tercera es la \textbf{latencia}, especialmente el TTFT y el percentil 95 de la latencia total, porque un router que reduce mucho el uso del teacher pero añade 500~ms de overhead puede no ser útil en una aplicación interactiva.

En el TFG se reportan las tres dimensiones simultáneamente, y se incluye además una comparativa relativa: cuánto porcentaje de calidad se pierde a cambio de cuánto porcentaje de coste se ahorra. Esta es la forma más honesta de caracterizar el compromiso.

\subsubsection{Dónde se sitúa este TFG}

El TFG no pretende inventar un nuevo algoritmo de routing, sino \textbf{caracterizar experimentalmente} las tres políticas bajo el mismo protocolo y en un entorno HPC real. La contribución no es teórica sino metodológica y cuantitativa. Se define un protocolo de evaluación común, reproducible, con warmup, ventanas de medida y repeticiones controladas. Se miden simultáneamente calidad y coste para las tres políticas, y se discute para qué franjas de calidad exigida la cascada es preferible al teacher directo y para cuáles no.

Dentro del mapa del estado del arte, la Política B es un método conocido y bien documentado; la Política C es la variante de routing tipo FrugalGPT/RouteLLM simplificada para un escenario con dos modelos (teacher y student) en lugar de un pool grande. La ambición no es ganar un benchmark sino aportar un estudio honesto y reproducible en MareNostrum~5.

%%%%%%%%%%%%%%%%%%%%%%%%%%%%%%%%%%%%%%%%%%%%%%%%%%%%%%%%%%%%%%%%%
\section{Estado del arte}
\label{sec:estado-del-arte}

Una vez presentados los fundamentos teóricos, conviene situar el trabajo dentro del panorama de investigación reciente. Esta sección revisa las contribuciones más relevantes en las tres líneas que convergen en el TFG (destilación de conocimiento en modelos generativos, serving eficiente de LLMs y routing o cascada entre modelos), con el doble objetivo de acreditar la pertinencia de la propuesta y de delimitar con claridad qué aporta este trabajo respecto a lo que ya existe.

\subsection{Destilación de conocimiento en modelos generativos}

La destilación de conocimiento, formulada originalmente por Hinton et al.~\cite{hinton2015distilling} para modelos de clasificación, ha experimentado una evolución acelerada desde que los LLMs se convirtieron en el paradigma dominante del procesamiento de lenguaje natural. La adaptación al contexto generativo no fue directa: en un clasificador, el teacher produce una distribución sobre un número fijo de clases; en un modelo autoregresivo, la distribución es sobre un vocabulario de decenas de miles de tokens y cambia en cada posición de la secuencia.

Los primeros trabajos en destilar modelos de lenguaje a escala se concentraron en encoders. DistilBERT~\cite{sanh2019distilbert} demostró que un modelo BERT con la mitad de capas, entrenado con una combinación de pérdida de distilación, supervisión clásica y alineamiento de representaciones, conservaba el 97\,\% del rendimiento del teacher en benchmarks de comprensión. Este resultado fue influyente porque estableció un protocolo práctico de destilación que se ha replicado en decenas de variantes posteriores.

La transición a modelos generativos trajo consigo la distinción, ya expuesta en la sección de Knowledge Distillation, entre \emph{word-level KD} y \emph{sequence-level KD}. Kim y Rush~\cite{kim2016sequence} formalizaron esta última para traducción automática y mostraron que entrenar un student con las traducciones del teacher como etiquetas era competitivo con la KD por logits y mucho más sencillo de implementar. Este enfoque fue el precursor directo de la oleada de modelos destilados que apareció en 2023--2024: Alpaca~\cite{taori2023alpaca} (student de LLaMA entrenado sobre salidas de GPT-3.5), Vicuna (sobre conversaciones de ShareGPT), Orca (que además incorpora cadenas de razonamiento del teacher), y la familia Phi de Microsoft (que combina destilación con curación cuidadosa de datos sintéticos).

En paralelo, Agarwal et al.~\cite{agarwal2024gkd} propusieron \textit{Generalized Knowledge Distillation} (GKD), un esquema on-policy donde el student genera y el teacher evalúa. GKD aborda el problema de \textit{train--test mismatch} propio de la destilación off-policy y reporta mejoras consistentes en tareas de generación abierta, aunque a costa de un presupuesto de cómputo significativamente mayor. La tensión entre off-policy (barato, reproducible, reutilizable) y on-policy (más fiel, más caro) es un eje central en la literatura actual y ha condicionado directamente la elección de este TFG, que opta por sequence-level off-policy por razones de presupuesto.

Un resultado reciente especialmente relevante es el de la familia Qwen 2.5~\cite{qwen2}, que incluye modelos desde 0,5B hasta 72B parámetros dentro de la misma arquitectura. Los modelos pequeños de esta familia se benefician de destilación durante su pipeline de entrenamiento, lo que los convierte en puntos de partida ya parcialmente destilados. Este hecho es importante para el TFG: el student de 1,5B que se usa como base no parte de cero en términos de capacidad, sino de un modelo que ya ha sido expuesto a señales de modelos mayores durante su preentrenamiento.

\subsection{Serving eficiente de modelos de lenguaje}

El serving de LLMs ha sido un área de innovación rápida impulsada por la necesidad industrial de reducir el coste por token generado. PagedAttention~\cite{vllm2023}, incorporado en vLLM, fue el primer salto cualitativo al aplicar la idea de paginación de memoria virtual al KV cache, multiplicando el throughput sin alterar la calidad. Desde entonces, el ecosistema se ha diversificado con motores como TensorRT-LLM (NVIDIA), SGLang, y TGI (HuggingFace), cada uno con optimizaciones especializadas.

FlashAttention~\cite{dao2022flashattention, dao2023flash2} atacó el cuello de botella de la atención por otro flanco: reorganizando el cómputo para minimizar los accesos a HBM mediante \textit{tiling} y fusión de operaciones. La combinación de FlashAttention (kernel rápido) con PagedAttention (gestión eficiente de memoria) es hoy el estándar de facto en serving de LLMs y es exactamente la configuración que usa el TFG a través de vLLM.

El \textit{speculative decoding}~\cite{leviathan2023speculative} introdujo un paradigma distinto, en el que un modelo borrador propone tokens que el modelo principal verifica en paralelo. Aunque guarda un parentesco conceptual con la cascada del TFG (dos modelos de distinto tamaño colaboran), opera a una escala diferente (dentro de una misma petición, token a token) y con el modelo grande siempre involucrado. No obstante, es técnicamente compatible con el enfoque de cascada, de modo que un student destilado podría servir simultáneamente como modelo borrador para speculative decoding cuando la petición se enruta al teacher.

Más recientemente, la investigación se ha orientado hacia la desagregación de prefill y decode en GPUs separadas, la cuantización dinámica del KV cache, y el uso de caches jerárquicas que persisten entre peticiones. Estas líneas son complementarias a la cascada y podrían integrarse en futuras iteraciones del sistema.

\subsection{Routing y cascadas de modelos}

La idea de asignar dinámicamente peticiones a modelos de distinto coste es anterior a los LLMs. Los \textit{cascade classifiers} de Viola y Jones~\cite{viola2001cascade} ya aplicaban este principio en detección de caras, con un clasificador rápido filtrando la mayoría de ventanas negativas antes de invocar al costoso. En el contexto de LLMs, FrugalGPT~\cite{chen2023frugalgpt} fue uno de los primeros trabajos en proponer explícitamente una cascada de modelos de lenguaje, combinando un modelo barato con un \textit{quality scorer} y un modelo caro, y demostrando ahorros de coste de hasta un 98\,\% con pérdidas de calidad inferiores al 1\,\% en ciertos benchmarks.

RouteLLM~\cite{dekoninck2024routing} formalizó la idea con un router entrenado mediante aprendizaje preferencial sobre pares de respuestas, y demostró que incluso un clasificador simple puede desviar una fracción significativa del tráfico al modelo pequeño sin pérdida apreciable de calidad percibida. El trabajo de Ong et al.~\cite{ong2025routellm} propuso un marco unificado que engloba tanto el routing (decisión ex-ante sobre a qué modelo enviar) como el cascading (escalado secuencial después de una primera respuesta), y definió una función de utilidad
\[
u(x) \;=\; q(x) \;-\; \lambda\, c(x),
\]
donde $q(x)$ es la calidad esperada, $c(x)$ el coste, y $\lambda$ un parámetro que controla el compromiso. Esta formulación es la base conceptual de la Política C del TFG.

La propuesta de este TFG se diferencia de los trabajos anteriores en un aspecto específico, la redefinición del coste $c(x)$ como \textbf{coste de servicio}, incorporando el estado dinámico del sistema (ocupación de GPUs, tamaño de batch, longitud de cola, KV cache residual) como variable del problema. Mientras que FrugalGPT y RouteLLM definen $c(x)$ como un coste estático asociado al modelo (precio por API call, o latencia media), la extensión propuesta redefine la utilidad como $u(x,z) = q(x) - \lambda\,c(x,z)$, donde $z$ captura el estado del sistema de serving en el instante de la decisión. Esta formulación busca demostrar que un router consciente del estado del sistema puede mantener la misma calidad y latencia con mejor throughput útil y menor consumo energético.

\subsection{Posicionamiento del TFG}

Con este panorama, el TFG se sitúa en la intersección de las tres líneas: utiliza la destilación de conocimiento para producir un student eficiente, despliega teacher y student sobre un motor de serving moderno en un entorno HPC real, y evalúa políticas de cascada que combinan ambos modelos. La contribución no pretende ser teórica en ninguna de las tres áreas, sino metodológica e integradora: construir un pipeline completo, reproducible y medible que conecte las tres piezas en MareNostrum~5, y caracterizar experimentalmente el compromiso calidad-coste-servicio bajo condiciones realistas. La extensión del coste a coste de servicio, sensible al estado del sistema, constituye la novedad conceptual más relevante del proyecto.

\subsection{Entorno HPC en MareNostrum~5 y SLURM}

La última pieza del marco teórico es el entorno donde todo esto se ejecuta. Este apartado describe, con el nivel de detalle necesario para entender el diseño del pipeline, cómo es MareNostrum~5 y qué implicaciones tiene sobre el proyecto.

\subsubsection{MareNostrum~5 a vista de pájaro}

MareNostrum~5 es el supercomputador del Barcelona Supercomputing Center (BSC), uno de los centros de referencia de computación de altas prestaciones en Europa~\cite{mn5}. Incluye particiones generales de CPU y particiones de aceleradores GPU. La relevante para este TFG es la partición de GPUs NVIDIA Hopper (H100): cada nodo combina CPUs de altas prestaciones con múltiples GPUs H100 de 64~GB de VRAM conectadas por NVLink, y a nivel de clúster los nodos se comunican por una red de alta velocidad tipo InfiniBand.

Desde el punto de vista del usuario, el clúster es un sistema Linux al que se entra por SSH y donde no se ejecutan comandos directamente en las GPUs: se envían \textit{jobs} al gestor de colas. Cuando el gestor encuentra recursos libres, lanza el job, lo ejecuta y devuelve los resultados. Esto introduce una latencia entre ``quiero ejecutar algo'' y ``se está ejecutando'' que puede ir desde segundos hasta horas, según la ocupación del clúster.

\subsubsection{SLURM como gestor de colas}

El gestor de colas es \textbf{SLURM}~\cite{slurm}. SLURM es el estándar de facto en la mayoría de supercomputadores académicos y ofrece todo lo esencial: solicitar recursos (nodos, CPUs, GPUs, memoria, tiempo), controlar la cola, monitorizar jobs, limitar el uso por usuario o grupo y registrar consumo de recursos (incluido \textit{energía} vía \texttt{sacct}). Sus comandos más habituales son \texttt{sbatch} para enviar un job, \texttt{squeue} para ver el estado de la cola, \texttt{scancel} para cancelar, \texttt{sacct} para ver el histórico y \texttt{salloc}/\texttt{srun} para sesiones interactivas.

Cada job se describe con un script de bash anotado con directivas \texttt{\#SBATCH}. Para este TFG, los scripts de experimentación solicitan una GPU H100 completa, una cantidad fija de CPUs y un tiempo límite ajustado al tipo de experimento (desde minutos para un benchmark hasta horas para una destilación). La política de la cola se rige por la cuota asignada al grupo de investigación Data Centric Computing.

\subsubsection{Implicaciones para el proyecto}

El entorno HPC no es solo ``ordenadores más grandes''; cambia cómo hay que diseñar los experimentos. Para este TFG, las implicaciones más importantes son las siguientes.

La primera es que \textbf{todo experimento es un job}. No se puede mantener una sesión interactiva abierta con GPU durante 19 semanas. Cada experimento se describe como un script SLURM reproducible, y cada ejecución queda identificada por su JobID. Esto obliga a tener automatización desde el primer día, pero a cambio proporciona trazabilidad completa de cada ejecución.

La segunda es que \textbf{la ocupación de la cola introduce varianza temporal}. El mismo job puede arrancar a los 2 minutos o a las 6 horas, dependiendo de la carga del sistema. Esto afecta tanto a la planificación como a las comparaciones. Si se quieren comparar dos configuraciones, es necesario asegurarse de que se ejecutan en condiciones equivalentes (mismo tipo de nodo, misma configuración de red, idealmente dentro de una ventana temporal cercana).

La tercera es que \textbf{la infraestructura de red del clúster impone sus propias restricciones}. Servir un LLM con una API HTTP dentro de un nodo del clúster es posible, pero el acceso desde fuera del nodo o desde fuera del clúster requiere túneles SSH o configuraciones específicas. Por ello, el pipeline de evaluación se diseña para que el cliente se ejecute dentro del mismo nodo que el servidor, evitando atravesar cualquier frontera de red.

La cuarta es que \textbf{la energía se puede medir}. SLURM registra \texttt{ConsumedEnergy} por job, lo que permite comparar el coste energético real de las distintas políticas de cascada. En la sección de sostenibilidad ya se avanzó que este valor se registra y se publica junto con los resultados experimentales.

Con esto se cierra el marco teórico. Quedan establecidas todas las piezas necesarias, entre ellas los fundamentos de redes neuronales y modelos de lenguaje, la arquitectura transformer y el coste de su inferencia, el \textit{fine-tuning} y el PEFT, la destilación de conocimiento, el serving moderno con vLLM, las cascadas y el routing entre modelos, y el entorno HPC donde se ejecuta todo. En la Parte~\ref{part:diseno-sistema} se utiliza esta base para describir el diseño del sistema.

%%%%%%%%%%%%%%%%%%%%%%%%%%%%%%%%%%%%%%%%%%%%%%%%%%%%%%%%%%%%%%%%%
\part{Diseño del sistema}
\label{part:diseno-sistema}
%%%%%%%%%%%%%%%%%%%%%%%%%%%%%%%%%%%%%%%%%%%%%%%%%%%%%%%%%%%%%%%%%

\section{Diseño del sistema experimental}
\label{sec:diseno-sistema-exp}

\subsection{Visión global del pipeline}

El sistema se organiza en cuatro pipelines concatenados: \textit{destilación}, \textit{serving}, \textit{evaluación} y \textit{routing}. Los dos primeros producen activos (un student destilado y un par de servidores de inferencia); el tercero somete esos activos a un protocolo de medida uniforme; el cuarto combina los servidores bajo una política concreta (A, B o C) y se vuelve a medir con el mismo protocolo del tercero. La Figura~\ref{fig:pipeline-global} muestra la vista integrada.

\begin{figure}[H]
\centering
\resizebox{\textwidth}{!}{%
\begin{tikzpicture}[
>=Latex,
node distance=6mm and 10mm,
stage/.style={draw, rounded corners=3pt, align=center, minimum width=38mm, minimum height=12mm, font=\footnotesize, thick},
asset/.style={draw, rounded corners=2pt, align=center, minimum width=26mm, minimum height=8mm, font=\scriptsize, fill=gray!12}
]
% Pipeline destilación
\node[stage, fill=bcdata!25] (p1) {\textbf{1. Destilación}\\teacher $\to$ student};
\node[asset, below=4mm of p1] (a1) {checkpoint student};

% Pipeline serving
\node[stage, right=14mm of p1, fill=bcstudent!25] (p2) {\textbf{2. Serving}\\vLLM (teacher + student)};
\node[asset, below=4mm of p2] (a2) {dos endpoints HTTP};

% Pipeline evaluación
\node[stage, right=14mm of p2, fill=bcmeasure!25] (p3) {\textbf{3. Evaluación}\\calidad + servicio};
\node[asset, below=4mm of p3] (a3) {métricas baseline};

% Pipeline routing
\node[stage, right=14mm of p3, fill=bcrouter!35] (p4) {\textbf{4. Routing}\\políticas A / B / C};
\node[asset, below=4mm of p4] (a4) {comparativa calidad-coste};

% flechas
\draw[->] (p1) -- (p2);
\draw[->] (p2) -- (p3);
\draw[->] (p3) -- (p4);
\draw[->, dashed] (a1) -- (p2);
\draw[->, dashed] (a2) -- (p3);
\draw[->, dashed] (a3) -- (p4);
\end{tikzpicture}%
}
\caption{Vista global del pipeline: cuatro bloques encadenados que producen activos reutilizables y métricas comparables entre sí bajo un único protocolo de medida. Fuente: elaboración propia.}
\label{fig:pipeline-global}
\end{figure}

Lo importante de esta vista es que \textbf{todas las mediciones comparten el mismo protocolo}, ya descrito en la Parte~\ref{part:gestion-tesis} (Sección~\ref{sec:metodologia-trabajo}, con warmup, ventana, repeticiones y control de variables). Eso permite comparar el teacher con el student, y cualquier política de cascada con la baseline, sin introducir sesgos por cambios de configuración. Cada uno de los cuatro pipelines se describe con detalle en las secciones siguientes.

\subsection{Pipeline de destilación}

Este pipeline produce el \textit{student} que usarán las Políticas B y C. Se diseña como la variante más sencilla posible compatible con el presupuesto del TFG, asumiendo explícitamente que no se trata de investigar KD sino de obtener un modelo suficientemente bueno para alimentar la cascada.

\subsubsection{Selección de teacher y student}

La elección de modelos viene condicionada por tres restricciones. Han de ser abiertos y permitir \textit{distillation}, vLLM debe servirlos sin modificaciones, y las combinaciones entre \textit{teacher} y \textit{student} deben tener sentido (misma familia, tokenizer idéntico y órdenes de magnitud de diferencia en parámetros). El candidato principal del proyecto es la familia Qwen2.5~\cite{qwen2}, con un teacher de entre 7B y 14B y un student más pequeño (1,5B o similar) dentro de la misma familia. Este emparejamiento garantiza tokenizer común, lo que simplifica la destilación por respuesta, y deja un diferencial de coste interesante entre los dos modelos.

En caso de que el plan de contingencia R2 se active (destilación inviable con la cuota disponible), se usa directamente un modelo preentrenado pequeño de la misma familia. No se considera una derrota del TFG sino una decisión documentada: la cascada y las políticas se pueden evaluar igual, aceptando un techo de calidad algo más bajo para el student.

\subsubsection{Datos para destilar}

Un punto clave del pipeline es \textbf{de qué datos sale el corpus de destilación}. En una destilación off-policy clásica, se toma un conjunto de prompts, se le pasan al teacher, y el student aprende a imitar las respuestas (o, mejor, las distribuciones) del teacher. Para este TFG el corpus tiene que cumplir dos condiciones.
\begin{enumerate}[nosep,leftmargin=*]
\item Ser representativo del \textit{tipo de peticiones} sobre el que se va a evaluar luego. Como la evaluación final se hace sobre razonamiento matemático, conviene que el corpus de destilación tenga también un peso importante de problemas de ese estilo, aunque no sean exactamente los mismos que aparecen en el benchmark.
\item No coincidir con los conjuntos de \textit{test} de los benchmarks que se usan en la evaluación. Cualquier contaminación invalida la comparación posterior: un student entrenado directamente con los test sets daría una falsa apariencia de calidad.
\end{enumerate}
En la práctica se usa una combinación de splits de entrenamiento públicos y datos sintéticos generados con el teacher sobre prompts semilla, con un filtrado básico de calidad (descartar respuestas vacías, respuestas con código incorrecto, respuestas muy cortas).

\subsubsection{Función de pérdida y criterios de diseño}

La pérdida es la destilación por respuesta descrita en la Sección~\ref{sec:knowledge-distillation}, y se escribe
\[
\mathcal{L}(x, y) = \alpha \, \mathrm{CE}\!\big(y,\, p^s(x;1)\big) + (1-\alpha)\, T^2 \, \mathrm{KL}\!\big(p^t(x;T)\,\|\,p^s(x;T)\big),
\]
con $T$ y $\alpha$ como hiperparámetros. Los valores exactos se eligen en los primeros experimentos del pipeline; la elección no es tanto un ``tuning'' exhaustivo como una decisión de diseño justificada, porque el presupuesto de GPU del TFG no da para hacer búsqueda masiva. Como punto de partida razonable se usan valores estándar de la literatura ($T\in[2, 4]$, $\alpha\in[0{,}1, 0{,}3]$) y se valida que la pérdida baja de forma estable.

\textbf{Criterio de parada.} El entrenamiento se detiene cuando la pérdida en validación deja de mejorar durante un número fijo de pasos, con un tope absoluto de pasos para acotar el uso de GPU. No se persigue \textit{squeezear} el último 1\,\% de accuracy del student: basta con que la destilación haya convergido a un comportamiento razonablemente parecido al teacher.

\textbf{Criterio de aceptación.} Para que el student destilado entre al pipeline de routing, tiene que cumplir (a) no colapsar ni generar outputs degenerados en una muestra de 100 prompts de prueba; (b) mantener accuracy por encima de un umbral mínimo absoluto (que se fija como una fracción de la accuracy del teacher) en el benchmark de calibración; (c) responder a la misma API que el teacher, para que la cascada se pueda montar sin lógica específica por modelo.

\subsubsection{Ramas experimentales de destilación}

Durante el proyecto se contemplan tres ramas en el pipeline de destilación, de menor a mayor ambición. No se ejecutan todas obligatoriamente: la decisión de explorar las siguientes depende del éxito y del presupuesto residual.

\begin{table}[H]
\centering
\caption{Ramas de destilación contempladas en el TFG, ordenadas por ambición creciente. Fuente: elaboración propia.}
\label{tab:ramas-kd}
\footnotesize
\begin{tabularx}{\textwidth}{>{\raggedright\arraybackslash}p{3.5cm} >{\raggedright\arraybackslash}X >{\raggedright\arraybackslash}p{3.2cm}}
\toprule
\textbf{Rama} & \textbf{Descripción} & \textbf{Decisión} \\
\midrule
K1. Off-policy básica & Corpus precomputado con el teacher; pérdida KD clásica con hard + soft labels. & \textbf{Elegida como base.} \\
\addlinespace
K2. Off-policy con énfasis en fallos del student & Se inyectan al corpus los ejemplos donde el student inicial falla, para reforzar esas zonas. & Solo si K1 deja un gap grande en un subconjunto identificable. \\
\addlinespace
K3. On-policy (GKD) & El student genera, el teacher puntúa, pérdida en secuencias del propio student~\cite{agarwal2024gkd}. & Descartada por presupuesto de GPU dentro del TFG. \\
\bottomrule
\end{tabularx}
\end{table}

La decisión de descartar la rama K3 para el TFG no quiere decir que sea peor; quiere decir que no es viable dentro del presupuesto actual. Se deja mencionada y referenciada para el capítulo de trabajo futuro.

\subsection{Pipeline de serving}

El pipeline de serving es el encargado de convertir dos checkpoints (teacher y student) en dos servicios de inferencia estables, accesibles por una API HTTP uniforme, que el resto de pipelines pueden usar como caja negra.

\subsubsection{Arquitectura de despliegue}

Cada modelo se sirve con una instancia independiente de vLLM, ejecutada como un job SLURM en una GPU H100 dedicada. Las dos instancias pueden vivir en el mismo nodo (si hay varias GPUs disponibles) o en nodos distintos. Ambas exponen una API compatible con el formato OpenAI (\texttt{/v1/completions} y \texttt{/v1/chat/completions}), de modo que el cliente de evaluación y el router las tratan como endpoints equivalentes salvo por el nombre del modelo.

\begin{figure}[H]
\centering
\begin{tikzpicture}[
>=Latex,
node distance=6mm and 10mm,
nodo/.style={draw, dashed, rounded corners=4pt, inner sep=6pt},
srv/.style={draw, rounded corners=2pt, align=center, minimum width=34mm, minimum height=11mm, font=\footnotesize, thick},
cli/.style={draw, rounded corners=2pt, align=center, minimum width=30mm, minimum height=9mm, font=\footnotesize, fill=yellow!22}
]
\node[cli] (cli) at (0,0) {Cliente de evaluación};

\node[srv, fill=bcteacher!30] (tea) at (7,1) {vLLM teacher\\(H100, API :8000)};
\node[srv, fill=bcstudent!30] (stu) at (7,-1) {vLLM student\\(H100, API :8001)};

\draw[->] (cli.east) -- (tea.west);
\draw[->] (cli.east) -- (stu.west);

\node[nodo, fit=(tea)(stu), label={[font=\scriptsize]below:Nodo GPU MareNostrum~5}] {};
\end{tikzpicture}
\caption{Arquitectura de despliegue: dos servidores vLLM independientes exponen APIs compatibles OpenAI y son atacados por el mismo cliente. Las métricas se toman con el cliente y el servidor corriendo en el mismo nodo para eliminar latencia de red externa. Fuente: elaboración propia.}
\label{fig:arq-serving}
\end{figure}

\subsubsection{Configuración del servidor}

La configuración de vLLM se fija una sola vez y se comparte entre teacher y student salvo por el modelo cargado. Esto es crucial para la comparabilidad, pues cualquier diferencia de parámetros de serving entre modelos contamina las métricas. Los parámetros principales incluyen los siguientes.
\begin{itemize}[nosep,leftmargin=*]
\item \texttt{dtype} = \texttt{bfloat16} (o FP16), uniforme entre ambos modelos.
\item \texttt{max\_model\_len} acotado a un valor razonable para GSM8K y MATH (suficiente para prompt+CoT+respuesta).
\item \texttt{gpu\_memory\_utilization} $\sim 0{,}9$ para dejar margen al sistema.
\item \texttt{swap\_space} = 0 (no se permite swap a CPU; un swap altera las métricas).
\item Parámetros de decoding fijos: \texttt{temperature}=0, \texttt{top\_p}=1, \texttt{max\_tokens} uniformes por benchmark.
\end{itemize}

La elección de \texttt{temperature}=0 es deliberada. No interesa estudiar sampling estocástico en este TFG, porque añade varianza entre ejecuciones que haría la comparativa mucho más difícil de interpretar. Para discutir la cascada basta con una política determinista.

\subsubsection{Instrumentación}

Cada servidor emite logs estructurados por petición que incluyen, entre otros campos, el timestamp de entrada, el del primer token, el del último token, el número de tokens de prompt y de generación, el modelo usado y un identificador de petición. El cliente, a su vez, marca sus propios timestamps y asocia cada petición a un experimento. A partir de estos dos logs se reconstruyen las métricas de servicio (TTFT, latencia p50/p95, throughput) sin necesidad de instrumentación adicional.

La memoria GPU se monitoriza con \texttt{nvidia-smi} en modo \texttt{dmon} durante la ventana del experimento, con un sampling de un segundo. Se registra el máximo de VRAM ocupada, que es el indicador útil. Si el máximo supera un umbral, el experimento no es seguro (riesgo de OOM en la siguiente ejecución) y se re-dimensiona.

\subsection{Pipeline de evaluación}

El pipeline de evaluación es el que aplica el protocolo de medida a un endpoint dado (teacher, student o una política de cascada) y devuelve un conjunto cerrado de métricas. Es la pieza que garantiza que todo se mide con la misma regla.

\subsubsection{Esquema del cliente}

El cliente de evaluación es un script Python que encadena tres pasos.

\begin{enumerate}[nosep,leftmargin=*]
\item \textbf{Carga del benchmark.} Lee el dataset (GSM8K, MATH) y aplica la plantilla de prompt fijada. La plantilla es la misma para todos los modelos; puede incluir \textit{few-shot examples} comunes.
\item \textbf{Generación.} Envía las peticiones al endpoint, con una política de concurrencia configurable. Para métricas de calidad la concurrencia se puede poner a 1 (no altera el resultado); para métricas de throughput se suelen hacer barridos con distintos niveles de concurrencia.
\item \textbf{Evaluación.} Extrae la respuesta final con un parser común (``el número después de \texttt{\#\#\#\#}'' en GSM8K, expresión en una caja \texttt{\textbackslash boxed\{\}} en MATH) y la compara con la verdad vía \textit{exact match}.
\end{enumerate}

El resultado por experimento es un fichero JSON con una entrada por petición (prompt, respuesta, métrica de calidad, tiempos) y un fichero CSV con el resumen agregado (accuracy, p50/p95 de TTFT y latencia, throughput medio, memoria GPU máxima).

\subsubsection{Benchmarks y su papel}

\textbf{GSM8K}~\cite{cobbe2021gsm8k} es un conjunto de unos 1\,300 problemas de aritmética escolar con razonamiento en varios pasos, pensado para evaluar capacidad matemática básica. La métrica es exact match sobre el número final. Es el benchmark de cabecera porque es estable, bien establecido y refleja una tarea donde el student tiene una probabilidad razonable de resolver una parte, pero falla en la cola más difícil.

\textbf{MATH}~\cite{hendrycks2021math} es más duro: problemas de olimpiada clasificados por dificultad (niveles 1--5). Muchos student pequeños fallan la mayoría de los niveles 4--5, lo que lo convierte en una piedra de toque útil para la cascada. Si el router tiene alguna capacidad, debería detectar que los niveles 4--5 merecen el teacher mucho más que los 1--2.

El TFG deja explícita la posibilidad de que MATH o GSM8K se sustituyan, combinen o complementen durante la ejecución del proyecto (por ejemplo, usar solo un subconjunto de MATH si el coste de evaluar el conjunto completo es prohibitivo). Cualquier decisión de este tipo queda documentada como parte de las ramas experimentales (Subsección~\ref{subsec:ramas-experimentales}).

\subsubsection{Métricas reportadas}

Por cada experimento se reportan las magnitudes de la Tabla~\ref{tab:metricas-reportadas}.

\begin{table}[H]
\centering
\caption{Métricas reportadas por el pipeline de evaluación para cada experimento. Fuente: elaboración propia.}
\label{tab:metricas-reportadas}
\footnotesize
\begin{tabularx}{\textwidth}{>{\raggedright\arraybackslash}p{3.5cm} >{\raggedright\arraybackslash}X >{\raggedright\arraybackslash}p{3.2cm}}
\toprule
\textbf{Métrica} & \textbf{Definición} & \textbf{Unidad} \\
\midrule
Accuracy / EM & Fracción de peticiones con respuesta correcta según el parser. & \% \\
TTFT p50 / p95 & Percentiles 50 y 95 del tiempo hasta el primer token. & ms \\
Latencia p50 / p95 & Percentiles 50 y 95 del tiempo total de servicio. & ms \\
Throughput medio & Tokens generados por segundo agregados del sistema. & tok/s \\
VRAM pico & Máximo de VRAM ocupada en la ventana de medida. & GB \\
\% escaladas al teacher & En políticas B y C: fracción de peticiones que terminan usando el teacher. & \% \\
Energía consumida & Registrada vía \texttt{sacct ConsumedEnergy} a nivel de job SLURM. & Wh \\
\bottomrule
\end{tabularx}
\end{table}

Con estas métricas, cada política (A, B o C) queda caracterizada por un punto en un espacio multidimensional calidad--coste--latencia. La discusión final de resultados consiste en comparar esos puntos y razonar bajo qué restricciones operativas una política domina a las otras.

\subsection{Pipeline de routing}

El pipeline de routing materializa la lógica de las Políticas B y C y las somete al mismo protocolo de medida que el teacher y el student aislados.

\subsubsection{Política A (baseline)}

No hay pipeline específico; se reutiliza el pipeline de evaluación atacando únicamente al endpoint del teacher. Sirve como referencia contra la que se miden B y C.

\subsubsection{Política B (cascada forzada)}

La Política B se implementa como un pequeño módulo cliente que, para cada petición, ejecuta la secuencia siguiente.
\begin{enumerate}[nosep,leftmargin=*]
\item Envía la petición al student y espera su respuesta.
\item Aplica una función \textit{quality gate} $g(\cdot)$ sobre la respuesta del student.
\item Si $g$ acepta, devuelve la respuesta del student.
\item Si $g$ rechaza, envía la misma petición al teacher y devuelve su respuesta.
\end{enumerate}

La clave del diseño es qué función se usa como \textit{quality gate}. Hay un abanico de opciones con distinto coste y distinto grado de información.
\begin{itemize}[nosep,leftmargin=*]
\item \textbf{Parser estructural.} Rechazar si la respuesta no tiene el formato esperado (por ejemplo, si no contiene un \texttt{\#\#\#\#} seguido de número en GSM8K, o una caja \texttt{\textbackslash boxed} en MATH). Es gratis pero solo detecta errores de formato, no de contenido.
\item \textbf{Señales de confianza del student.} Usar log-probabilidades del último token, entropía media de la distribución, margen top-1 vs top-2. Más informativo, sigue siendo barato.
\item \textbf{Segunda muestra de verificación.} Pedir al student que se \textit{autoverifique}, o comparar dos generaciones con temperatura baja y rechazar si difieren. Más caro.
\item \textbf{LLM juez externo.} Preguntar a un tercer modelo (o al propio teacher con un prompt ligero) si la respuesta parece correcta. Caro pero preciso.
\end{itemize}

El TFG adopta la combinación ``parser estructural + señales de confianza'' porque es la más defendible como \textit{cascada forzada} en el sentido estricto. El \textit{quality gate} es barato, no requiere ejecutar el teacher ni entrenar modelos adicionales. Las alternativas más caras se dejan como extensiones para trabajo futuro.

\subsubsection{Política C (cascada por confianza)}

La Política C introduce un \textbf{predictor} que decide \textit{a priori} si una petición se envía al student o directamente al teacher. El objetivo es evitar el coste extra de las peticiones que, en la Política B, acaban escalando tras ejecutarse en el student.

\textbf{Features del predictor.} Se consideran tres grupos.
\begin{enumerate}[nosep,leftmargin=*]
\item \textit{Del prompt:} longitud en tokens, presencia de símbolos matemáticos, ratio de palabras por línea, densidad de números, idioma detectado, nivel de dificultad si es un benchmark que lo anota.
\item \textit{Del student (con pasada corta):} log-prob del primer token generado, entropía media de los primeros $k$ tokens, perplejidad del prompt. Esto implica pagar una pequeña parte del student; no es gratis pero es barato.
\item \textit{Mixtas.} Combinaciones de las anteriores con un clasificador (regresión logística o gradient boosting).
\end{enumerate}

\textbf{Etiquetas de entrenamiento.} Para cada prompt del conjunto de entrenamiento del predictor se ejecutan \textit{offline} el student y el teacher, y se etiqueta como ``fácil'' si el student acierta (y por tanto se puede servir con el pequeño) y como ``difícil'' si no. Sobre este dataset se entrena el clasificador.

\textbf{Umbral de decisión.} El predictor devuelve una probabilidad de ``fácil''; el umbral que decide si se usa student o teacher se calibra sobre un conjunto de validación para balancear ahorro y calidad. Se reporta la curva completa (ROC-like) en los resultados, no un único punto.

\textbf{Latencia del router.} El router corre en CPU con features baratas o con una pasada corta del student, pero aun así añade una latencia no nula al pipeline. El objetivo del TFG es caracterizarla explícitamente, porque forma parte del coste de la Política C.

\subsubsection{Criterio de comparación entre políticas}

Para decidir si una política es mejor que otra no basta con mirar un número. Las tres conviven en un espacio multiobjetivo con al menos tres ejes: accuracy, fracción de peticiones al teacher (proxy de coste) y latencia p95. Una política domina a otra si es mejor o igual en los tres ejes y mejor en al menos uno. Lo más probable es que no haya dominación estricta: la Política C será mejor en coste y latencia a costa de un pequeño margen de accuracy, y la Política B un caso intermedio.

La discusión de resultados se organiza en torno a esa frontera de Pareto calidad--coste y no en torno a un ``ganador''. En el capítulo de resultados se reportarán las tres políticas lado a lado con el mismo conjunto de métricas, y se razonará en qué escenario operativo cada una es preferible.

\subsection{Ramas experimentales y decisiones de diseño}
\label{subsec:ramas-experimentales}

Un TFG experimental no se reduce a la rama que funciona. Mucha información útil está en las ramas que se probaron, se descartaron o se aparcaron, y en los motivos de cada decisión. Esta sección recoge de forma estructurada esas decisiones y actúa como registro permanente; se irá ampliando a medida que el trabajo avance, pero su esqueleto ya es conocido.

\begin{table}[H]
\centering
\caption{Ramas experimentales y decisiones de diseño del TFG. Elaboración propia.}
\label{tab:ramas}
\footnotesize
\begin{tabularx}{\textwidth}{>{\raggedright\arraybackslash}p{3cm} >{\raggedright\arraybackslash}X >{\raggedright\arraybackslash}p{3cm}}
\toprule
\textbf{Decisión} & \textbf{Alternativas y justificación} & \textbf{Estado} \\
\midrule
Familia de modelos & Qwen2.5 (misma familia teacher/student, tokenizer común); alternativas: Llama3, Mistral. Qwen2.5 elegida por madurez del ecosistema y cobertura de tamaños. & \textbf{Fijada.} \\
\addlinespace
Motor de serving & vLLM vs TGI vs llama.cpp. Elegido vLLM por PagedAttention, madurez, uso estándar en BSC y API compatible OpenAI. & \textbf{Fijada.} \\
\addlinespace
Tipo de KD & Off-policy response-based (K1) vs on-policy GKD (K3) vs feature-based. Elegida K1 por coste y simplicidad; K3 descartada por GPU. & \textbf{Fijada.} \\
\addlinespace
Benchmark principal & GSM8K + MATH vs HumanEval + HELM vs mezcla. Elegidos GSM8K+MATH por cobertura de razonamiento y estabilidad. & \textbf{Fijada, sujeta a revisión.} \\
\addlinespace
Quality gate Política~B & Parser + señales de confianza vs verificador externo vs self-consistency. Elegida la primera por coste. & \textbf{Fijada.} \\
\addlinespace
Predictor Política~C & Features mixtas con clasificador ligero vs LLM juez vs heurística manual. Elegido clasificador ligero por interpretabilidad y coste. & \textbf{Fijada, tuning pendiente.} \\
\addlinespace
Temperatura de decoding & 0 vs sampling con $T>0$. Fijada a 0 para eliminar varianza y simplificar comparación. & \textbf{Fijada.} \\
\addlinespace
Concurrencia de evaluación & 1 para calidad, sweep para throughput. Evita que el mismo experimento mezcle dos regímenes. & \textbf{Fijada.} \\
\bottomrule
\end{tabularx}
\end{table}

Cada una de estas decisiones podría cambiar a lo largo del proyecto; cuando cambie, se actualizará esta tabla y se documentará el motivo en \texttt{CHANGELOG.md}. El objetivo es que la memoria, en su versión final, cuente no solo qué se hizo sino también qué se probó, qué se descartó y por qué.

%%%%%%%%%%%%%%%%%%%%%%%%%%%%%%%%%%%%%%%%%%%%%%%%%%%%%%%%%%%%%%%%%
\part{Implementación}
\label{part:implementacion}
%%%%%%%%%%%%%%%%%%%%%%%%%%%%%%%%%%%%%%%%%%%%%%%%%%%%%%%%%%%%%%%%%

\section{Implementación del sistema}
\label{sec:implementacion-sistema}

\subsection{Estructura del repositorio y componentes}

\textit{Sección por completar a medida que el código estabilice su estructura definitiva. Placeholder estructural con lo que ya es firme.}

El código del TFG se organiza en torno a cuatro directorios principales: \texttt{distill/} (scripts de destilación), \texttt{serving/} (configuraciones y launchers de vLLM), \texttt{eval/} (cliente de evaluación y scripts de post-proceso) y \texttt{routing/} (implementación de las Políticas B y C y el predictor). Los experimentos reproducibles viven en \texttt{experiments/}, cada uno con un fichero de configuración, un script SLURM, y un directorio de resultados trazable al JobID que lo generó.

Los detalles técnicos de cada componente (librerías concretas, versiones, estructura de las configuraciones, formato de logs) se desarrollarán en este capítulo cuando el código del TFG esté estabilizado, para no tener que reescribirlo dos veces si aún hay decisiones abiertas.

\subsection{Integración con SLURM y entorno de ejecución}

\textit{Por completar con scripts SLURM finales.}

Esta sección describirá cómo se despliega el entorno en MareNostrum~5, incluyendo la estrategia de módulos, el entorno Python, los scripts SLURM tipo para destilación, serving y evaluación, y los mecanismos de registro automático de consumo energético vía \texttt{sacct}.

\subsection{Problemas técnicos relevantes y soluciones}

\textit{Se irá completando con los hallazgos reales a medida que aparezcan.}

Este bloque recoge los problemas no triviales encontrados durante la implementación (configuraciones de vLLM, versiones incompatibles, límites de la cola, timeouts, etc.) y cómo se resolvieron, con el objetivo de servir de guía para quien continúe esta línea.

%%%%%%%%%%%%%%%%%%%%%%%%%%%%%%%%%%%%%%%%%%%%%%%%%%%%%%%%%%%%%%%%%
\part{Experimentación y resultados}
\label{part:experimentacion}
%%%%%%%%%%%%%%%%%%%%%%%%%%%%%%%%%%%%%%%%%%%%%%%%%%%%%%%%%%%%%%%%%

\section{Experimentación y análisis de resultados}
\label{sec:experimentacion-resultados}

\subsection{Protocolo experimental}

\textit{Versión resumida del protocolo ya definido en la Parte~\ref{part:gestion-tesis}, con referencias cruzadas a las decisiones de la Parte~\ref{part:diseno-sistema}. Se consolidará aquí cuando se cierre el diseño experimental definitivo.}

\subsection{Resultados}

\textit{Capítulo a redactar cuando los experimentos estén cerrados. Se incluirán:} 
\begin{itemize}[nosep,leftmargin=*]
\item Caracterización del teacher (calidad y servicio).
\item Caracterización del student (calidad y servicio).
\item Comparativa de las tres políticas A, B, C en calidad, coste y latencia.
\item Análisis de la frontera de Pareto calidad-coste.
\item Consumo energético medido frente a las estimaciones iniciales de la Parte~\ref{part:gestion-tesis} (sostenibilidad).
\end{itemize}

\subsection{Discusión}

\textit{A redactar con los resultados en la mano: qué se confirma de las hipótesis del diseño, qué sorprende, qué limitaciones salen a la luz.}

%%%%%%%%%%%%%%%%%%%%%%%%%%%%%%%%%%%%%%%%%%%%%%%%%%%%%%%%%%%%%%%%%
\part{Conclusiones}
\label{part:conclusiones}
%%%%%%%%%%%%%%%%%%%%%%%%%%%%%%%%%%%%%%%%%%%%%%%%%%%%%%%%%%%%%%%%%

\section{Conclusiones y trabajo futuro}
\label{sec:conclusiones-futuro}

\subsection{Cumplimiento de objetivos}

\textit{Revisión OE1--OE6 frente a resultados reales.}

\subsection{Conclusiones finales}

\textit{Balance global del TFG al cierre.}

\subsection{Trabajo futuro}

\textit{Líneas naturales de continuación: GKD on-policy, routing con LLM juez, cascadas con más de dos niveles, integración con speculative decoding.}

%%%%%%%%%%%%%%%%%%%%%%%%%%%%%%%%%%%%%%%%%%%%%%%%%%%%%%%%%%%%%%%%%
% BIBLIOGRAFÍA
%%%%%%%%%%%%%%%%%%%%%%%%%%%%%%%%%%%%%%%%%%%%%%%%%%%%%%%%%%%%%%%%%
\begin{thebibliography}{99}

% --- Referencias bibliográficas (incl. planificación y normativa académica) ---
\bibitem{hinton2015distilling}
G.~Hinton, O.~Vinyals, and J.~Dean, ``Distilling the Knowledge in a Neural Network,'' arXiv preprint arXiv:1503.02531, 2015.

\bibitem{vllm2023}
W.~Kwon \textit{et al.}, ``Efficient Memory Management for Large Language Model Serving with PagedAttention,'' in \textit{Proceedings of the 29th Symposium on Operating Systems Principles (SOSP)}, ACM, 2023.

\bibitem{chen2023frugalgpt}
L.~Chen, M.~Zaharia, and J.~Zou, ``FrugalGPT: How to Use Large Language Models while Reducing Cost and Improving Performance,'' arXiv preprint arXiv:2305.05176, 2023.

\bibitem{dekoninck2024routing}
J.~Dekoninck \textit{et al.}, ``RouteLLM: Learning to Route LLMs with Preference Data,'' arXiv preprint arXiv:2406.18665, 2024.

\bibitem{touvron2023llama}
H.~Touvron \textit{et al.}, ``LLaMA: Open and Efficient Foundation Language Models,'' arXiv preprint arXiv:2302.13971, 2023.

\bibitem{jiang2023mistral}
A.~Q.~Jiang \textit{et al.}, ``Mistral 7B,'' arXiv preprint arXiv:2310.06825, 2023.

\bibitem{cobbe2021gsm8k}
K.~Cobbe \textit{et al.}, ``Training Verifiers to Solve Math Word Problems,'' arXiv preprint arXiv:2110.14168, 2021.

\bibitem{hendrycks2021math}
D.~Hendrycks \textit{et al.}, ``Measuring Mathematical Problem Solving with the MATH Dataset,'' in \textit{Proceedings of NeurIPS}, 2021.

\bibitem{vllm-docs}
vLLM Project, ``vLLM Documentation,'' 2024. [En línea]. Disponible en: \url{https://docs.vllm.ai/}. Consultado: 22 de enero de 2026.

\bibitem{gep23}
Módulo 2.3 -- Gestión del tiempo. Asignatura GEP, UPC, 2026.

\bibitem{gpi}
J.\,R.~Rodríguez, J.~García, and I.~Lamarca, \textit{Gestión de Proyectos Informáticos: métodos, herramientas y casos}. Editorial UOC.

\bibitem{mcconnell06}
S.~McConnell, \textit{Software Estimation: Demystifying the Black Art}. Microsoft Press, 2006.

\bibitem{pmbok}
Project Management Institute, \textit{A Guide to the Project Management Body of Knowledge (PMBOK Guide)}. 7.ª~ed., PMI, 2021.

\bibitem{schwaber2020scrum}
K.~Schwaber and J.~Sutherland, \textit{The Scrum Guide: The Definitive Guide to Scrum}. Scrum.org, 2020.

\bibitem{iso21502}
International Organization for Standardization, \textit{ISO 21502:2020 Project, Programme and Portfolio Management -- Guidance on Project Management}. ISO, 2020.

\bibitem{mn5}
Barcelona Supercomputing Center, \textit{MareNostrum~5 -- System Overview}. BSC, 2024. [En línea]. Disponible en: \url{https://www.bsc.es/marenostrum/marenostrum-5}. Consultado: febrero de 2026.

\bibitem{slurm}
A.\,B.~Yoo, M.\,A.~Jette, and M.~Grondona, ``SLURM: Simple Linux Utility for Resource Management,'' in \textit{Job Scheduling Strategies for Parallel Processing (JSSPP)}, Springer, 2003.

\bibitem{gep23e2}
M.~Romera Peral, \textit{Segunda entrega GEP: Planificación temporal}. UPC, 2026.

\bibitem{hays2026}
Hays, \textit{Guía del Mercado Laboral 2026 -- Tecnologías de la Información} (tablas salariales). Hays España, 2026.

\bibitem{ree2024}
Red Eléctrica de España, \textit{Informe del sistema eléctrico español 2024}. REE, 2024. [En línea]. Disponible en: \url{https://www.ree.es/}. Consultado: febrero de 2026.

\bibitem{lambda2026}
Lambda Labs, \textit{GPU Cloud Pricing -- 1xH100 SXM5 80\,GB}. [En línea]. Disponible en: \url{https://lambdalabs.com/service/gpu-cloud}. Consultado: febrero de 2026.

\bibitem{aws2026}
Amazon Web Services, \textit{Amazon EC2 P5 Instances Pricing (p5.xlarge, us-east-1)}. [En línea]. Disponible en: \url{https://aws.amazon.com/ec2/pricing/on-demand/}. Consultado: febrero de 2026.

% --- Añadidas marco teórico / diseño ---
\bibitem{vaswani2017attention}
A.~Vaswani \textit{et al.}, ``Attention Is All You Need,'' in \textit{Advances in Neural Information Processing Systems (NeurIPS)}, 2017.

\bibitem{sanh2019distilbert}
V.~Sanh, L.~Debut, J.~Chaumond, and T.~Wolf, ``DistilBERT, a distilled version of BERT: smaller, faster, cheaper and lighter,'' arXiv preprint arXiv:1910.01108, 2019.

\bibitem{gou2021survey}
J.~Gou, B.~Yu, S.\,J.~Maybank, and D.~Tao, ``Knowledge Distillation: A Survey,'' \textit{International Journal of Computer Vision}, vol.~129, no.~6, pp.~1789--1819, 2021.

\bibitem{agarwal2024gkd}
R.~Agarwal \textit{et al.}, ``On-Policy Distillation of Language Models: Learning from Self-Generated Mistakes,'' arXiv preprint arXiv:2306.13649, 2024.

\bibitem{dao2022flashattention}
T.~Dao, D.\,Y.~Fu, S.~Ermon, A.~Rudra, and C.~Ré, ``FlashAttention: Fast and Memory-Efficient Exact Attention with IO-Awareness,'' in \textit{Advances in Neural Information Processing Systems (NeurIPS)}, 2022.

\bibitem{dao2023flash2}
T.~Dao, ``FlashAttention-2: Faster Attention with Better Parallelism and Work Partitioning,'' arXiv preprint arXiv:2307.08691, 2023.

\bibitem{leviathan2023speculative}
Y.~Leviathan, M.~Kalman, and Y.~Matias, ``Fast Inference from Transformers via Speculative Decoding,'' in \textit{Proceedings of the 40th International Conference on Machine Learning (ICML)}, 2023.

\bibitem{viola2001cascade}
P.~Viola and M.~Jones, ``Rapid Object Detection using a Boosted Cascade of Simple Features,'' in \textit{Proceedings of CVPR}, 2001.

\bibitem{qwen2}
A.~Yang \textit{et al.}, ``Qwen2.5 Technical Report,'' arXiv preprint arXiv:2412.15115, 2024.

\bibitem{su2021roformer}
J.~Su, Y.~Lu, S.~Pan, B.~Wen, and Y.~Liu, ``RoFormer: Enhanced Transformer with Rotary Position Embedding,'' arXiv preprint arXiv:2104.09864, 2021.

\bibitem{hu2021lora}
E.\,J.~Hu \textit{et al.}, ``LoRA: Low-Rank Adaptation of Large Language Models,'' in \textit{International Conference on Learning Representations (ICLR)}, 2022.

\bibitem{dettmers2023qlora}
T.~Dettmers, A.~Pagnoni, A.~Holtzman, and L.~Zettlemoyer, ``QLoRA: Efficient Finetuning of Quantized LLMs,'' in \textit{Advances in Neural Information Processing Systems (NeurIPS)}, 2023.

\bibitem{kim2016sequence}
Y.~Kim and A.\,M.~Rush, ``Sequence-Level Knowledge Distillation,'' in \textit{Proceedings of EMNLP}, 2016.

\bibitem{taori2023alpaca}
R.~Taori \textit{et al.}, ``Stanford Alpaca: An Instruction-following LLaMA model,'' Stanford CRFM, 2023.

\bibitem{ong2025routellm}
I.~Ong, A.~Almahairi, V.~Wu, W.-L.~Chiang, T.~Wu, J.\,E.~Gonzalez, M.\,W.~Kadous, and I.~Stoica, ``RouteLLM: Learning to Route LLMs with Preference Data,'' arXiv preprint arXiv:2406.18665, 2024.

\end{thebibliography}

\end{document}
